% =====================================================================
%  Gemeinsame Vorlage der Arbeitsblätter – Grundwissen Q13
%  – Informatik 12 (gleiches Layout wie informatik/10/java/arbeitsblaetter/
%  und informatik/11/graphen/arbeitsblaetter/ab-vorlage.tex)
%  Einbinden in jedem Blatt direkt nach \documentclass: % =====================================================================
%  Gemeinsame Vorlage der Arbeitsblätter – Grundwissen Q13
%  – Informatik 12 (gleiches Layout wie informatik/10/java/arbeitsblaetter/
%  und informatik/11/graphen/arbeitsblaetter/ab-vorlage.tex)
%  Einbinden in jedem Blatt direkt nach \documentclass: % =====================================================================
%  Gemeinsame Vorlage der Arbeitsblätter – Grundwissen Q13
%  – Informatik 12 (gleiches Layout wie informatik/10/java/arbeitsblaetter/
%  und informatik/11/graphen/arbeitsblaetter/ab-vorlage.tex)
%  Einbinden in jedem Blatt direkt nach \documentclass: % =====================================================================
%  Gemeinsame Vorlage der Arbeitsblätter – Grundwissen Q13
%  – Informatik 12 (gleiches Layout wie informatik/10/java/arbeitsblaetter/
%  und informatik/11/graphen/arbeitsblaetter/ab-vorlage.tex)
%  Einbinden in jedem Blatt direkt nach \documentclass: \input{ab-vorlage}
%  Übersetzen mit XeLaTeX, LuaLaTeX, pdfLaTeX oder Tectonic.
%
%  Blatt:   xelatex ab-0-grundwissen-13.tex
%  Lösung:  xelatex -jobname=ab-0-grundwissen-13-lsg "\def\mitloesung{}\input{ab-0-grundwissen-13}"
%           (füllt alle Lücken rot aus; siehe \lsg, \lsglinien, \lsgoder)
% =====================================================================
\usepackage[a4paper,left=18mm,right=18mm,top=26mm,bottom=17mm,headheight=15mm,headsep=4mm,footskip=8mm]{geometry}
\usepackage{iftex}
\ifPDFTeX
  \usepackage[T1]{fontenc}
  \usepackage[utf8]{inputenc}
  \usepackage{helvet}
  \usepackage[scaled=0.86]{beramono}
  \renewcommand{\familydefault}{\sfdefault}
\else
  \usepackage{fontspec}
  \setmainfont{texgyreheros}[Extension=.otf,UprightFont=*-regular,BoldFont=*-bold,ItalicFont=*-italic,BoldItalicFont=*-bolditalic]
  \setmonofont{DejaVuSansMono}[Extension=.ttf,UprightFont=*,BoldFont=*-Bold,ItalicFont=*-Oblique,BoldItalicFont=*-BoldOblique,Scale=0.86]
\fi
\usepackage[ngerman,shorthands=off]{babel}
\usepackage{xcolor,graphicx,tabularx,array,enumitem,fancyhdr,tikz,listings,multicol,calc,etoolbox,pifont}
\usepackage[most]{tcolorbox}
\usetikzlibrary{arrows.meta,positioning,calc,decorations.pathreplacing,shapes.multipart}
\usepackage[hidelinks]{hyperref}
\setlength{\parindent}{0pt}
\setlength{\parskip}{3pt}
\setlist{nosep,leftmargin=*}

% ---------- Lösungsschalter ----------
\newif\ifloesung
\ifdefined\mitloesung\loesungtrue\fi
\newcommand{\lsgfarbe}{\color{red!75!black}}

% ---------- Kopf- und Fußzeile (einheitliche Form) ----------
\newcommand{\lehrkraft}{Hau}
\newcommand{\themenbereich}{Grundwissen 6 bis 12}
\newcommand{\abnummer}{}
\pagestyle{fancy}\fancyhf{}
\renewcommand{\headrulewidth}{0.4pt}
\fancyhead[L]{\small Klasse 13\\Datum:}
\fancyhead[C]{\small Themenbereich:\\\themenbereich}
\fancyhead[R]{\small \lehrkraft\ifloesung\\{\lsgfarbe\bfseries Lösung}\fi}
\fancyfoot[C]{\scriptsize edu-mrh.de\,·\,Informatik 13\,·\,Blatt \abnummer\ifloesung\ (Lösung)\fi\,·\,Seite \thepage}
\newcommand{\abtitel}[2]{\renewcommand{\abnummer}{#1}\begin{center}\Large\bfseries #1\quad\underline{#2}\end{center}\vspace{-1mm}}

% ---------- Boxen ----------
\newenvironment{aufgabe}[2][]{\begin{tcolorbox}[enhanced,breakable,colback=white,colframe=black,boxrule=0.7pt,arc=3mm,left=2mm,right=2mm,top=1.5mm,bottom=1.5mm,before skip=2.5mm,after skip=2.5mm]\textbf{Aufgabe #2)}\ifstrempty{#1}{}{ \textit{#1}}\par\vspace{0.6mm}}{\end{tcolorbox}}
\newtcolorbox{merke}[1][Merke]{enhanced,breakable,colback=green!5,colframe=green!40!black,boxrule=0.8pt,arc=2mm,left=2mm,right=2mm,top=1.5mm,bottom=1.5mm,before skip=2.5mm,after skip=2.5mm,title={#1},fonttitle=\bfseries,coltitle=white}
\newtcolorbox{hinweis}{enhanced,breakable,colback=yellow!12,colframe=orange!70!black,boxrule=0.6pt,arc=2mm,left=2mm,right=2mm,top=1mm,bottom=1mm,before skip=2mm,after skip=2mm}
\newcommand{\web}[1]{{\small\color{gray}\textit{Website: #1}}}

% ---------- Lücken, Linien, Karos ----------
\newcommand{\luecke}[1][3.5cm]{\rule[-0.5ex]{#1}{0.4pt}}
\newcommand{\linien}[1]{\par\vspace{1.5mm}\foreach \i in {1,...,#1}{\noindent\rule[0pt]{\linewidth}{0.3pt}\par\vspace{5.5mm}}\vspace{-3mm}}
\newcommand{\kariert}[1]{\par\vspace{1.5mm}\noindent\begin{tikzpicture}\draw[step=5mm,gray!55,thin] (0,0) grid (\linewidth-1pt,#1*5mm);\end{tikzpicture}\par}
\newcommand{\karobox}[2]{\begin{tikzpicture}\draw[step=5mm,gray!55,thin] (0,0) grid (#1,#2);\end{tikzpicture}}
\newcommand{\klassenkarte}[3]{\begin{tabular}{|l|}\hline\multicolumn{1}{|c|}{\textbf{#1}}\\\hline #2\\\hline #3\\\hline\end{tabular}}
\newcommand{\code}[1]{\texttt{#1}}

% Lücke mit Lösung: \lsg[Mindestbreite]{Lösungstext}
\newlength{\lsgbreite}
\newcommand{\lsg}[2][3.5cm]{\ifloesung\settowidth{\lsgbreite}{\,#2\,}\ifdim\lsgbreite<#1\setlength{\lsgbreite}{#1}\fi\underline{\makebox[\lsgbreite][l]{\lsgfarbe\,#2}}\else\luecke[#1]\fi}
% Schreiblinien mit Lösung: \lsglinien{Anzahl}{Lösungstext} (gleiche Höhe wie \linien)
\newcommand{\lsglinien}[2]{\ifloesung\par\vspace{1.5mm}\noindent\begin{minipage}[t][#1\dimexpr 6.9mm\relax][t]{\linewidth}\lsgfarbe #2\end{minipage}\par\vspace{-1.5mm}\else\linien{#1}\fi}
% Beliebiger Inhalt: \lsgoder{nur in der Lösung}{nur im Blatt}
\newcommand{\lsgoder}[2]{\ifloesung\leavevmode#1\else#2\fi}

% ---------- Verkettete Strukturen zeichnen ----------
% Objekt mit Kopf und Attributen:  \node[obj] (k1) {k1: Knoten \nodepart{second} daten = Hans\\ nachfolger};
\tikzset{
  obj/.style={draw,rectangle split,rectangle split parts=2,rectangle split part fill={orange!25,orange!8},
              rounded corners=2pt,font=\scriptsize,inner sep=2pt,align=left},
  qobj/.style={obj,rectangle split part fill={teal!25,teal!8}},
  pobj/.style={draw,rounded corners=2pt,fill=gray!8,font=\scriptsize,inner sep=2pt,align=left},
  zeiger/.style={thick,-{Stealth[length=2.2mm]}},
  lsgzeiger/.style={zeiger,red!75!black},
  nullref/.style={font=\scriptsize\ttfamily,inner sep=1pt},
  schritt/.style={circle,draw,fill=white,inner sep=0.6pt,font=\scriptsize\bfseries},
  lsgschritt/.style={schritt,draw=red!75!black,text=red!75!black},
}

% ---------- Java-Listings ----------
\lstset{language=Java,basicstyle=\ttfamily\small,keywordstyle=\bfseries,commentstyle=\color{gray!80!black},stringstyle=\color{black},showstringspaces=false,columns=flexible,keepspaces=true,tabsize=3,frame=single,framerule=0.4pt,rulecolor=\color{black},xleftmargin=2mm,framexleftmargin=1mm,aboveskip=2mm,belowskip=2mm,breaklines=true,escapechar=§,
 literate={ä}{{\"a}}1 {ö}{{\"o}}1 {ü}{{\"u}}1 {Ä}{{\"A}}1 {Ö}{{\"O}}1 {Ü}{{\"U}}1 {ß}{{\ss}}1 {–}{{--}}1 {„}{{\glqq}}1 {“}{{\grqq}}1 {→}{{$\to$}}1 {①}{{\ding{172}}}1 {②}{{\ding{173}}}1 {③}{{\ding{174}}}1 {④}{{\ding{175}}}1}
\lstnewenvironment{java}[1][]{\lstset{#1}}{}

%  Übersetzen mit XeLaTeX, LuaLaTeX, pdfLaTeX oder Tectonic.
%
%  Blatt:   xelatex ab-0-grundwissen-13.tex
%  Lösung:  xelatex -jobname=ab-0-grundwissen-13-lsg "\def\mitloesung{}\documentclass[11pt,a4paper]{article}
\input{ab-vorlage}
% Blatt:   tectonic ab-0-grundwissen-13.tex
% Lösung:  tectonic ab-0-grundwissen-13-lsg.tex  (Datei enthält nur: \def\mitloesung{}\input{ab-0-grundwissen-13})
% Seite 1: Grundwissen 6 bis 11 als ausgefüllte Übersicht (identisch in Blatt und Lösung)
% Seite 2: Übersichtsaufgaben zu den Inhalten der Jahrgangsstufe 12
\usetikzlibrary{shapes.geometric}
\newgeometry{left=14mm,right=14mm,top=21mm,bottom=12mm,headheight=11mm,headsep=3mm,footskip=5mm}
\ifloesung\tikzset{lsgbild/.style={red!75!black}}\else\tikzset{lsgbild/.style={opacity=0}}\fi
\newcommand{\gw}[1]{\par\vspace{0.6mm}{\bfseries\color{green!40!black}#1}\par}
\newcommand{\teil}[1]{{\bfseries #1}\par\vspace{0.5mm}}
\setlength{\columnsep}{5mm}
\setlength{\columnseprule}{0.3pt}
\newcommand{\kopf}[1]{\multicolumn{1}{c}{\textbf{#1}}}
\renewenvironment{aufgabe}[2][]{\begin{tcolorbox}[enhanced,breakable,colback=white,colframe=black,boxrule=0.7pt,arc=3mm,left=2mm,right=2mm,top=1mm,bottom=1mm,before skip=1.8mm,after skip=1.8mm]\textbf{Aufgabe #2)}\ifstrempty{#1}{}{ \textit{#1}}\par\vspace{0.3mm}}{\end{tcolorbox}}
\lstdefinestyle{klein}{basicstyle=\ttfamily\scriptsize,aboveskip=0.8mm,belowskip=0.8mm}
\begin{document}
\abtitel{0}{Grundwissen}
\vspace{-3mm}
\teil{Teil A: Grundwissen der Jahrgangsstufen 6 bis 11 (zum Nachschlagen)}

\begin{minipage}[t]{0.575\linewidth}
\vspace{0pt}
\begin{tikzpicture}[x=1cm,y=1cm,font=\footnotesize]
\node[anchor=north west] (kk) at (0,0) {\klassenkarte{Dreieck}{seitenlaenge\\farbe}{setzeFarbe(neu)}};
\node[anchor=south west,font=\footnotesize\bfseries] at (0,0.05) {Klasse (Bauplan)};
\node[anchor=north west,draw,rounded corners=2.5mm,inner sep=0pt] (ok) at (5.35,0) {\begin{tabular}{l}\rule{0pt}{3.5mm}\textbf{form1: Dreieck}\\[0.5mm]\hline\rule{0pt}{3.5mm}seitenlaenge = 20\\farbe = "blau"\\[0.5mm]\end{tabular}};
\node[anchor=south west,font=\footnotesize\bfseries] at (5.35,0.05) {Objekt (Instanz)};
\draw[decorate,decoration={brace,amplitude=3pt}] ([xshift=2pt]kk.north east) -- ([xshift=2pt]kk.north east|-{(0,-0.5)}) node[midway,right=4pt,font=\scriptsize] {Klassenname};
\draw[decorate,decoration={brace,amplitude=3pt}] ([xshift=2pt]kk.north east|-{(0,-0.58)}) -- ([xshift=2pt]kk.north east|-{(0,-1.35)}) node[midway,right=4pt,font=\scriptsize] {Attribute};
\draw[decorate,decoration={brace,amplitude=3pt}] ([xshift=2pt]kk.north east|-{(0,-1.43)}) -- ([xshift=2pt]kk.south east) node[midway,right=4pt,font=\scriptsize] {Methoden};
\draw[decorate,decoration={brace,amplitude=3pt}] ([xshift=2pt]ok.north east) -- ([xshift=2pt]ok.north east|-{(0,-0.55)}) node[midway,right=4pt,align=left,font=\scriptsize] {Objektname:\\Klasse};
\draw[decorate,decoration={brace,amplitude=3pt}] ([xshift=2pt]ok.north east|-{(0,-0.63)}) -- ([xshift=2pt]ok.south east) node[midway,right=4pt,align=left,font=\scriptsize] {Attribut-\\werte};
\end{tikzpicture}
\end{minipage}\hfill
\begin{minipage}[t]{0.405\linewidth}
\vspace{0pt}\footnotesize
\textbf{Objektorientierung}: Problemlösen durch Analyse der beteiligten \textbf{Objekte} und ihres Zusammenspiels. Die \textbf{Klasse} ist der Bauplan, ein \textbf{Objekt} eine konkrete Ausprägung (Instanz) der Klasse.\\[0.5mm]
\textbf{Punktnotation}: \code{objekt.methode(parameter);}\\ \code{objekt.attribut = wert;}\quad z.\,B. \code{karol.schritt();}
\end{minipage}

\vspace{-1.5mm}
\begin{multicols}{2}\footnotesize\setlength{\parskip}{0.8pt}
\gw{0.1 Beziehungen, Datentypen, Zugriffsrechte}
Objekte können \textbf{Referenzen} auf andere Objekte enthalten (\textbf{Referenzattribute}); die referenzierte Klasse wird zum \textbf{Datentyp}, z.\,B. \code{private Person eigentuemer;} in \code{Haus}:\\[0.3mm]
\begin{tikzpicture}[font=\scriptsize]
\node[draw,inner sep=2pt] (h) {\textbf{Haus}};
\node[draw,inner sep=2pt,right=17mm of h] (p) {\textbf{Person}};
\draw[-{Stealth}] (h) -- node[above,inner sep=1pt] {eigentuemer} node[below,pos=0.85,inner sep=1pt] {1} (p);
\end{tikzpicture}\quad {\scriptsize(Klassendiagramm ohne \code{String}, \code{int})}\\[0.3mm]
Weitere Datentypen: \code{boolean}, \code{char}, \code{double}, \code{float}, \code{long}, \code{byte}, Felder, jede Klasse.
Objekte kommunizieren über \textbf{Methodenaufrufe}. \textbf{Datenkapselung}: Attribute sind \code{private}, Zugriff nur über \textbf{Getter} und \textbf{Setter}.\\[0.5mm]
\renewcommand{\arraystretch}{1}
\begin{tabular}{@{}|c|l|l|@{}}\hline
\code{+} & \code{public} & öffentlich: jede Klasse darf zugreifen\\\hline
\code{\#} & \code{protected} & eigene Klasse und Unterklassen\\\hline
\code{-} & \code{private} & nur innerhalb der eigenen Klasse\\\hline
\end{tabular}

\gw{0.2 Bausteine von Algorithmen: Struktogramm und Java}
\begin{tikzpicture}[x=1cm,y=1cm,font=\scriptsize]
\draw (0,0) rectangle (3.95,1.1); \draw (0,1.1) -- (1.97,0.55) -- (3.95,1.1); \draw (0,0.55) -- (3.95,0.55); \draw (1.97,0) -- (1.97,0.55);
\node at (1.97,0.93) {\ttfamily\itshape istWand()}; \node at (0.33,0.66) {wahr}; \node at (3.55,0.66) {falsch};
\node at (0.98,0.27) {\ttfamily linksDrehen()}; \node at (2.96,0.27) {\ttfamily schritt()};
\begin{scope}[xshift=4.3cm]
\draw (0,0) rectangle (4.05,1.1); \draw (0.4,0) -- (0.4,0.75) -- (4.05,0.75);
\node[anchor=west] at (0.05,0.93) {Wiederhole solange {\ttfamily\itshape !istWand()}};
\node[anchor=west] at (0.5,0.53) {\ttfamily hinlegen()}; \node[anchor=west] at (0.5,0.2) {\ttfamily schritt()};
\end{scope}
\end{tikzpicture}\\[-1.2mm]
\begin{minipage}[t]{0.47\linewidth}
\begin{java}[style=klein]
if (istWand()) {
   linksDrehen();
} else {
   schritt();
}
\end{java}
\end{minipage}\hfill
\begin{minipage}[t]{0.49\linewidth}
\begin{java}[style=klein]
while (!istWand()) {
   hinlegen();
   schritt();
}
\end{java}
\end{minipage}\\[0.3mm]
\textbf{Sequenz}: Anweisungen untereinander; \textbf{einseitig}: \code{if} ohne \code{else}; „Wiederhole 5-mal“: \code{for (int i = 0; i < 5; i++) \{…\}}

\gw{0.3 Feld (Array)}
Folge \underline{fester Länge} von Werten bzw. Objekten \underline{desselben Typs}. Zugriff über den \textbf{Index}; die Zählung beginnt bei 0.
\begin{java}[style=klein]
int zahl[];            // Deklaration
zahl = new int[1000];  // Instanziierung (Index 0..999)
zahl[2] = 10;          // Wertzuweisung (3. Element)
\end{java}

\gw{0.4 Das EVA-Prinzip}
\begin{tikzpicture}[font=\scriptsize,node distance=2.5mm,every node/.style={align=center}]
\node[draw,minimum height=7mm] (e) {\textbf{Eingabe}\\\code{gewicht, groesse}};
\node[draw,ellipse,inner sep=0pt,right=of e] (v) {\textbf{Verarbeitung}\\\code{bmi = gewicht/}\\\code{(groesse*groesse)}};
\node[draw,minimum height=7mm,right=of v] (a) {\textbf{Ausgabe}\\\code{println(…)}};
\draw[-{Stealth}] (e) -- (v); \draw[-{Stealth}] (v) -- (a);
\end{tikzpicture}\\
Bei Methoden: Eingabe über \textbf{Parameter}, Ausgabe als \textbf{Rückgabewert} oder auf dem Bildschirm.

\gw{0.5 Vererbung, @Override und Polymorphismus}
Eine \textbf{Unterklasse} erbt Attribute und Methoden ihrer \textbf{Oberklasse}; in Java von genau \underline{einer} Klasse (keine Mehrfachvererbung).\\[0.3mm]
\begin{minipage}[t]{0.21\linewidth}
\vspace{0pt}
\begin{tikzpicture}[font=\scriptsize]
\node[draw,rectangle split,rectangle split parts=2,align=left,inner sep=2pt] (p) {\textbf{Person}\nodepart{second}alter\\name};
\node[draw,rectangle split,rectangle split parts=2,align=left,inner sep=2pt,below=4mm of p] (s) {\textbf{Schueler}\nodepart{second}klasse};
\draw[-{Triangle[open,length=2mm]}] (s) -- (p);
\end{tikzpicture}
\end{minipage}\hfill
\begin{minipage}[t]{0.78\linewidth}
\vspace{0pt}
\begin{java}[style=klein,aboveskip=0mm,belowskip=0mm]
public class Schueler extends Person {
   private String klasse;
   Schueler(int a, String n, String k) {
      super(a, n);  klasse = k;
   }
}
\end{java}
\end{minipage}\\[0.8mm]
\begin{tabular}{@{}|l|p{6.25cm}|@{}}\hline
\code{extends} & gibt die Oberklasse an\\\hline
\code{super(…)} & ruft den Konstruktor der Oberklasse auf\\\hline
\code{@Override} & Unterklasse \textbf{überschreibt} eine geerbte Methode (gleiche Signatur, neues Verhalten)\\\hline
Polymorphie & Variable vom Typ der Oberklasse kann auf Unterklassen-Objekte verweisen; ausgeführt wird die Methode der \emph{tatsächlichen} Klasse.\\\hline
\end{tabular}

\gw{0.6 Datenbanken und SQL}
Datenbanken speichern strukturierte Daten in \textbf{Tabellen}: Ein \textbf{Datensatz} (Zeile, Tupel) enthält die Attributwerte (Spalten $\to$ \textbf{Attribute}) eines Objekts. Der \textbf{Primärschlüssel} (unterstrichen; ein oder mehrere Attribute) identifiziert jeden Datensatz eindeutig. Beziehungen zwischen Tabellen: \textbf{Fremdschlüssel}.\\[0.3mm]
\textbf{SQL} (Structured Query Language):\\
\begin{tabular}{@{}l@{\quad}l@{}}
\code{SELECT [DISTINCT] a1, a2} & welche Spalten (ohne Duplikate)\\
\code{FROM tabelle} & aus welcher Tabelle\\
\code{WHERE bedingung} & welche Zeilen\\
\code{ORDER BY a [ASC|DESC]} & auf- bzw. absteigend sortieren\\
\end{tabular}\\[0.3mm]
Beispiel: Schüler:innen der 9a mit Geburtstag im Oktober:
\begin{java}[style=klein,language=SQL,aboveskip=0.3mm]
SELECT Nachname, Vorname, Geburtstag FROM Schueler
WHERE Klasse = '9a' AND Geburtstag LIKE '%-10-%'
ORDER BY Geburtstag ASC;
\end{java}

\gw{0.7 Graphen}
Ein Graph besteht aus \textbf{Knoten} und \textbf{Kanten}. \textbf{Pfad}: Folge durch Kanten verbundener Knoten, keine Kante doppelt. \textbf{Zyklus}: Pfad mit gleichem Start- und Endknoten ($\to$ zyklischer/azyklischer Graph). \textbf{Gerichtet}: Pfeile (eine Richtung), \textbf{ungerichtet}: Striche (beide Richtungen). \textbf{Gewichtet}: Kanten tragen Werte (z.\,B. Entfernung). \textbf{Länge eines Pfades}: Anzahl der Kanten bzw. Summe der Gewichte. Algorithmen: Breiten-, Tiefensuche, Dijkstra.

\gw{0.8 Client-Server-Modell}
Als Graph gerichtet und sternförmig mit dem \textbf{Server} in der Mitte: Jeder \textbf{Client} stellt eine \textbf{Anfrage}, der Server schickt eine \textbf{Antwort} (Zyklus der Länge 2). Protokolle: HTTP(S), DNS, SMTP, IMAP, FTP; darunter TCP/IP. Adressierung über die \textbf{IP-Adresse}, auf Hardwareseite über die \textbf{MAC-Adresse}.

\gw{0.9 Codierung und Verschlüsselung}
\textbf{Codierung}: Information nach festen, öffentlichen Regeln anders darstellen (Binärzahl, ASCII, Morsecode, QR-Code). \textbf{Verschlüsselung}: Nur wer den \textbf{Schlüssel} kennt, kann den Klartext lesen. Das Verfahren darf bekannt sein; es muss so viele Schlüssel geben, dass Ausprobieren (Brute Force) aussichtslos ist. \textbf{Symmetrisch}: ein Schlüssel zum Ver- und Entschlüsseln (Caesar, Vigenère, AES); Problem: Schlüsselaustausch. \textbf{Asymmetrisch}: \textbf{öffentlicher} Schlüssel verschlüsselt, nur der \textbf{private} entschlüsselt (RSA).

\gw{0.10 Künstliche Intelligenz und Maschinelles Lernen}
Informatik $\supset$ Data Science $\supset$ KI $\supset$ Maschinelles Lernen $\supset$ Deep Learning.
\textbf{Schwache KI} löst eng begrenzte Aufgaben (Bilderkennung, Chatbot) ohne Bewusstsein, wie alle heutigen Systeme. \textbf{Starke KI} wäre allgemein intelligent wie ein Mensch (gibt es nicht).\\[0.3mm]
\begin{minipage}[c]{0.37\linewidth}
\begin{tikzpicture}[font=\scriptsize,x=1cm,y=1cm]
\node[draw,circle,inner sep=1pt] (n) at (1.3,0) {$\Sigma\,|\,f$};
\node (x1) at (0,0.4) {$x_1$}; \node (x2) at (0,-0.4) {$x_2$}; \node (b) at (1.3,-0.8) {$b$};
\draw[-{Stealth}] (x1) -- node[above,pos=0.4] {$w_1$} (n); \draw[-{Stealth}] (x2) -- node[below,pos=0.4] {$w_2$} (n);
\draw[-{Stealth}] (b) -- (n); \draw[-{Stealth}] (n) -- ++(0.8,0) node[right] {$y$};
\end{tikzpicture}
\end{minipage}\hfill
\begin{minipage}[c]{0.62\linewidth}
\textbf{Neuron}: Eingaben $x_i$ mal \textbf{Gewichte} $w_i$, aufsummiert, plus \textbf{Bias} $b$; die \textbf{Aktivierungsfunktion} $f$ (z.\,B. Schwellenwert) liefert die Ausgabe $y$.
\end{minipage}\\[0.3mm]
\textbf{Überwachtes} Lernen: aus Beispielen mit Lösung (Label), z.\,B. Spamfilter. \textbf{Unüberwachtes}: findet Muster/Gruppen in Daten ohne Label. \textbf{Bestärkendes}: Agent lernt durch Belohnung und Bestrafung.
\end{multicols}

\newpage
\teil{Teil B: Neu in Jahrgangsstufe 12 – kurze Übersichtsaufgaben}
\footnotesize\setlength{\parskip}{1.5pt}
\begin{aufgabe}[Rekursion]{1}
\begin{minipage}[t]{0.35\linewidth}
\vspace{0pt}
\begin{java}[style=klein,aboveskip=0mm,belowskip=0mm]
public int summe(int n) {
   if (n == 0) {
      return 0;             // §\ding{172}§
   }
   return n + summe(n - 1); // §\ding{173}§
}
\end{java}
\end{minipage}\hfill
\begin{minipage}[t]{0.63\linewidth}
\vspace{0pt}
\textbf{a)} Benenne die beiden Bestandteile jeder rekursiven Methode:\\[0.5mm]
\ding{172} \lsg[4.2cm]{Rekursionsanfang (Abbruch)}\enspace \ding{173} \lsg[5.4cm]{rekursiver Aufruf, kleineres Problem}\\[0.5mm]
\textbf{b)} Was passiert, wenn \ding{172} fehlt? \lsg[6.6cm]{endlos viele Aufrufe $\to$ StackOverflowError}\\[0.5mm]
\textbf{c)} Berechne: \code{summe(3)} = \lsg[7.7cm]{3 + summe(2) = … = 3 + 2 + 1 + summe(0) = 6}
\end{minipage}
\end{aufgabe}

\begin{aufgabe}[Rekursive Datenstrukturen]{2}
\textbf{a)} Ergänze die Übersicht.\\[0.5mm]
\renewcommand{\arraystretch}{1.15}
\begin{tabularx}{\linewidth}{|p{2.9cm}|X|X|X|}\hline
 & \textbf{Warteschlange (Queue)} & \textbf{Liste} & \textbf{Binärbaum}\\\hline
Nachfolger je Knoten\rule[-2mm]{0pt}{6.2mm} & genau einer (am Ende \code{null}) & \lsgoder{\lsgfarbe genau einer}{} & \lsgoder{\lsgfarbe bis zu zwei (links, rechts)}{}\\\hline
Einfügen/Entnehmen\rule[-2mm]{0pt}{6.2mm} & \code{add} hinten, \code{poll} vorne: \lsg[1cm]{FIFO} & \lsgoder{\lsgfarbe an beliebiger Stelle}{} & \lsgoder{\lsgfarbe links kleiner, rechts größer}{}\\\hline
\end{tabularx}\\[0.5mm]
\begin{minipage}[t]{0.75\linewidth}
\vspace{0pt}
\textbf{b)} Beim Entwurfsmuster \textbf{Kompositum} erben \code{Knoten} und \code{Abschluss} von der abstrakten Klasse \code{Listenelement}. Welche Aufgabe hat der \code{Abschluss}, welchen Vorteil bringt er?
\lsglinien{2}{Er ersetzt \code{null} am Ende: Jeder Knoten hat immer einen Nachfolger und ruft die Methode dort auf; der Abschluss beendet die Rekursion. Die Abfragen auf \code{null} entfallen.}
\textbf{c)} Füge \textbf{45} in den Suchbaum ein. \textbf{Inorder}-Reihenfolge (links -- Wurzel -- rechts):\\[0.5mm] \lsg[\linewidth]{20, 30, 40, 45, 50, 60, 70 (sortiert!)}\\[0.8mm]
\textbf{d)} Warum findet man Daten im Suchbaum schneller als in einer Liste?\\[0.5mm] \lsg[\linewidth]{Jeder Vergleich halbiert etwa den Suchbereich: ca. $\log_2 n$ statt $n$ Schritte.}
\end{minipage}\hfill
\begin{minipage}[t]{0.23\linewidth}
\vspace{1mm}\centering
\begin{tikzpicture}[level distance=9mm,level 1/.style={sibling distance=19mm},level 2/.style={sibling distance=9.5mm},every node/.style={draw,circle,inner sep=0.5pt,minimum size=5.5mm,font=\scriptsize}]
\node {50}
  child {node {30} child {node {20}} child {node (v) {40}}}
  child {node {70} child {node {60}} child[missing]};
\node[lsgbild,draw,circle,minimum size=5.5mm,font=\scriptsize,below right=4mm and 0mm of v] (n) {45};
\draw[lsgbild] (v) -- (n);
\end{tikzpicture}
\end{minipage}
\end{aufgabe}

\begin{aufgabe}[Nebenläufigkeit]{3}
\begin{minipage}[t]{0.42\linewidth}
\vspace{0pt}
Zwei Kassen (Threads) buchen \textbf{gleichzeitig} ab. Guthaben: 50\,€; Kasse A bucht 5\,€ ab, Kasse B 10\,€.
\begin{java}[style=klein,belowskip=0mm]
public void abbuchen(int betrag) {
   int neu = guthaben - betrag;
   guthaben = neu;
}
\end{java}
\end{minipage}\hfill
\begin{minipage}[t]{0.55\linewidth}
\vspace{0pt}
\textbf{a)} Erkläre, wie am Ende 40\,€ statt 35\,€ auf dem Konto stehen können.
\lsglinien{2}{A und B lesen beide 50\,€. A schreibt 45\,€, danach überschreibt B mit 40\,€: Die Abbuchung von A geht verloren (\emph{Lost Update}, Race Condition).}
\end{minipage}\\[1mm]
\textbf{b)} Der Codeabschnitt, in dem auf eine gemeinsam genutzte Ressource zugegriffen wird, heißt \lsg[3.3cm]{kritischer Abschnitt}. Darin darf sich immer nur \lsg[1.9cm]{ein Thread} befinden (\lsg[3.7cm]{gegenseitiger Ausschluss}); in Java mit dem Schlüsselwort \lsg[2.2cm]{synchronized}. Wartet jeder von zwei Prozessen auf ein Betriebsmittel, das der andere belegt, blockieren beide für immer: \lsg[3.7cm]{Verklemmung (Deadlock)}.
\end{aufgabe}

\begin{aufgabe}[Softwareengineering]{4}
\textbf{a) Wasserfallmodell}: Nummeriere die Phasen:\enspace
\lsg[0.45cm]{\,4} Test\enspace \lsg[0.45cm]{\,1} Analyse\enspace \lsg[0.45cm]{\,5} Wartung\enspace \lsg[0.45cm]{\,3} Implementierung\enspace \lsg[0.45cm]{\,2} Entwurf\\[0.8mm]
Nachteil: \lsg[15.8cm]{Jede Phase wird nur einmal durchlaufen: Fehler und Änderungswünsche fallen erst spät auf und sind dann teuer.}\\[0.8mm]
\textbf{b) Agile Methoden (Scrum)}: Ein \textbf{Sprint} ist ein \lsg[5.3cm]{kurzer, fester Zeitraum (1–4 Wochen)}. Ergebnis: \lsg[4cm]{lauffähiges Teilprodukt}\\[0.8mm]
Vorteil gegenüber dem Wasserfallmodell: \lsg[11.2cm]{frühe Rückmeldung der Kunden; Anforderungen dürfen sich ändern}\\[0.8mm]
\textbf{c)} Erkläre kurz:\\[0.3mm]
\renewcommand{\arraystretch}{1.12}
\begin{tabularx}{\linewidth}{|p{2.3cm}|X|}\hline
Modularisierung\rule[-3.2mm]{0pt}{7.5mm} & \lsgoder{\lsgfarbe Zerlegung in unabhängige Teile mit klarer Aufgabe: Arbeitsteilung, einzeln testbar, austauschbar}{}\\\hline
Schnittstelle\rule[-3.2mm]{0pt}{7.5mm} & \lsgoder{\lsgfarbe legt fest, welche Methoden (Signaturen) ein Modul anbietet, ohne die Umsetzung; in Java \code{interface}}{}\\\hline
MVC\rule[-3.2mm]{0pt}{7.5mm} & \lsgoder{\lsgfarbe \textbf{Model}: Daten und Logik, \textbf{View}: Anzeige, \textbf{Controller}: verarbeitet Eingaben, steuert Model und View}{}\\\hline
\end{tabularx}
\end{aufgabe}

\begin{aufgabe}[Informationssicherheit]{5}
Welches \textbf{Schutzziel} (Vertraulichkeit, Integrität, Verfügbarkeit, Authentizität) ist verletzt? Nenne eine Maßnahme.\\[0.5mm]
\renewcommand{\arraystretch}{1.12}
\begin{tabularx}{\linewidth}{|p{7.6cm}|p{2.4cm}|X|}\hline
\textbf{Vorfall} & \textbf{Schutzziel} & \textbf{Maßnahme}\\\hline
\rule[-2mm]{0pt}{6.2mm}Ein Bagger durchtrennt die Internetleitung der Schule. & \lsgoder{\lsgfarbe Verfügbarkeit}{} & \lsgoder{\lsgfarbe zweite Leitung, Backups}{}\\\hline
\rule[-2mm]{0pt}{6.2mm}Jemand ändert heimlich Noten in der Datenbank. & \lsgoder{\lsgfarbe Integrität}{} & \lsgoder{\lsgfarbe Zugriffsrechte, Protokollierung}{}\\\hline
\rule[-2mm]{0pt}{6.2mm}Eine Phishing-Mail gibt sich als Schulleitung aus. & \lsgoder{\lsgfarbe Authentizität}{} & \lsgoder{\lsgfarbe digitale Signatur, Absender prüfen}{}\\\hline
\rule[-2mm]{0pt}{6.2mm}Fremde lesen Schülerdaten auf einem offenen Laptop. & \lsgoder{\lsgfarbe Vertraulichkeit}{} & \lsgoder{\lsgfarbe Bildschirmsperre, Verschlüsselung}{}\\\hline
\end{tabularx}
\end{aufgabe}
\end{document}
"
%           (füllt alle Lücken rot aus; siehe \lsg, \lsglinien, \lsgoder)
% =====================================================================
\usepackage[a4paper,left=18mm,right=18mm,top=26mm,bottom=17mm,headheight=15mm,headsep=4mm,footskip=8mm]{geometry}
\usepackage{iftex}
\ifPDFTeX
  \usepackage[T1]{fontenc}
  \usepackage[utf8]{inputenc}
  \usepackage{helvet}
  \usepackage[scaled=0.86]{beramono}
  \renewcommand{\familydefault}{\sfdefault}
\else
  \usepackage{fontspec}
  \setmainfont{texgyreheros}[Extension=.otf,UprightFont=*-regular,BoldFont=*-bold,ItalicFont=*-italic,BoldItalicFont=*-bolditalic]
  \setmonofont{DejaVuSansMono}[Extension=.ttf,UprightFont=*,BoldFont=*-Bold,ItalicFont=*-Oblique,BoldItalicFont=*-BoldOblique,Scale=0.86]
\fi
\usepackage[ngerman,shorthands=off]{babel}
\usepackage{xcolor,graphicx,tabularx,array,enumitem,fancyhdr,tikz,listings,multicol,calc,etoolbox,pifont}
\usepackage[most]{tcolorbox}
\usetikzlibrary{arrows.meta,positioning,calc,decorations.pathreplacing,shapes.multipart}
\usepackage[hidelinks]{hyperref}
\setlength{\parindent}{0pt}
\setlength{\parskip}{3pt}
\setlist{nosep,leftmargin=*}

% ---------- Lösungsschalter ----------
\newif\ifloesung
\ifdefined\mitloesung\loesungtrue\fi
\newcommand{\lsgfarbe}{\color{red!75!black}}

% ---------- Kopf- und Fußzeile (einheitliche Form) ----------
\newcommand{\lehrkraft}{Hau}
\newcommand{\themenbereich}{Grundwissen 6 bis 12}
\newcommand{\abnummer}{}
\pagestyle{fancy}\fancyhf{}
\renewcommand{\headrulewidth}{0.4pt}
\fancyhead[L]{\small Klasse 13\\Datum:}
\fancyhead[C]{\small Themenbereich:\\\themenbereich}
\fancyhead[R]{\small \lehrkraft\ifloesung\\{\lsgfarbe\bfseries Lösung}\fi}
\fancyfoot[C]{\scriptsize edu-mrh.de\,·\,Informatik 13\,·\,Blatt \abnummer\ifloesung\ (Lösung)\fi\,·\,Seite \thepage}
\newcommand{\abtitel}[2]{\renewcommand{\abnummer}{#1}\begin{center}\Large\bfseries #1\quad\underline{#2}\end{center}\vspace{-1mm}}

% ---------- Boxen ----------
\newenvironment{aufgabe}[2][]{\begin{tcolorbox}[enhanced,breakable,colback=white,colframe=black,boxrule=0.7pt,arc=3mm,left=2mm,right=2mm,top=1.5mm,bottom=1.5mm,before skip=2.5mm,after skip=2.5mm]\textbf{Aufgabe #2)}\ifstrempty{#1}{}{ \textit{#1}}\par\vspace{0.6mm}}{\end{tcolorbox}}
\newtcolorbox{merke}[1][Merke]{enhanced,breakable,colback=green!5,colframe=green!40!black,boxrule=0.8pt,arc=2mm,left=2mm,right=2mm,top=1.5mm,bottom=1.5mm,before skip=2.5mm,after skip=2.5mm,title={#1},fonttitle=\bfseries,coltitle=white}
\newtcolorbox{hinweis}{enhanced,breakable,colback=yellow!12,colframe=orange!70!black,boxrule=0.6pt,arc=2mm,left=2mm,right=2mm,top=1mm,bottom=1mm,before skip=2mm,after skip=2mm}
\newcommand{\web}[1]{{\small\color{gray}\textit{Website: #1}}}

% ---------- Lücken, Linien, Karos ----------
\newcommand{\luecke}[1][3.5cm]{\rule[-0.5ex]{#1}{0.4pt}}
\newcommand{\linien}[1]{\par\vspace{1.5mm}\foreach \i in {1,...,#1}{\noindent\rule[0pt]{\linewidth}{0.3pt}\par\vspace{5.5mm}}\vspace{-3mm}}
\newcommand{\kariert}[1]{\par\vspace{1.5mm}\noindent\begin{tikzpicture}\draw[step=5mm,gray!55,thin] (0,0) grid (\linewidth-1pt,#1*5mm);\end{tikzpicture}\par}
\newcommand{\karobox}[2]{\begin{tikzpicture}\draw[step=5mm,gray!55,thin] (0,0) grid (#1,#2);\end{tikzpicture}}
\newcommand{\klassenkarte}[3]{\begin{tabular}{|l|}\hline\multicolumn{1}{|c|}{\textbf{#1}}\\\hline #2\\\hline #3\\\hline\end{tabular}}
\newcommand{\code}[1]{\texttt{#1}}

% Lücke mit Lösung: \lsg[Mindestbreite]{Lösungstext}
\newlength{\lsgbreite}
\newcommand{\lsg}[2][3.5cm]{\ifloesung\settowidth{\lsgbreite}{\,#2\,}\ifdim\lsgbreite<#1\setlength{\lsgbreite}{#1}\fi\underline{\makebox[\lsgbreite][l]{\lsgfarbe\,#2}}\else\luecke[#1]\fi}
% Schreiblinien mit Lösung: \lsglinien{Anzahl}{Lösungstext} (gleiche Höhe wie \linien)
\newcommand{\lsglinien}[2]{\ifloesung\par\vspace{1.5mm}\noindent\begin{minipage}[t][#1\dimexpr 6.9mm\relax][t]{\linewidth}\lsgfarbe #2\end{minipage}\par\vspace{-1.5mm}\else\linien{#1}\fi}
% Beliebiger Inhalt: \lsgoder{nur in der Lösung}{nur im Blatt}
\newcommand{\lsgoder}[2]{\ifloesung\leavevmode#1\else#2\fi}

% ---------- Verkettete Strukturen zeichnen ----------
% Objekt mit Kopf und Attributen:  \node[obj] (k1) {k1: Knoten \nodepart{second} daten = Hans\\ nachfolger};
\tikzset{
  obj/.style={draw,rectangle split,rectangle split parts=2,rectangle split part fill={orange!25,orange!8},
              rounded corners=2pt,font=\scriptsize,inner sep=2pt,align=left},
  qobj/.style={obj,rectangle split part fill={teal!25,teal!8}},
  pobj/.style={draw,rounded corners=2pt,fill=gray!8,font=\scriptsize,inner sep=2pt,align=left},
  zeiger/.style={thick,-{Stealth[length=2.2mm]}},
  lsgzeiger/.style={zeiger,red!75!black},
  nullref/.style={font=\scriptsize\ttfamily,inner sep=1pt},
  schritt/.style={circle,draw,fill=white,inner sep=0.6pt,font=\scriptsize\bfseries},
  lsgschritt/.style={schritt,draw=red!75!black,text=red!75!black},
}

% ---------- Java-Listings ----------
\lstset{language=Java,basicstyle=\ttfamily\small,keywordstyle=\bfseries,commentstyle=\color{gray!80!black},stringstyle=\color{black},showstringspaces=false,columns=flexible,keepspaces=true,tabsize=3,frame=single,framerule=0.4pt,rulecolor=\color{black},xleftmargin=2mm,framexleftmargin=1mm,aboveskip=2mm,belowskip=2mm,breaklines=true,escapechar=§,
 literate={ä}{{\"a}}1 {ö}{{\"o}}1 {ü}{{\"u}}1 {Ä}{{\"A}}1 {Ö}{{\"O}}1 {Ü}{{\"U}}1 {ß}{{\ss}}1 {–}{{--}}1 {„}{{\glqq}}1 {“}{{\grqq}}1 {→}{{$\to$}}1 {①}{{\ding{172}}}1 {②}{{\ding{173}}}1 {③}{{\ding{174}}}1 {④}{{\ding{175}}}1}
\lstnewenvironment{java}[1][]{\lstset{#1}}{}

%  Übersetzen mit XeLaTeX, LuaLaTeX, pdfLaTeX oder Tectonic.
%
%  Blatt:   xelatex ab-0-grundwissen-13.tex
%  Lösung:  xelatex -jobname=ab-0-grundwissen-13-lsg "\def\mitloesung{}\documentclass[11pt,a4paper]{article}
% =====================================================================
%  Gemeinsame Vorlage der Arbeitsblätter – Grundwissen Q13
%  – Informatik 12 (gleiches Layout wie informatik/10/java/arbeitsblaetter/
%  und informatik/11/graphen/arbeitsblaetter/ab-vorlage.tex)
%  Einbinden in jedem Blatt direkt nach \documentclass: \input{ab-vorlage}
%  Übersetzen mit XeLaTeX, LuaLaTeX, pdfLaTeX oder Tectonic.
%
%  Blatt:   xelatex ab-0-grundwissen-13.tex
%  Lösung:  xelatex -jobname=ab-0-grundwissen-13-lsg "\def\mitloesung{}\input{ab-0-grundwissen-13}"
%           (füllt alle Lücken rot aus; siehe \lsg, \lsglinien, \lsgoder)
% =====================================================================
\usepackage[a4paper,left=18mm,right=18mm,top=26mm,bottom=17mm,headheight=15mm,headsep=4mm,footskip=8mm]{geometry}
\usepackage{iftex}
\ifPDFTeX
  \usepackage[T1]{fontenc}
  \usepackage[utf8]{inputenc}
  \usepackage{helvet}
  \usepackage[scaled=0.86]{beramono}
  \renewcommand{\familydefault}{\sfdefault}
\else
  \usepackage{fontspec}
  \setmainfont{texgyreheros}[Extension=.otf,UprightFont=*-regular,BoldFont=*-bold,ItalicFont=*-italic,BoldItalicFont=*-bolditalic]
  \setmonofont{DejaVuSansMono}[Extension=.ttf,UprightFont=*,BoldFont=*-Bold,ItalicFont=*-Oblique,BoldItalicFont=*-BoldOblique,Scale=0.86]
\fi
\usepackage[ngerman,shorthands=off]{babel}
\usepackage{xcolor,graphicx,tabularx,array,enumitem,fancyhdr,tikz,listings,multicol,calc,etoolbox,pifont}
\usepackage[most]{tcolorbox}
\usetikzlibrary{arrows.meta,positioning,calc,decorations.pathreplacing,shapes.multipart}
\usepackage[hidelinks]{hyperref}
\setlength{\parindent}{0pt}
\setlength{\parskip}{3pt}
\setlist{nosep,leftmargin=*}

% ---------- Lösungsschalter ----------
\newif\ifloesung
\ifdefined\mitloesung\loesungtrue\fi
\newcommand{\lsgfarbe}{\color{red!75!black}}

% ---------- Kopf- und Fußzeile (einheitliche Form) ----------
\newcommand{\lehrkraft}{Hau}
\newcommand{\themenbereich}{Grundwissen 6 bis 12}
\newcommand{\abnummer}{}
\pagestyle{fancy}\fancyhf{}
\renewcommand{\headrulewidth}{0.4pt}
\fancyhead[L]{\small Klasse 13\\Datum:}
\fancyhead[C]{\small Themenbereich:\\\themenbereich}
\fancyhead[R]{\small \lehrkraft\ifloesung\\{\lsgfarbe\bfseries Lösung}\fi}
\fancyfoot[C]{\scriptsize edu-mrh.de\,·\,Informatik 13\,·\,Blatt \abnummer\ifloesung\ (Lösung)\fi\,·\,Seite \thepage}
\newcommand{\abtitel}[2]{\renewcommand{\abnummer}{#1}\begin{center}\Large\bfseries #1\quad\underline{#2}\end{center}\vspace{-1mm}}

% ---------- Boxen ----------
\newenvironment{aufgabe}[2][]{\begin{tcolorbox}[enhanced,breakable,colback=white,colframe=black,boxrule=0.7pt,arc=3mm,left=2mm,right=2mm,top=1.5mm,bottom=1.5mm,before skip=2.5mm,after skip=2.5mm]\textbf{Aufgabe #2)}\ifstrempty{#1}{}{ \textit{#1}}\par\vspace{0.6mm}}{\end{tcolorbox}}
\newtcolorbox{merke}[1][Merke]{enhanced,breakable,colback=green!5,colframe=green!40!black,boxrule=0.8pt,arc=2mm,left=2mm,right=2mm,top=1.5mm,bottom=1.5mm,before skip=2.5mm,after skip=2.5mm,title={#1},fonttitle=\bfseries,coltitle=white}
\newtcolorbox{hinweis}{enhanced,breakable,colback=yellow!12,colframe=orange!70!black,boxrule=0.6pt,arc=2mm,left=2mm,right=2mm,top=1mm,bottom=1mm,before skip=2mm,after skip=2mm}
\newcommand{\web}[1]{{\small\color{gray}\textit{Website: #1}}}

% ---------- Lücken, Linien, Karos ----------
\newcommand{\luecke}[1][3.5cm]{\rule[-0.5ex]{#1}{0.4pt}}
\newcommand{\linien}[1]{\par\vspace{1.5mm}\foreach \i in {1,...,#1}{\noindent\rule[0pt]{\linewidth}{0.3pt}\par\vspace{5.5mm}}\vspace{-3mm}}
\newcommand{\kariert}[1]{\par\vspace{1.5mm}\noindent\begin{tikzpicture}\draw[step=5mm,gray!55,thin] (0,0) grid (\linewidth-1pt,#1*5mm);\end{tikzpicture}\par}
\newcommand{\karobox}[2]{\begin{tikzpicture}\draw[step=5mm,gray!55,thin] (0,0) grid (#1,#2);\end{tikzpicture}}
\newcommand{\klassenkarte}[3]{\begin{tabular}{|l|}\hline\multicolumn{1}{|c|}{\textbf{#1}}\\\hline #2\\\hline #3\\\hline\end{tabular}}
\newcommand{\code}[1]{\texttt{#1}}

% Lücke mit Lösung: \lsg[Mindestbreite]{Lösungstext}
\newlength{\lsgbreite}
\newcommand{\lsg}[2][3.5cm]{\ifloesung\settowidth{\lsgbreite}{\,#2\,}\ifdim\lsgbreite<#1\setlength{\lsgbreite}{#1}\fi\underline{\makebox[\lsgbreite][l]{\lsgfarbe\,#2}}\else\luecke[#1]\fi}
% Schreiblinien mit Lösung: \lsglinien{Anzahl}{Lösungstext} (gleiche Höhe wie \linien)
\newcommand{\lsglinien}[2]{\ifloesung\par\vspace{1.5mm}\noindent\begin{minipage}[t][#1\dimexpr 6.9mm\relax][t]{\linewidth}\lsgfarbe #2\end{minipage}\par\vspace{-1.5mm}\else\linien{#1}\fi}
% Beliebiger Inhalt: \lsgoder{nur in der Lösung}{nur im Blatt}
\newcommand{\lsgoder}[2]{\ifloesung\leavevmode#1\else#2\fi}

% ---------- Verkettete Strukturen zeichnen ----------
% Objekt mit Kopf und Attributen:  \node[obj] (k1) {k1: Knoten \nodepart{second} daten = Hans\\ nachfolger};
\tikzset{
  obj/.style={draw,rectangle split,rectangle split parts=2,rectangle split part fill={orange!25,orange!8},
              rounded corners=2pt,font=\scriptsize,inner sep=2pt,align=left},
  qobj/.style={obj,rectangle split part fill={teal!25,teal!8}},
  pobj/.style={draw,rounded corners=2pt,fill=gray!8,font=\scriptsize,inner sep=2pt,align=left},
  zeiger/.style={thick,-{Stealth[length=2.2mm]}},
  lsgzeiger/.style={zeiger,red!75!black},
  nullref/.style={font=\scriptsize\ttfamily,inner sep=1pt},
  schritt/.style={circle,draw,fill=white,inner sep=0.6pt,font=\scriptsize\bfseries},
  lsgschritt/.style={schritt,draw=red!75!black,text=red!75!black},
}

% ---------- Java-Listings ----------
\lstset{language=Java,basicstyle=\ttfamily\small,keywordstyle=\bfseries,commentstyle=\color{gray!80!black},stringstyle=\color{black},showstringspaces=false,columns=flexible,keepspaces=true,tabsize=3,frame=single,framerule=0.4pt,rulecolor=\color{black},xleftmargin=2mm,framexleftmargin=1mm,aboveskip=2mm,belowskip=2mm,breaklines=true,escapechar=§,
 literate={ä}{{\"a}}1 {ö}{{\"o}}1 {ü}{{\"u}}1 {Ä}{{\"A}}1 {Ö}{{\"O}}1 {Ü}{{\"U}}1 {ß}{{\ss}}1 {–}{{--}}1 {„}{{\glqq}}1 {“}{{\grqq}}1 {→}{{$\to$}}1 {①}{{\ding{172}}}1 {②}{{\ding{173}}}1 {③}{{\ding{174}}}1 {④}{{\ding{175}}}1}
\lstnewenvironment{java}[1][]{\lstset{#1}}{}

% Blatt:   tectonic ab-0-grundwissen-13.tex
% Lösung:  tectonic ab-0-grundwissen-13-lsg.tex  (Datei enthält nur: \def\mitloesung{}\documentclass[11pt,a4paper]{article}
\input{ab-vorlage}
% Blatt:   tectonic ab-0-grundwissen-13.tex
% Lösung:  tectonic ab-0-grundwissen-13-lsg.tex  (Datei enthält nur: \def\mitloesung{}\input{ab-0-grundwissen-13})
% Seite 1: Grundwissen 6 bis 11 als ausgefüllte Übersicht (identisch in Blatt und Lösung)
% Seite 2: Übersichtsaufgaben zu den Inhalten der Jahrgangsstufe 12
\usetikzlibrary{shapes.geometric}
\newgeometry{left=14mm,right=14mm,top=21mm,bottom=12mm,headheight=11mm,headsep=3mm,footskip=5mm}
\ifloesung\tikzset{lsgbild/.style={red!75!black}}\else\tikzset{lsgbild/.style={opacity=0}}\fi
\newcommand{\gw}[1]{\par\vspace{0.6mm}{\bfseries\color{green!40!black}#1}\par}
\newcommand{\teil}[1]{{\bfseries #1}\par\vspace{0.5mm}}
\setlength{\columnsep}{5mm}
\setlength{\columnseprule}{0.3pt}
\newcommand{\kopf}[1]{\multicolumn{1}{c}{\textbf{#1}}}
\renewenvironment{aufgabe}[2][]{\begin{tcolorbox}[enhanced,breakable,colback=white,colframe=black,boxrule=0.7pt,arc=3mm,left=2mm,right=2mm,top=1mm,bottom=1mm,before skip=1.8mm,after skip=1.8mm]\textbf{Aufgabe #2)}\ifstrempty{#1}{}{ \textit{#1}}\par\vspace{0.3mm}}{\end{tcolorbox}}
\lstdefinestyle{klein}{basicstyle=\ttfamily\scriptsize,aboveskip=0.8mm,belowskip=0.8mm}
\begin{document}
\abtitel{0}{Grundwissen}
\vspace{-3mm}
\teil{Teil A: Grundwissen der Jahrgangsstufen 6 bis 11 (zum Nachschlagen)}

\begin{minipage}[t]{0.575\linewidth}
\vspace{0pt}
\begin{tikzpicture}[x=1cm,y=1cm,font=\footnotesize]
\node[anchor=north west] (kk) at (0,0) {\klassenkarte{Dreieck}{seitenlaenge\\farbe}{setzeFarbe(neu)}};
\node[anchor=south west,font=\footnotesize\bfseries] at (0,0.05) {Klasse (Bauplan)};
\node[anchor=north west,draw,rounded corners=2.5mm,inner sep=0pt] (ok) at (5.35,0) {\begin{tabular}{l}\rule{0pt}{3.5mm}\textbf{form1: Dreieck}\\[0.5mm]\hline\rule{0pt}{3.5mm}seitenlaenge = 20\\farbe = "blau"\\[0.5mm]\end{tabular}};
\node[anchor=south west,font=\footnotesize\bfseries] at (5.35,0.05) {Objekt (Instanz)};
\draw[decorate,decoration={brace,amplitude=3pt}] ([xshift=2pt]kk.north east) -- ([xshift=2pt]kk.north east|-{(0,-0.5)}) node[midway,right=4pt,font=\scriptsize] {Klassenname};
\draw[decorate,decoration={brace,amplitude=3pt}] ([xshift=2pt]kk.north east|-{(0,-0.58)}) -- ([xshift=2pt]kk.north east|-{(0,-1.35)}) node[midway,right=4pt,font=\scriptsize] {Attribute};
\draw[decorate,decoration={brace,amplitude=3pt}] ([xshift=2pt]kk.north east|-{(0,-1.43)}) -- ([xshift=2pt]kk.south east) node[midway,right=4pt,font=\scriptsize] {Methoden};
\draw[decorate,decoration={brace,amplitude=3pt}] ([xshift=2pt]ok.north east) -- ([xshift=2pt]ok.north east|-{(0,-0.55)}) node[midway,right=4pt,align=left,font=\scriptsize] {Objektname:\\Klasse};
\draw[decorate,decoration={brace,amplitude=3pt}] ([xshift=2pt]ok.north east|-{(0,-0.63)}) -- ([xshift=2pt]ok.south east) node[midway,right=4pt,align=left,font=\scriptsize] {Attribut-\\werte};
\end{tikzpicture}
\end{minipage}\hfill
\begin{minipage}[t]{0.405\linewidth}
\vspace{0pt}\footnotesize
\textbf{Objektorientierung}: Problemlösen durch Analyse der beteiligten \textbf{Objekte} und ihres Zusammenspiels. Die \textbf{Klasse} ist der Bauplan, ein \textbf{Objekt} eine konkrete Ausprägung (Instanz) der Klasse.\\[0.5mm]
\textbf{Punktnotation}: \code{objekt.methode(parameter);}\\ \code{objekt.attribut = wert;}\quad z.\,B. \code{karol.schritt();}
\end{minipage}

\vspace{-1.5mm}
\begin{multicols}{2}\footnotesize\setlength{\parskip}{0.8pt}
\gw{0.1 Beziehungen, Datentypen, Zugriffsrechte}
Objekte können \textbf{Referenzen} auf andere Objekte enthalten (\textbf{Referenzattribute}); die referenzierte Klasse wird zum \textbf{Datentyp}, z.\,B. \code{private Person eigentuemer;} in \code{Haus}:\\[0.3mm]
\begin{tikzpicture}[font=\scriptsize]
\node[draw,inner sep=2pt] (h) {\textbf{Haus}};
\node[draw,inner sep=2pt,right=17mm of h] (p) {\textbf{Person}};
\draw[-{Stealth}] (h) -- node[above,inner sep=1pt] {eigentuemer} node[below,pos=0.85,inner sep=1pt] {1} (p);
\end{tikzpicture}\quad {\scriptsize(Klassendiagramm ohne \code{String}, \code{int})}\\[0.3mm]
Weitere Datentypen: \code{boolean}, \code{char}, \code{double}, \code{float}, \code{long}, \code{byte}, Felder, jede Klasse.
Objekte kommunizieren über \textbf{Methodenaufrufe}. \textbf{Datenkapselung}: Attribute sind \code{private}, Zugriff nur über \textbf{Getter} und \textbf{Setter}.\\[0.5mm]
\renewcommand{\arraystretch}{1}
\begin{tabular}{@{}|c|l|l|@{}}\hline
\code{+} & \code{public} & öffentlich: jede Klasse darf zugreifen\\\hline
\code{\#} & \code{protected} & eigene Klasse und Unterklassen\\\hline
\code{-} & \code{private} & nur innerhalb der eigenen Klasse\\\hline
\end{tabular}

\gw{0.2 Bausteine von Algorithmen: Struktogramm und Java}
\begin{tikzpicture}[x=1cm,y=1cm,font=\scriptsize]
\draw (0,0) rectangle (3.95,1.1); \draw (0,1.1) -- (1.97,0.55) -- (3.95,1.1); \draw (0,0.55) -- (3.95,0.55); \draw (1.97,0) -- (1.97,0.55);
\node at (1.97,0.93) {\ttfamily\itshape istWand()}; \node at (0.33,0.66) {wahr}; \node at (3.55,0.66) {falsch};
\node at (0.98,0.27) {\ttfamily linksDrehen()}; \node at (2.96,0.27) {\ttfamily schritt()};
\begin{scope}[xshift=4.3cm]
\draw (0,0) rectangle (4.05,1.1); \draw (0.4,0) -- (0.4,0.75) -- (4.05,0.75);
\node[anchor=west] at (0.05,0.93) {Wiederhole solange {\ttfamily\itshape !istWand()}};
\node[anchor=west] at (0.5,0.53) {\ttfamily hinlegen()}; \node[anchor=west] at (0.5,0.2) {\ttfamily schritt()};
\end{scope}
\end{tikzpicture}\\[-1.2mm]
\begin{minipage}[t]{0.47\linewidth}
\begin{java}[style=klein]
if (istWand()) {
   linksDrehen();
} else {
   schritt();
}
\end{java}
\end{minipage}\hfill
\begin{minipage}[t]{0.49\linewidth}
\begin{java}[style=klein]
while (!istWand()) {
   hinlegen();
   schritt();
}
\end{java}
\end{minipage}\\[0.3mm]
\textbf{Sequenz}: Anweisungen untereinander; \textbf{einseitig}: \code{if} ohne \code{else}; „Wiederhole 5-mal“: \code{for (int i = 0; i < 5; i++) \{…\}}

\gw{0.3 Feld (Array)}
Folge \underline{fester Länge} von Werten bzw. Objekten \underline{desselben Typs}. Zugriff über den \textbf{Index}; die Zählung beginnt bei 0.
\begin{java}[style=klein]
int zahl[];            // Deklaration
zahl = new int[1000];  // Instanziierung (Index 0..999)
zahl[2] = 10;          // Wertzuweisung (3. Element)
\end{java}

\gw{0.4 Das EVA-Prinzip}
\begin{tikzpicture}[font=\scriptsize,node distance=2.5mm,every node/.style={align=center}]
\node[draw,minimum height=7mm] (e) {\textbf{Eingabe}\\\code{gewicht, groesse}};
\node[draw,ellipse,inner sep=0pt,right=of e] (v) {\textbf{Verarbeitung}\\\code{bmi = gewicht/}\\\code{(groesse*groesse)}};
\node[draw,minimum height=7mm,right=of v] (a) {\textbf{Ausgabe}\\\code{println(…)}};
\draw[-{Stealth}] (e) -- (v); \draw[-{Stealth}] (v) -- (a);
\end{tikzpicture}\\
Bei Methoden: Eingabe über \textbf{Parameter}, Ausgabe als \textbf{Rückgabewert} oder auf dem Bildschirm.

\gw{0.5 Vererbung, @Override und Polymorphismus}
Eine \textbf{Unterklasse} erbt Attribute und Methoden ihrer \textbf{Oberklasse}; in Java von genau \underline{einer} Klasse (keine Mehrfachvererbung).\\[0.3mm]
\begin{minipage}[t]{0.21\linewidth}
\vspace{0pt}
\begin{tikzpicture}[font=\scriptsize]
\node[draw,rectangle split,rectangle split parts=2,align=left,inner sep=2pt] (p) {\textbf{Person}\nodepart{second}alter\\name};
\node[draw,rectangle split,rectangle split parts=2,align=left,inner sep=2pt,below=4mm of p] (s) {\textbf{Schueler}\nodepart{second}klasse};
\draw[-{Triangle[open,length=2mm]}] (s) -- (p);
\end{tikzpicture}
\end{minipage}\hfill
\begin{minipage}[t]{0.78\linewidth}
\vspace{0pt}
\begin{java}[style=klein,aboveskip=0mm,belowskip=0mm]
public class Schueler extends Person {
   private String klasse;
   Schueler(int a, String n, String k) {
      super(a, n);  klasse = k;
   }
}
\end{java}
\end{minipage}\\[0.8mm]
\begin{tabular}{@{}|l|p{6.25cm}|@{}}\hline
\code{extends} & gibt die Oberklasse an\\\hline
\code{super(…)} & ruft den Konstruktor der Oberklasse auf\\\hline
\code{@Override} & Unterklasse \textbf{überschreibt} eine geerbte Methode (gleiche Signatur, neues Verhalten)\\\hline
Polymorphie & Variable vom Typ der Oberklasse kann auf Unterklassen-Objekte verweisen; ausgeführt wird die Methode der \emph{tatsächlichen} Klasse.\\\hline
\end{tabular}

\gw{0.6 Datenbanken und SQL}
Datenbanken speichern strukturierte Daten in \textbf{Tabellen}: Ein \textbf{Datensatz} (Zeile, Tupel) enthält die Attributwerte (Spalten $\to$ \textbf{Attribute}) eines Objekts. Der \textbf{Primärschlüssel} (unterstrichen; ein oder mehrere Attribute) identifiziert jeden Datensatz eindeutig. Beziehungen zwischen Tabellen: \textbf{Fremdschlüssel}.\\[0.3mm]
\textbf{SQL} (Structured Query Language):\\
\begin{tabular}{@{}l@{\quad}l@{}}
\code{SELECT [DISTINCT] a1, a2} & welche Spalten (ohne Duplikate)\\
\code{FROM tabelle} & aus welcher Tabelle\\
\code{WHERE bedingung} & welche Zeilen\\
\code{ORDER BY a [ASC|DESC]} & auf- bzw. absteigend sortieren\\
\end{tabular}\\[0.3mm]
Beispiel: Schüler:innen der 9a mit Geburtstag im Oktober:
\begin{java}[style=klein,language=SQL,aboveskip=0.3mm]
SELECT Nachname, Vorname, Geburtstag FROM Schueler
WHERE Klasse = '9a' AND Geburtstag LIKE '%-10-%'
ORDER BY Geburtstag ASC;
\end{java}

\gw{0.7 Graphen}
Ein Graph besteht aus \textbf{Knoten} und \textbf{Kanten}. \textbf{Pfad}: Folge durch Kanten verbundener Knoten, keine Kante doppelt. \textbf{Zyklus}: Pfad mit gleichem Start- und Endknoten ($\to$ zyklischer/azyklischer Graph). \textbf{Gerichtet}: Pfeile (eine Richtung), \textbf{ungerichtet}: Striche (beide Richtungen). \textbf{Gewichtet}: Kanten tragen Werte (z.\,B. Entfernung). \textbf{Länge eines Pfades}: Anzahl der Kanten bzw. Summe der Gewichte. Algorithmen: Breiten-, Tiefensuche, Dijkstra.

\gw{0.8 Client-Server-Modell}
Als Graph gerichtet und sternförmig mit dem \textbf{Server} in der Mitte: Jeder \textbf{Client} stellt eine \textbf{Anfrage}, der Server schickt eine \textbf{Antwort} (Zyklus der Länge 2). Protokolle: HTTP(S), DNS, SMTP, IMAP, FTP; darunter TCP/IP. Adressierung über die \textbf{IP-Adresse}, auf Hardwareseite über die \textbf{MAC-Adresse}.

\gw{0.9 Codierung und Verschlüsselung}
\textbf{Codierung}: Information nach festen, öffentlichen Regeln anders darstellen (Binärzahl, ASCII, Morsecode, QR-Code). \textbf{Verschlüsselung}: Nur wer den \textbf{Schlüssel} kennt, kann den Klartext lesen. Das Verfahren darf bekannt sein; es muss so viele Schlüssel geben, dass Ausprobieren (Brute Force) aussichtslos ist. \textbf{Symmetrisch}: ein Schlüssel zum Ver- und Entschlüsseln (Caesar, Vigenère, AES); Problem: Schlüsselaustausch. \textbf{Asymmetrisch}: \textbf{öffentlicher} Schlüssel verschlüsselt, nur der \textbf{private} entschlüsselt (RSA).

\gw{0.10 Künstliche Intelligenz und Maschinelles Lernen}
Informatik $\supset$ Data Science $\supset$ KI $\supset$ Maschinelles Lernen $\supset$ Deep Learning.
\textbf{Schwache KI} löst eng begrenzte Aufgaben (Bilderkennung, Chatbot) ohne Bewusstsein, wie alle heutigen Systeme. \textbf{Starke KI} wäre allgemein intelligent wie ein Mensch (gibt es nicht).\\[0.3mm]
\begin{minipage}[c]{0.37\linewidth}
\begin{tikzpicture}[font=\scriptsize,x=1cm,y=1cm]
\node[draw,circle,inner sep=1pt] (n) at (1.3,0) {$\Sigma\,|\,f$};
\node (x1) at (0,0.4) {$x_1$}; \node (x2) at (0,-0.4) {$x_2$}; \node (b) at (1.3,-0.8) {$b$};
\draw[-{Stealth}] (x1) -- node[above,pos=0.4] {$w_1$} (n); \draw[-{Stealth}] (x2) -- node[below,pos=0.4] {$w_2$} (n);
\draw[-{Stealth}] (b) -- (n); \draw[-{Stealth}] (n) -- ++(0.8,0) node[right] {$y$};
\end{tikzpicture}
\end{minipage}\hfill
\begin{minipage}[c]{0.62\linewidth}
\textbf{Neuron}: Eingaben $x_i$ mal \textbf{Gewichte} $w_i$, aufsummiert, plus \textbf{Bias} $b$; die \textbf{Aktivierungsfunktion} $f$ (z.\,B. Schwellenwert) liefert die Ausgabe $y$.
\end{minipage}\\[0.3mm]
\textbf{Überwachtes} Lernen: aus Beispielen mit Lösung (Label), z.\,B. Spamfilter. \textbf{Unüberwachtes}: findet Muster/Gruppen in Daten ohne Label. \textbf{Bestärkendes}: Agent lernt durch Belohnung und Bestrafung.
\end{multicols}

\newpage
\teil{Teil B: Neu in Jahrgangsstufe 12 – kurze Übersichtsaufgaben}
\footnotesize\setlength{\parskip}{1.5pt}
\begin{aufgabe}[Rekursion]{1}
\begin{minipage}[t]{0.35\linewidth}
\vspace{0pt}
\begin{java}[style=klein,aboveskip=0mm,belowskip=0mm]
public int summe(int n) {
   if (n == 0) {
      return 0;             // §\ding{172}§
   }
   return n + summe(n - 1); // §\ding{173}§
}
\end{java}
\end{minipage}\hfill
\begin{minipage}[t]{0.63\linewidth}
\vspace{0pt}
\textbf{a)} Benenne die beiden Bestandteile jeder rekursiven Methode:\\[0.5mm]
\ding{172} \lsg[4.2cm]{Rekursionsanfang (Abbruch)}\enspace \ding{173} \lsg[5.4cm]{rekursiver Aufruf, kleineres Problem}\\[0.5mm]
\textbf{b)} Was passiert, wenn \ding{172} fehlt? \lsg[6.6cm]{endlos viele Aufrufe $\to$ StackOverflowError}\\[0.5mm]
\textbf{c)} Berechne: \code{summe(3)} = \lsg[7.7cm]{3 + summe(2) = … = 3 + 2 + 1 + summe(0) = 6}
\end{minipage}
\end{aufgabe}

\begin{aufgabe}[Rekursive Datenstrukturen]{2}
\textbf{a)} Ergänze die Übersicht.\\[0.5mm]
\renewcommand{\arraystretch}{1.15}
\begin{tabularx}{\linewidth}{|p{2.9cm}|X|X|X|}\hline
 & \textbf{Warteschlange (Queue)} & \textbf{Liste} & \textbf{Binärbaum}\\\hline
Nachfolger je Knoten\rule[-2mm]{0pt}{6.2mm} & genau einer (am Ende \code{null}) & \lsgoder{\lsgfarbe genau einer}{} & \lsgoder{\lsgfarbe bis zu zwei (links, rechts)}{}\\\hline
Einfügen/Entnehmen\rule[-2mm]{0pt}{6.2mm} & \code{add} hinten, \code{poll} vorne: \lsg[1cm]{FIFO} & \lsgoder{\lsgfarbe an beliebiger Stelle}{} & \lsgoder{\lsgfarbe links kleiner, rechts größer}{}\\\hline
\end{tabularx}\\[0.5mm]
\begin{minipage}[t]{0.75\linewidth}
\vspace{0pt}
\textbf{b)} Beim Entwurfsmuster \textbf{Kompositum} erben \code{Knoten} und \code{Abschluss} von der abstrakten Klasse \code{Listenelement}. Welche Aufgabe hat der \code{Abschluss}, welchen Vorteil bringt er?
\lsglinien{2}{Er ersetzt \code{null} am Ende: Jeder Knoten hat immer einen Nachfolger und ruft die Methode dort auf; der Abschluss beendet die Rekursion. Die Abfragen auf \code{null} entfallen.}
\textbf{c)} Füge \textbf{45} in den Suchbaum ein. \textbf{Inorder}-Reihenfolge (links -- Wurzel -- rechts):\\[0.5mm] \lsg[\linewidth]{20, 30, 40, 45, 50, 60, 70 (sortiert!)}\\[0.8mm]
\textbf{d)} Warum findet man Daten im Suchbaum schneller als in einer Liste?\\[0.5mm] \lsg[\linewidth]{Jeder Vergleich halbiert etwa den Suchbereich: ca. $\log_2 n$ statt $n$ Schritte.}
\end{minipage}\hfill
\begin{minipage}[t]{0.23\linewidth}
\vspace{1mm}\centering
\begin{tikzpicture}[level distance=9mm,level 1/.style={sibling distance=19mm},level 2/.style={sibling distance=9.5mm},every node/.style={draw,circle,inner sep=0.5pt,minimum size=5.5mm,font=\scriptsize}]
\node {50}
  child {node {30} child {node {20}} child {node (v) {40}}}
  child {node {70} child {node {60}} child[missing]};
\node[lsgbild,draw,circle,minimum size=5.5mm,font=\scriptsize,below right=4mm and 0mm of v] (n) {45};
\draw[lsgbild] (v) -- (n);
\end{tikzpicture}
\end{minipage}
\end{aufgabe}

\begin{aufgabe}[Nebenläufigkeit]{3}
\begin{minipage}[t]{0.42\linewidth}
\vspace{0pt}
Zwei Kassen (Threads) buchen \textbf{gleichzeitig} ab. Guthaben: 50\,€; Kasse A bucht 5\,€ ab, Kasse B 10\,€.
\begin{java}[style=klein,belowskip=0mm]
public void abbuchen(int betrag) {
   int neu = guthaben - betrag;
   guthaben = neu;
}
\end{java}
\end{minipage}\hfill
\begin{minipage}[t]{0.55\linewidth}
\vspace{0pt}
\textbf{a)} Erkläre, wie am Ende 40\,€ statt 35\,€ auf dem Konto stehen können.
\lsglinien{2}{A und B lesen beide 50\,€. A schreibt 45\,€, danach überschreibt B mit 40\,€: Die Abbuchung von A geht verloren (\emph{Lost Update}, Race Condition).}
\end{minipage}\\[1mm]
\textbf{b)} Der Codeabschnitt, in dem auf eine gemeinsam genutzte Ressource zugegriffen wird, heißt \lsg[3.3cm]{kritischer Abschnitt}. Darin darf sich immer nur \lsg[1.9cm]{ein Thread} befinden (\lsg[3.7cm]{gegenseitiger Ausschluss}); in Java mit dem Schlüsselwort \lsg[2.2cm]{synchronized}. Wartet jeder von zwei Prozessen auf ein Betriebsmittel, das der andere belegt, blockieren beide für immer: \lsg[3.7cm]{Verklemmung (Deadlock)}.
\end{aufgabe}

\begin{aufgabe}[Softwareengineering]{4}
\textbf{a) Wasserfallmodell}: Nummeriere die Phasen:\enspace
\lsg[0.45cm]{\,4} Test\enspace \lsg[0.45cm]{\,1} Analyse\enspace \lsg[0.45cm]{\,5} Wartung\enspace \lsg[0.45cm]{\,3} Implementierung\enspace \lsg[0.45cm]{\,2} Entwurf\\[0.8mm]
Nachteil: \lsg[15.8cm]{Jede Phase wird nur einmal durchlaufen: Fehler und Änderungswünsche fallen erst spät auf und sind dann teuer.}\\[0.8mm]
\textbf{b) Agile Methoden (Scrum)}: Ein \textbf{Sprint} ist ein \lsg[5.3cm]{kurzer, fester Zeitraum (1–4 Wochen)}. Ergebnis: \lsg[4cm]{lauffähiges Teilprodukt}\\[0.8mm]
Vorteil gegenüber dem Wasserfallmodell: \lsg[11.2cm]{frühe Rückmeldung der Kunden; Anforderungen dürfen sich ändern}\\[0.8mm]
\textbf{c)} Erkläre kurz:\\[0.3mm]
\renewcommand{\arraystretch}{1.12}
\begin{tabularx}{\linewidth}{|p{2.3cm}|X|}\hline
Modularisierung\rule[-3.2mm]{0pt}{7.5mm} & \lsgoder{\lsgfarbe Zerlegung in unabhängige Teile mit klarer Aufgabe: Arbeitsteilung, einzeln testbar, austauschbar}{}\\\hline
Schnittstelle\rule[-3.2mm]{0pt}{7.5mm} & \lsgoder{\lsgfarbe legt fest, welche Methoden (Signaturen) ein Modul anbietet, ohne die Umsetzung; in Java \code{interface}}{}\\\hline
MVC\rule[-3.2mm]{0pt}{7.5mm} & \lsgoder{\lsgfarbe \textbf{Model}: Daten und Logik, \textbf{View}: Anzeige, \textbf{Controller}: verarbeitet Eingaben, steuert Model und View}{}\\\hline
\end{tabularx}
\end{aufgabe}

\begin{aufgabe}[Informationssicherheit]{5}
Welches \textbf{Schutzziel} (Vertraulichkeit, Integrität, Verfügbarkeit, Authentizität) ist verletzt? Nenne eine Maßnahme.\\[0.5mm]
\renewcommand{\arraystretch}{1.12}
\begin{tabularx}{\linewidth}{|p{7.6cm}|p{2.4cm}|X|}\hline
\textbf{Vorfall} & \textbf{Schutzziel} & \textbf{Maßnahme}\\\hline
\rule[-2mm]{0pt}{6.2mm}Ein Bagger durchtrennt die Internetleitung der Schule. & \lsgoder{\lsgfarbe Verfügbarkeit}{} & \lsgoder{\lsgfarbe zweite Leitung, Backups}{}\\\hline
\rule[-2mm]{0pt}{6.2mm}Jemand ändert heimlich Noten in der Datenbank. & \lsgoder{\lsgfarbe Integrität}{} & \lsgoder{\lsgfarbe Zugriffsrechte, Protokollierung}{}\\\hline
\rule[-2mm]{0pt}{6.2mm}Eine Phishing-Mail gibt sich als Schulleitung aus. & \lsgoder{\lsgfarbe Authentizität}{} & \lsgoder{\lsgfarbe digitale Signatur, Absender prüfen}{}\\\hline
\rule[-2mm]{0pt}{6.2mm}Fremde lesen Schülerdaten auf einem offenen Laptop. & \lsgoder{\lsgfarbe Vertraulichkeit}{} & \lsgoder{\lsgfarbe Bildschirmsperre, Verschlüsselung}{}\\\hline
\end{tabularx}
\end{aufgabe}
\end{document}
)
% Seite 1: Grundwissen 6 bis 11 als ausgefüllte Übersicht (identisch in Blatt und Lösung)
% Seite 2: Übersichtsaufgaben zu den Inhalten der Jahrgangsstufe 12
\usetikzlibrary{shapes.geometric}
\newgeometry{left=14mm,right=14mm,top=21mm,bottom=12mm,headheight=11mm,headsep=3mm,footskip=5mm}
\ifloesung\tikzset{lsgbild/.style={red!75!black}}\else\tikzset{lsgbild/.style={opacity=0}}\fi
\newcommand{\gw}[1]{\par\vspace{0.6mm}{\bfseries\color{green!40!black}#1}\par}
\newcommand{\teil}[1]{{\bfseries #1}\par\vspace{0.5mm}}
\setlength{\columnsep}{5mm}
\setlength{\columnseprule}{0.3pt}
\newcommand{\kopf}[1]{\multicolumn{1}{c}{\textbf{#1}}}
\renewenvironment{aufgabe}[2][]{\begin{tcolorbox}[enhanced,breakable,colback=white,colframe=black,boxrule=0.7pt,arc=3mm,left=2mm,right=2mm,top=1mm,bottom=1mm,before skip=1.8mm,after skip=1.8mm]\textbf{Aufgabe #2)}\ifstrempty{#1}{}{ \textit{#1}}\par\vspace{0.3mm}}{\end{tcolorbox}}
\lstdefinestyle{klein}{basicstyle=\ttfamily\scriptsize,aboveskip=0.8mm,belowskip=0.8mm}
\begin{document}
\abtitel{0}{Grundwissen}
\vspace{-3mm}
\teil{Teil A: Grundwissen der Jahrgangsstufen 6 bis 11 (zum Nachschlagen)}

\begin{minipage}[t]{0.575\linewidth}
\vspace{0pt}
\begin{tikzpicture}[x=1cm,y=1cm,font=\footnotesize]
\node[anchor=north west] (kk) at (0,0) {\klassenkarte{Dreieck}{seitenlaenge\\farbe}{setzeFarbe(neu)}};
\node[anchor=south west,font=\footnotesize\bfseries] at (0,0.05) {Klasse (Bauplan)};
\node[anchor=north west,draw,rounded corners=2.5mm,inner sep=0pt] (ok) at (5.35,0) {\begin{tabular}{l}\rule{0pt}{3.5mm}\textbf{form1: Dreieck}\\[0.5mm]\hline\rule{0pt}{3.5mm}seitenlaenge = 20\\farbe = "blau"\\[0.5mm]\end{tabular}};
\node[anchor=south west,font=\footnotesize\bfseries] at (5.35,0.05) {Objekt (Instanz)};
\draw[decorate,decoration={brace,amplitude=3pt}] ([xshift=2pt]kk.north east) -- ([xshift=2pt]kk.north east|-{(0,-0.5)}) node[midway,right=4pt,font=\scriptsize] {Klassenname};
\draw[decorate,decoration={brace,amplitude=3pt}] ([xshift=2pt]kk.north east|-{(0,-0.58)}) -- ([xshift=2pt]kk.north east|-{(0,-1.35)}) node[midway,right=4pt,font=\scriptsize] {Attribute};
\draw[decorate,decoration={brace,amplitude=3pt}] ([xshift=2pt]kk.north east|-{(0,-1.43)}) -- ([xshift=2pt]kk.south east) node[midway,right=4pt,font=\scriptsize] {Methoden};
\draw[decorate,decoration={brace,amplitude=3pt}] ([xshift=2pt]ok.north east) -- ([xshift=2pt]ok.north east|-{(0,-0.55)}) node[midway,right=4pt,align=left,font=\scriptsize] {Objektname:\\Klasse};
\draw[decorate,decoration={brace,amplitude=3pt}] ([xshift=2pt]ok.north east|-{(0,-0.63)}) -- ([xshift=2pt]ok.south east) node[midway,right=4pt,align=left,font=\scriptsize] {Attribut-\\werte};
\end{tikzpicture}
\end{minipage}\hfill
\begin{minipage}[t]{0.405\linewidth}
\vspace{0pt}\footnotesize
\textbf{Objektorientierung}: Problemlösen durch Analyse der beteiligten \textbf{Objekte} und ihres Zusammenspiels. Die \textbf{Klasse} ist der Bauplan, ein \textbf{Objekt} eine konkrete Ausprägung (Instanz) der Klasse.\\[0.5mm]
\textbf{Punktnotation}: \code{objekt.methode(parameter);}\\ \code{objekt.attribut = wert;}\quad z.\,B. \code{karol.schritt();}
\end{minipage}

\vspace{-1.5mm}
\begin{multicols}{2}\footnotesize\setlength{\parskip}{0.8pt}
\gw{0.1 Beziehungen, Datentypen, Zugriffsrechte}
Objekte können \textbf{Referenzen} auf andere Objekte enthalten (\textbf{Referenzattribute}); die referenzierte Klasse wird zum \textbf{Datentyp}, z.\,B. \code{private Person eigentuemer;} in \code{Haus}:\\[0.3mm]
\begin{tikzpicture}[font=\scriptsize]
\node[draw,inner sep=2pt] (h) {\textbf{Haus}};
\node[draw,inner sep=2pt,right=17mm of h] (p) {\textbf{Person}};
\draw[-{Stealth}] (h) -- node[above,inner sep=1pt] {eigentuemer} node[below,pos=0.85,inner sep=1pt] {1} (p);
\end{tikzpicture}\quad {\scriptsize(Klassendiagramm ohne \code{String}, \code{int})}\\[0.3mm]
Weitere Datentypen: \code{boolean}, \code{char}, \code{double}, \code{float}, \code{long}, \code{byte}, Felder, jede Klasse.
Objekte kommunizieren über \textbf{Methodenaufrufe}. \textbf{Datenkapselung}: Attribute sind \code{private}, Zugriff nur über \textbf{Getter} und \textbf{Setter}.\\[0.5mm]
\renewcommand{\arraystretch}{1}
\begin{tabular}{@{}|c|l|l|@{}}\hline
\code{+} & \code{public} & öffentlich: jede Klasse darf zugreifen\\\hline
\code{\#} & \code{protected} & eigene Klasse und Unterklassen\\\hline
\code{-} & \code{private} & nur innerhalb der eigenen Klasse\\\hline
\end{tabular}

\gw{0.2 Bausteine von Algorithmen: Struktogramm und Java}
\begin{tikzpicture}[x=1cm,y=1cm,font=\scriptsize]
\draw (0,0) rectangle (3.95,1.1); \draw (0,1.1) -- (1.97,0.55) -- (3.95,1.1); \draw (0,0.55) -- (3.95,0.55); \draw (1.97,0) -- (1.97,0.55);
\node at (1.97,0.93) {\ttfamily\itshape istWand()}; \node at (0.33,0.66) {wahr}; \node at (3.55,0.66) {falsch};
\node at (0.98,0.27) {\ttfamily linksDrehen()}; \node at (2.96,0.27) {\ttfamily schritt()};
\begin{scope}[xshift=4.3cm]
\draw (0,0) rectangle (4.05,1.1); \draw (0.4,0) -- (0.4,0.75) -- (4.05,0.75);
\node[anchor=west] at (0.05,0.93) {Wiederhole solange {\ttfamily\itshape !istWand()}};
\node[anchor=west] at (0.5,0.53) {\ttfamily hinlegen()}; \node[anchor=west] at (0.5,0.2) {\ttfamily schritt()};
\end{scope}
\end{tikzpicture}\\[-1.2mm]
\begin{minipage}[t]{0.47\linewidth}
\begin{java}[style=klein]
if (istWand()) {
   linksDrehen();
} else {
   schritt();
}
\end{java}
\end{minipage}\hfill
\begin{minipage}[t]{0.49\linewidth}
\begin{java}[style=klein]
while (!istWand()) {
   hinlegen();
   schritt();
}
\end{java}
\end{minipage}\\[0.3mm]
\textbf{Sequenz}: Anweisungen untereinander; \textbf{einseitig}: \code{if} ohne \code{else}; „Wiederhole 5-mal“: \code{for (int i = 0; i < 5; i++) \{…\}}

\gw{0.3 Feld (Array)}
Folge \underline{fester Länge} von Werten bzw. Objekten \underline{desselben Typs}. Zugriff über den \textbf{Index}; die Zählung beginnt bei 0.
\begin{java}[style=klein]
int zahl[];            // Deklaration
zahl = new int[1000];  // Instanziierung (Index 0..999)
zahl[2] = 10;          // Wertzuweisung (3. Element)
\end{java}

\gw{0.4 Das EVA-Prinzip}
\begin{tikzpicture}[font=\scriptsize,node distance=2.5mm,every node/.style={align=center}]
\node[draw,minimum height=7mm] (e) {\textbf{Eingabe}\\\code{gewicht, groesse}};
\node[draw,ellipse,inner sep=0pt,right=of e] (v) {\textbf{Verarbeitung}\\\code{bmi = gewicht/}\\\code{(groesse*groesse)}};
\node[draw,minimum height=7mm,right=of v] (a) {\textbf{Ausgabe}\\\code{println(…)}};
\draw[-{Stealth}] (e) -- (v); \draw[-{Stealth}] (v) -- (a);
\end{tikzpicture}\\
Bei Methoden: Eingabe über \textbf{Parameter}, Ausgabe als \textbf{Rückgabewert} oder auf dem Bildschirm.

\gw{0.5 Vererbung, @Override und Polymorphismus}
Eine \textbf{Unterklasse} erbt Attribute und Methoden ihrer \textbf{Oberklasse}; in Java von genau \underline{einer} Klasse (keine Mehrfachvererbung).\\[0.3mm]
\begin{minipage}[t]{0.21\linewidth}
\vspace{0pt}
\begin{tikzpicture}[font=\scriptsize]
\node[draw,rectangle split,rectangle split parts=2,align=left,inner sep=2pt] (p) {\textbf{Person}\nodepart{second}alter\\name};
\node[draw,rectangle split,rectangle split parts=2,align=left,inner sep=2pt,below=4mm of p] (s) {\textbf{Schueler}\nodepart{second}klasse};
\draw[-{Triangle[open,length=2mm]}] (s) -- (p);
\end{tikzpicture}
\end{minipage}\hfill
\begin{minipage}[t]{0.78\linewidth}
\vspace{0pt}
\begin{java}[style=klein,aboveskip=0mm,belowskip=0mm]
public class Schueler extends Person {
   private String klasse;
   Schueler(int a, String n, String k) {
      super(a, n);  klasse = k;
   }
}
\end{java}
\end{minipage}\\[0.8mm]
\begin{tabular}{@{}|l|p{6.25cm}|@{}}\hline
\code{extends} & gibt die Oberklasse an\\\hline
\code{super(…)} & ruft den Konstruktor der Oberklasse auf\\\hline
\code{@Override} & Unterklasse \textbf{überschreibt} eine geerbte Methode (gleiche Signatur, neues Verhalten)\\\hline
Polymorphie & Variable vom Typ der Oberklasse kann auf Unterklassen-Objekte verweisen; ausgeführt wird die Methode der \emph{tatsächlichen} Klasse.\\\hline
\end{tabular}

\gw{0.6 Datenbanken und SQL}
Datenbanken speichern strukturierte Daten in \textbf{Tabellen}: Ein \textbf{Datensatz} (Zeile, Tupel) enthält die Attributwerte (Spalten $\to$ \textbf{Attribute}) eines Objekts. Der \textbf{Primärschlüssel} (unterstrichen; ein oder mehrere Attribute) identifiziert jeden Datensatz eindeutig. Beziehungen zwischen Tabellen: \textbf{Fremdschlüssel}.\\[0.3mm]
\textbf{SQL} (Structured Query Language):\\
\begin{tabular}{@{}l@{\quad}l@{}}
\code{SELECT [DISTINCT] a1, a2} & welche Spalten (ohne Duplikate)\\
\code{FROM tabelle} & aus welcher Tabelle\\
\code{WHERE bedingung} & welche Zeilen\\
\code{ORDER BY a [ASC|DESC]} & auf- bzw. absteigend sortieren\\
\end{tabular}\\[0.3mm]
Beispiel: Schüler:innen der 9a mit Geburtstag im Oktober:
\begin{java}[style=klein,language=SQL,aboveskip=0.3mm]
SELECT Nachname, Vorname, Geburtstag FROM Schueler
WHERE Klasse = '9a' AND Geburtstag LIKE '%-10-%'
ORDER BY Geburtstag ASC;
\end{java}

\gw{0.7 Graphen}
Ein Graph besteht aus \textbf{Knoten} und \textbf{Kanten}. \textbf{Pfad}: Folge durch Kanten verbundener Knoten, keine Kante doppelt. \textbf{Zyklus}: Pfad mit gleichem Start- und Endknoten ($\to$ zyklischer/azyklischer Graph). \textbf{Gerichtet}: Pfeile (eine Richtung), \textbf{ungerichtet}: Striche (beide Richtungen). \textbf{Gewichtet}: Kanten tragen Werte (z.\,B. Entfernung). \textbf{Länge eines Pfades}: Anzahl der Kanten bzw. Summe der Gewichte. Algorithmen: Breiten-, Tiefensuche, Dijkstra.

\gw{0.8 Client-Server-Modell}
Als Graph gerichtet und sternförmig mit dem \textbf{Server} in der Mitte: Jeder \textbf{Client} stellt eine \textbf{Anfrage}, der Server schickt eine \textbf{Antwort} (Zyklus der Länge 2). Protokolle: HTTP(S), DNS, SMTP, IMAP, FTP; darunter TCP/IP. Adressierung über die \textbf{IP-Adresse}, auf Hardwareseite über die \textbf{MAC-Adresse}.

\gw{0.9 Codierung und Verschlüsselung}
\textbf{Codierung}: Information nach festen, öffentlichen Regeln anders darstellen (Binärzahl, ASCII, Morsecode, QR-Code). \textbf{Verschlüsselung}: Nur wer den \textbf{Schlüssel} kennt, kann den Klartext lesen. Das Verfahren darf bekannt sein; es muss so viele Schlüssel geben, dass Ausprobieren (Brute Force) aussichtslos ist. \textbf{Symmetrisch}: ein Schlüssel zum Ver- und Entschlüsseln (Caesar, Vigenère, AES); Problem: Schlüsselaustausch. \textbf{Asymmetrisch}: \textbf{öffentlicher} Schlüssel verschlüsselt, nur der \textbf{private} entschlüsselt (RSA).

\gw{0.10 Künstliche Intelligenz und Maschinelles Lernen}
Informatik $\supset$ Data Science $\supset$ KI $\supset$ Maschinelles Lernen $\supset$ Deep Learning.
\textbf{Schwache KI} löst eng begrenzte Aufgaben (Bilderkennung, Chatbot) ohne Bewusstsein, wie alle heutigen Systeme. \textbf{Starke KI} wäre allgemein intelligent wie ein Mensch (gibt es nicht).\\[0.3mm]
\begin{minipage}[c]{0.37\linewidth}
\begin{tikzpicture}[font=\scriptsize,x=1cm,y=1cm]
\node[draw,circle,inner sep=1pt] (n) at (1.3,0) {$\Sigma\,|\,f$};
\node (x1) at (0,0.4) {$x_1$}; \node (x2) at (0,-0.4) {$x_2$}; \node (b) at (1.3,-0.8) {$b$};
\draw[-{Stealth}] (x1) -- node[above,pos=0.4] {$w_1$} (n); \draw[-{Stealth}] (x2) -- node[below,pos=0.4] {$w_2$} (n);
\draw[-{Stealth}] (b) -- (n); \draw[-{Stealth}] (n) -- ++(0.8,0) node[right] {$y$};
\end{tikzpicture}
\end{minipage}\hfill
\begin{minipage}[c]{0.62\linewidth}
\textbf{Neuron}: Eingaben $x_i$ mal \textbf{Gewichte} $w_i$, aufsummiert, plus \textbf{Bias} $b$; die \textbf{Aktivierungsfunktion} $f$ (z.\,B. Schwellenwert) liefert die Ausgabe $y$.
\end{minipage}\\[0.3mm]
\textbf{Überwachtes} Lernen: aus Beispielen mit Lösung (Label), z.\,B. Spamfilter. \textbf{Unüberwachtes}: findet Muster/Gruppen in Daten ohne Label. \textbf{Bestärkendes}: Agent lernt durch Belohnung und Bestrafung.
\end{multicols}

\newpage
\teil{Teil B: Neu in Jahrgangsstufe 12 – kurze Übersichtsaufgaben}
\footnotesize\setlength{\parskip}{1.5pt}
\begin{aufgabe}[Rekursion]{1}
\begin{minipage}[t]{0.35\linewidth}
\vspace{0pt}
\begin{java}[style=klein,aboveskip=0mm,belowskip=0mm]
public int summe(int n) {
   if (n == 0) {
      return 0;             // §\ding{172}§
   }
   return n + summe(n - 1); // §\ding{173}§
}
\end{java}
\end{minipage}\hfill
\begin{minipage}[t]{0.63\linewidth}
\vspace{0pt}
\textbf{a)} Benenne die beiden Bestandteile jeder rekursiven Methode:\\[0.5mm]
\ding{172} \lsg[4.2cm]{Rekursionsanfang (Abbruch)}\enspace \ding{173} \lsg[5.4cm]{rekursiver Aufruf, kleineres Problem}\\[0.5mm]
\textbf{b)} Was passiert, wenn \ding{172} fehlt? \lsg[6.6cm]{endlos viele Aufrufe $\to$ StackOverflowError}\\[0.5mm]
\textbf{c)} Berechne: \code{summe(3)} = \lsg[7.7cm]{3 + summe(2) = … = 3 + 2 + 1 + summe(0) = 6}
\end{minipage}
\end{aufgabe}

\begin{aufgabe}[Rekursive Datenstrukturen]{2}
\textbf{a)} Ergänze die Übersicht.\\[0.5mm]
\renewcommand{\arraystretch}{1.15}
\begin{tabularx}{\linewidth}{|p{2.9cm}|X|X|X|}\hline
 & \textbf{Warteschlange (Queue)} & \textbf{Liste} & \textbf{Binärbaum}\\\hline
Nachfolger je Knoten\rule[-2mm]{0pt}{6.2mm} & genau einer (am Ende \code{null}) & \lsgoder{\lsgfarbe genau einer}{} & \lsgoder{\lsgfarbe bis zu zwei (links, rechts)}{}\\\hline
Einfügen/Entnehmen\rule[-2mm]{0pt}{6.2mm} & \code{add} hinten, \code{poll} vorne: \lsg[1cm]{FIFO} & \lsgoder{\lsgfarbe an beliebiger Stelle}{} & \lsgoder{\lsgfarbe links kleiner, rechts größer}{}\\\hline
\end{tabularx}\\[0.5mm]
\begin{minipage}[t]{0.75\linewidth}
\vspace{0pt}
\textbf{b)} Beim Entwurfsmuster \textbf{Kompositum} erben \code{Knoten} und \code{Abschluss} von der abstrakten Klasse \code{Listenelement}. Welche Aufgabe hat der \code{Abschluss}, welchen Vorteil bringt er?
\lsglinien{2}{Er ersetzt \code{null} am Ende: Jeder Knoten hat immer einen Nachfolger und ruft die Methode dort auf; der Abschluss beendet die Rekursion. Die Abfragen auf \code{null} entfallen.}
\textbf{c)} Füge \textbf{45} in den Suchbaum ein. \textbf{Inorder}-Reihenfolge (links -- Wurzel -- rechts):\\[0.5mm] \lsg[\linewidth]{20, 30, 40, 45, 50, 60, 70 (sortiert!)}\\[0.8mm]
\textbf{d)} Warum findet man Daten im Suchbaum schneller als in einer Liste?\\[0.5mm] \lsg[\linewidth]{Jeder Vergleich halbiert etwa den Suchbereich: ca. $\log_2 n$ statt $n$ Schritte.}
\end{minipage}\hfill
\begin{minipage}[t]{0.23\linewidth}
\vspace{1mm}\centering
\begin{tikzpicture}[level distance=9mm,level 1/.style={sibling distance=19mm},level 2/.style={sibling distance=9.5mm},every node/.style={draw,circle,inner sep=0.5pt,minimum size=5.5mm,font=\scriptsize}]
\node {50}
  child {node {30} child {node {20}} child {node (v) {40}}}
  child {node {70} child {node {60}} child[missing]};
\node[lsgbild,draw,circle,minimum size=5.5mm,font=\scriptsize,below right=4mm and 0mm of v] (n) {45};
\draw[lsgbild] (v) -- (n);
\end{tikzpicture}
\end{minipage}
\end{aufgabe}

\begin{aufgabe}[Nebenläufigkeit]{3}
\begin{minipage}[t]{0.42\linewidth}
\vspace{0pt}
Zwei Kassen (Threads) buchen \textbf{gleichzeitig} ab. Guthaben: 50\,€; Kasse A bucht 5\,€ ab, Kasse B 10\,€.
\begin{java}[style=klein,belowskip=0mm]
public void abbuchen(int betrag) {
   int neu = guthaben - betrag;
   guthaben = neu;
}
\end{java}
\end{minipage}\hfill
\begin{minipage}[t]{0.55\linewidth}
\vspace{0pt}
\textbf{a)} Erkläre, wie am Ende 40\,€ statt 35\,€ auf dem Konto stehen können.
\lsglinien{2}{A und B lesen beide 50\,€. A schreibt 45\,€, danach überschreibt B mit 40\,€: Die Abbuchung von A geht verloren (\emph{Lost Update}, Race Condition).}
\end{minipage}\\[1mm]
\textbf{b)} Der Codeabschnitt, in dem auf eine gemeinsam genutzte Ressource zugegriffen wird, heißt \lsg[3.3cm]{kritischer Abschnitt}. Darin darf sich immer nur \lsg[1.9cm]{ein Thread} befinden (\lsg[3.7cm]{gegenseitiger Ausschluss}); in Java mit dem Schlüsselwort \lsg[2.2cm]{synchronized}. Wartet jeder von zwei Prozessen auf ein Betriebsmittel, das der andere belegt, blockieren beide für immer: \lsg[3.7cm]{Verklemmung (Deadlock)}.
\end{aufgabe}

\begin{aufgabe}[Softwareengineering]{4}
\textbf{a) Wasserfallmodell}: Nummeriere die Phasen:\enspace
\lsg[0.45cm]{\,4} Test\enspace \lsg[0.45cm]{\,1} Analyse\enspace \lsg[0.45cm]{\,5} Wartung\enspace \lsg[0.45cm]{\,3} Implementierung\enspace \lsg[0.45cm]{\,2} Entwurf\\[0.8mm]
Nachteil: \lsg[15.8cm]{Jede Phase wird nur einmal durchlaufen: Fehler und Änderungswünsche fallen erst spät auf und sind dann teuer.}\\[0.8mm]
\textbf{b) Agile Methoden (Scrum)}: Ein \textbf{Sprint} ist ein \lsg[5.3cm]{kurzer, fester Zeitraum (1–4 Wochen)}. Ergebnis: \lsg[4cm]{lauffähiges Teilprodukt}\\[0.8mm]
Vorteil gegenüber dem Wasserfallmodell: \lsg[11.2cm]{frühe Rückmeldung der Kunden; Anforderungen dürfen sich ändern}\\[0.8mm]
\textbf{c)} Erkläre kurz:\\[0.3mm]
\renewcommand{\arraystretch}{1.12}
\begin{tabularx}{\linewidth}{|p{2.3cm}|X|}\hline
Modularisierung\rule[-3.2mm]{0pt}{7.5mm} & \lsgoder{\lsgfarbe Zerlegung in unabhängige Teile mit klarer Aufgabe: Arbeitsteilung, einzeln testbar, austauschbar}{}\\\hline
Schnittstelle\rule[-3.2mm]{0pt}{7.5mm} & \lsgoder{\lsgfarbe legt fest, welche Methoden (Signaturen) ein Modul anbietet, ohne die Umsetzung; in Java \code{interface}}{}\\\hline
MVC\rule[-3.2mm]{0pt}{7.5mm} & \lsgoder{\lsgfarbe \textbf{Model}: Daten und Logik, \textbf{View}: Anzeige, \textbf{Controller}: verarbeitet Eingaben, steuert Model und View}{}\\\hline
\end{tabularx}
\end{aufgabe}

\begin{aufgabe}[Informationssicherheit]{5}
Welches \textbf{Schutzziel} (Vertraulichkeit, Integrität, Verfügbarkeit, Authentizität) ist verletzt? Nenne eine Maßnahme.\\[0.5mm]
\renewcommand{\arraystretch}{1.12}
\begin{tabularx}{\linewidth}{|p{7.6cm}|p{2.4cm}|X|}\hline
\textbf{Vorfall} & \textbf{Schutzziel} & \textbf{Maßnahme}\\\hline
\rule[-2mm]{0pt}{6.2mm}Ein Bagger durchtrennt die Internetleitung der Schule. & \lsgoder{\lsgfarbe Verfügbarkeit}{} & \lsgoder{\lsgfarbe zweite Leitung, Backups}{}\\\hline
\rule[-2mm]{0pt}{6.2mm}Jemand ändert heimlich Noten in der Datenbank. & \lsgoder{\lsgfarbe Integrität}{} & \lsgoder{\lsgfarbe Zugriffsrechte, Protokollierung}{}\\\hline
\rule[-2mm]{0pt}{6.2mm}Eine Phishing-Mail gibt sich als Schulleitung aus. & \lsgoder{\lsgfarbe Authentizität}{} & \lsgoder{\lsgfarbe digitale Signatur, Absender prüfen}{}\\\hline
\rule[-2mm]{0pt}{6.2mm}Fremde lesen Schülerdaten auf einem offenen Laptop. & \lsgoder{\lsgfarbe Vertraulichkeit}{} & \lsgoder{\lsgfarbe Bildschirmsperre, Verschlüsselung}{}\\\hline
\end{tabularx}
\end{aufgabe}
\end{document}
"
%           (füllt alle Lücken rot aus; siehe \lsg, \lsglinien, \lsgoder)
% =====================================================================
\usepackage[a4paper,left=18mm,right=18mm,top=26mm,bottom=17mm,headheight=15mm,headsep=4mm,footskip=8mm]{geometry}
\usepackage{iftex}
\ifPDFTeX
  \usepackage[T1]{fontenc}
  \usepackage[utf8]{inputenc}
  \usepackage{helvet}
  \usepackage[scaled=0.86]{beramono}
  \renewcommand{\familydefault}{\sfdefault}
\else
  \usepackage{fontspec}
  \setmainfont{texgyreheros}[Extension=.otf,UprightFont=*-regular,BoldFont=*-bold,ItalicFont=*-italic,BoldItalicFont=*-bolditalic]
  \setmonofont{DejaVuSansMono}[Extension=.ttf,UprightFont=*,BoldFont=*-Bold,ItalicFont=*-Oblique,BoldItalicFont=*-BoldOblique,Scale=0.86]
\fi
\usepackage[ngerman,shorthands=off]{babel}
\usepackage{xcolor,graphicx,tabularx,array,enumitem,fancyhdr,tikz,listings,multicol,calc,etoolbox,pifont}
\usepackage[most]{tcolorbox}
\usetikzlibrary{arrows.meta,positioning,calc,decorations.pathreplacing,shapes.multipart}
\usepackage[hidelinks]{hyperref}
\setlength{\parindent}{0pt}
\setlength{\parskip}{3pt}
\setlist{nosep,leftmargin=*}

% ---------- Lösungsschalter ----------
\newif\ifloesung
\ifdefined\mitloesung\loesungtrue\fi
\newcommand{\lsgfarbe}{\color{red!75!black}}

% ---------- Kopf- und Fußzeile (einheitliche Form) ----------
\newcommand{\lehrkraft}{Hau}
\newcommand{\themenbereich}{Grundwissen 6 bis 12}
\newcommand{\abnummer}{}
\pagestyle{fancy}\fancyhf{}
\renewcommand{\headrulewidth}{0.4pt}
\fancyhead[L]{\small Klasse 13\\Datum:}
\fancyhead[C]{\small Themenbereich:\\\themenbereich}
\fancyhead[R]{\small \lehrkraft\ifloesung\\{\lsgfarbe\bfseries Lösung}\fi}
\fancyfoot[C]{\scriptsize edu-mrh.de\,·\,Informatik 13\,·\,Blatt \abnummer\ifloesung\ (Lösung)\fi\,·\,Seite \thepage}
\newcommand{\abtitel}[2]{\renewcommand{\abnummer}{#1}\begin{center}\Large\bfseries #1\quad\underline{#2}\end{center}\vspace{-1mm}}

% ---------- Boxen ----------
\newenvironment{aufgabe}[2][]{\begin{tcolorbox}[enhanced,breakable,colback=white,colframe=black,boxrule=0.7pt,arc=3mm,left=2mm,right=2mm,top=1.5mm,bottom=1.5mm,before skip=2.5mm,after skip=2.5mm]\textbf{Aufgabe #2)}\ifstrempty{#1}{}{ \textit{#1}}\par\vspace{0.6mm}}{\end{tcolorbox}}
\newtcolorbox{merke}[1][Merke]{enhanced,breakable,colback=green!5,colframe=green!40!black,boxrule=0.8pt,arc=2mm,left=2mm,right=2mm,top=1.5mm,bottom=1.5mm,before skip=2.5mm,after skip=2.5mm,title={#1},fonttitle=\bfseries,coltitle=white}
\newtcolorbox{hinweis}{enhanced,breakable,colback=yellow!12,colframe=orange!70!black,boxrule=0.6pt,arc=2mm,left=2mm,right=2mm,top=1mm,bottom=1mm,before skip=2mm,after skip=2mm}
\newcommand{\web}[1]{{\small\color{gray}\textit{Website: #1}}}

% ---------- Lücken, Linien, Karos ----------
\newcommand{\luecke}[1][3.5cm]{\rule[-0.5ex]{#1}{0.4pt}}
\newcommand{\linien}[1]{\par\vspace{1.5mm}\foreach \i in {1,...,#1}{\noindent\rule[0pt]{\linewidth}{0.3pt}\par\vspace{5.5mm}}\vspace{-3mm}}
\newcommand{\kariert}[1]{\par\vspace{1.5mm}\noindent\begin{tikzpicture}\draw[step=5mm,gray!55,thin] (0,0) grid (\linewidth-1pt,#1*5mm);\end{tikzpicture}\par}
\newcommand{\karobox}[2]{\begin{tikzpicture}\draw[step=5mm,gray!55,thin] (0,0) grid (#1,#2);\end{tikzpicture}}
\newcommand{\klassenkarte}[3]{\begin{tabular}{|l|}\hline\multicolumn{1}{|c|}{\textbf{#1}}\\\hline #2\\\hline #3\\\hline\end{tabular}}
\newcommand{\code}[1]{\texttt{#1}}

% Lücke mit Lösung: \lsg[Mindestbreite]{Lösungstext}
\newlength{\lsgbreite}
\newcommand{\lsg}[2][3.5cm]{\ifloesung\settowidth{\lsgbreite}{\,#2\,}\ifdim\lsgbreite<#1\setlength{\lsgbreite}{#1}\fi\underline{\makebox[\lsgbreite][l]{\lsgfarbe\,#2}}\else\luecke[#1]\fi}
% Schreiblinien mit Lösung: \lsglinien{Anzahl}{Lösungstext} (gleiche Höhe wie \linien)
\newcommand{\lsglinien}[2]{\ifloesung\par\vspace{1.5mm}\noindent\begin{minipage}[t][#1\dimexpr 6.9mm\relax][t]{\linewidth}\lsgfarbe #2\end{minipage}\par\vspace{-1.5mm}\else\linien{#1}\fi}
% Beliebiger Inhalt: \lsgoder{nur in der Lösung}{nur im Blatt}
\newcommand{\lsgoder}[2]{\ifloesung\leavevmode#1\else#2\fi}

% ---------- Verkettete Strukturen zeichnen ----------
% Objekt mit Kopf und Attributen:  \node[obj] (k1) {k1: Knoten \nodepart{second} daten = Hans\\ nachfolger};
\tikzset{
  obj/.style={draw,rectangle split,rectangle split parts=2,rectangle split part fill={orange!25,orange!8},
              rounded corners=2pt,font=\scriptsize,inner sep=2pt,align=left},
  qobj/.style={obj,rectangle split part fill={teal!25,teal!8}},
  pobj/.style={draw,rounded corners=2pt,fill=gray!8,font=\scriptsize,inner sep=2pt,align=left},
  zeiger/.style={thick,-{Stealth[length=2.2mm]}},
  lsgzeiger/.style={zeiger,red!75!black},
  nullref/.style={font=\scriptsize\ttfamily,inner sep=1pt},
  schritt/.style={circle,draw,fill=white,inner sep=0.6pt,font=\scriptsize\bfseries},
  lsgschritt/.style={schritt,draw=red!75!black,text=red!75!black},
}

% ---------- Java-Listings ----------
\lstset{language=Java,basicstyle=\ttfamily\small,keywordstyle=\bfseries,commentstyle=\color{gray!80!black},stringstyle=\color{black},showstringspaces=false,columns=flexible,keepspaces=true,tabsize=3,frame=single,framerule=0.4pt,rulecolor=\color{black},xleftmargin=2mm,framexleftmargin=1mm,aboveskip=2mm,belowskip=2mm,breaklines=true,escapechar=§,
 literate={ä}{{\"a}}1 {ö}{{\"o}}1 {ü}{{\"u}}1 {Ä}{{\"A}}1 {Ö}{{\"O}}1 {Ü}{{\"U}}1 {ß}{{\ss}}1 {–}{{--}}1 {„}{{\glqq}}1 {“}{{\grqq}}1 {→}{{$\to$}}1 {①}{{\ding{172}}}1 {②}{{\ding{173}}}1 {③}{{\ding{174}}}1 {④}{{\ding{175}}}1}
\lstnewenvironment{java}[1][]{\lstset{#1}}{}

%  Übersetzen mit XeLaTeX, LuaLaTeX, pdfLaTeX oder Tectonic.
%
%  Blatt:   xelatex ab-0-grundwissen-13.tex
%  Lösung:  xelatex -jobname=ab-0-grundwissen-13-lsg "\def\mitloesung{}\documentclass[11pt,a4paper]{article}
% =====================================================================
%  Gemeinsame Vorlage der Arbeitsblätter – Grundwissen Q13
%  – Informatik 12 (gleiches Layout wie informatik/10/java/arbeitsblaetter/
%  und informatik/11/graphen/arbeitsblaetter/ab-vorlage.tex)
%  Einbinden in jedem Blatt direkt nach \documentclass: % =====================================================================
%  Gemeinsame Vorlage der Arbeitsblätter – Grundwissen Q13
%  – Informatik 12 (gleiches Layout wie informatik/10/java/arbeitsblaetter/
%  und informatik/11/graphen/arbeitsblaetter/ab-vorlage.tex)
%  Einbinden in jedem Blatt direkt nach \documentclass: \input{ab-vorlage}
%  Übersetzen mit XeLaTeX, LuaLaTeX, pdfLaTeX oder Tectonic.
%
%  Blatt:   xelatex ab-0-grundwissen-13.tex
%  Lösung:  xelatex -jobname=ab-0-grundwissen-13-lsg "\def\mitloesung{}\input{ab-0-grundwissen-13}"
%           (füllt alle Lücken rot aus; siehe \lsg, \lsglinien, \lsgoder)
% =====================================================================
\usepackage[a4paper,left=18mm,right=18mm,top=26mm,bottom=17mm,headheight=15mm,headsep=4mm,footskip=8mm]{geometry}
\usepackage{iftex}
\ifPDFTeX
  \usepackage[T1]{fontenc}
  \usepackage[utf8]{inputenc}
  \usepackage{helvet}
  \usepackage[scaled=0.86]{beramono}
  \renewcommand{\familydefault}{\sfdefault}
\else
  \usepackage{fontspec}
  \setmainfont{texgyreheros}[Extension=.otf,UprightFont=*-regular,BoldFont=*-bold,ItalicFont=*-italic,BoldItalicFont=*-bolditalic]
  \setmonofont{DejaVuSansMono}[Extension=.ttf,UprightFont=*,BoldFont=*-Bold,ItalicFont=*-Oblique,BoldItalicFont=*-BoldOblique,Scale=0.86]
\fi
\usepackage[ngerman,shorthands=off]{babel}
\usepackage{xcolor,graphicx,tabularx,array,enumitem,fancyhdr,tikz,listings,multicol,calc,etoolbox,pifont}
\usepackage[most]{tcolorbox}
\usetikzlibrary{arrows.meta,positioning,calc,decorations.pathreplacing,shapes.multipart}
\usepackage[hidelinks]{hyperref}
\setlength{\parindent}{0pt}
\setlength{\parskip}{3pt}
\setlist{nosep,leftmargin=*}

% ---------- Lösungsschalter ----------
\newif\ifloesung
\ifdefined\mitloesung\loesungtrue\fi
\newcommand{\lsgfarbe}{\color{red!75!black}}

% ---------- Kopf- und Fußzeile (einheitliche Form) ----------
\newcommand{\lehrkraft}{Hau}
\newcommand{\themenbereich}{Grundwissen 6 bis 12}
\newcommand{\abnummer}{}
\pagestyle{fancy}\fancyhf{}
\renewcommand{\headrulewidth}{0.4pt}
\fancyhead[L]{\small Klasse 13\\Datum:}
\fancyhead[C]{\small Themenbereich:\\\themenbereich}
\fancyhead[R]{\small \lehrkraft\ifloesung\\{\lsgfarbe\bfseries Lösung}\fi}
\fancyfoot[C]{\scriptsize edu-mrh.de\,·\,Informatik 13\,·\,Blatt \abnummer\ifloesung\ (Lösung)\fi\,·\,Seite \thepage}
\newcommand{\abtitel}[2]{\renewcommand{\abnummer}{#1}\begin{center}\Large\bfseries #1\quad\underline{#2}\end{center}\vspace{-1mm}}

% ---------- Boxen ----------
\newenvironment{aufgabe}[2][]{\begin{tcolorbox}[enhanced,breakable,colback=white,colframe=black,boxrule=0.7pt,arc=3mm,left=2mm,right=2mm,top=1.5mm,bottom=1.5mm,before skip=2.5mm,after skip=2.5mm]\textbf{Aufgabe #2)}\ifstrempty{#1}{}{ \textit{#1}}\par\vspace{0.6mm}}{\end{tcolorbox}}
\newtcolorbox{merke}[1][Merke]{enhanced,breakable,colback=green!5,colframe=green!40!black,boxrule=0.8pt,arc=2mm,left=2mm,right=2mm,top=1.5mm,bottom=1.5mm,before skip=2.5mm,after skip=2.5mm,title={#1},fonttitle=\bfseries,coltitle=white}
\newtcolorbox{hinweis}{enhanced,breakable,colback=yellow!12,colframe=orange!70!black,boxrule=0.6pt,arc=2mm,left=2mm,right=2mm,top=1mm,bottom=1mm,before skip=2mm,after skip=2mm}
\newcommand{\web}[1]{{\small\color{gray}\textit{Website: #1}}}

% ---------- Lücken, Linien, Karos ----------
\newcommand{\luecke}[1][3.5cm]{\rule[-0.5ex]{#1}{0.4pt}}
\newcommand{\linien}[1]{\par\vspace{1.5mm}\foreach \i in {1,...,#1}{\noindent\rule[0pt]{\linewidth}{0.3pt}\par\vspace{5.5mm}}\vspace{-3mm}}
\newcommand{\kariert}[1]{\par\vspace{1.5mm}\noindent\begin{tikzpicture}\draw[step=5mm,gray!55,thin] (0,0) grid (\linewidth-1pt,#1*5mm);\end{tikzpicture}\par}
\newcommand{\karobox}[2]{\begin{tikzpicture}\draw[step=5mm,gray!55,thin] (0,0) grid (#1,#2);\end{tikzpicture}}
\newcommand{\klassenkarte}[3]{\begin{tabular}{|l|}\hline\multicolumn{1}{|c|}{\textbf{#1}}\\\hline #2\\\hline #3\\\hline\end{tabular}}
\newcommand{\code}[1]{\texttt{#1}}

% Lücke mit Lösung: \lsg[Mindestbreite]{Lösungstext}
\newlength{\lsgbreite}
\newcommand{\lsg}[2][3.5cm]{\ifloesung\settowidth{\lsgbreite}{\,#2\,}\ifdim\lsgbreite<#1\setlength{\lsgbreite}{#1}\fi\underline{\makebox[\lsgbreite][l]{\lsgfarbe\,#2}}\else\luecke[#1]\fi}
% Schreiblinien mit Lösung: \lsglinien{Anzahl}{Lösungstext} (gleiche Höhe wie \linien)
\newcommand{\lsglinien}[2]{\ifloesung\par\vspace{1.5mm}\noindent\begin{minipage}[t][#1\dimexpr 6.9mm\relax][t]{\linewidth}\lsgfarbe #2\end{minipage}\par\vspace{-1.5mm}\else\linien{#1}\fi}
% Beliebiger Inhalt: \lsgoder{nur in der Lösung}{nur im Blatt}
\newcommand{\lsgoder}[2]{\ifloesung\leavevmode#1\else#2\fi}

% ---------- Verkettete Strukturen zeichnen ----------
% Objekt mit Kopf und Attributen:  \node[obj] (k1) {k1: Knoten \nodepart{second} daten = Hans\\ nachfolger};
\tikzset{
  obj/.style={draw,rectangle split,rectangle split parts=2,rectangle split part fill={orange!25,orange!8},
              rounded corners=2pt,font=\scriptsize,inner sep=2pt,align=left},
  qobj/.style={obj,rectangle split part fill={teal!25,teal!8}},
  pobj/.style={draw,rounded corners=2pt,fill=gray!8,font=\scriptsize,inner sep=2pt,align=left},
  zeiger/.style={thick,-{Stealth[length=2.2mm]}},
  lsgzeiger/.style={zeiger,red!75!black},
  nullref/.style={font=\scriptsize\ttfamily,inner sep=1pt},
  schritt/.style={circle,draw,fill=white,inner sep=0.6pt,font=\scriptsize\bfseries},
  lsgschritt/.style={schritt,draw=red!75!black,text=red!75!black},
}

% ---------- Java-Listings ----------
\lstset{language=Java,basicstyle=\ttfamily\small,keywordstyle=\bfseries,commentstyle=\color{gray!80!black},stringstyle=\color{black},showstringspaces=false,columns=flexible,keepspaces=true,tabsize=3,frame=single,framerule=0.4pt,rulecolor=\color{black},xleftmargin=2mm,framexleftmargin=1mm,aboveskip=2mm,belowskip=2mm,breaklines=true,escapechar=§,
 literate={ä}{{\"a}}1 {ö}{{\"o}}1 {ü}{{\"u}}1 {Ä}{{\"A}}1 {Ö}{{\"O}}1 {Ü}{{\"U}}1 {ß}{{\ss}}1 {–}{{--}}1 {„}{{\glqq}}1 {“}{{\grqq}}1 {→}{{$\to$}}1 {①}{{\ding{172}}}1 {②}{{\ding{173}}}1 {③}{{\ding{174}}}1 {④}{{\ding{175}}}1}
\lstnewenvironment{java}[1][]{\lstset{#1}}{}

%  Übersetzen mit XeLaTeX, LuaLaTeX, pdfLaTeX oder Tectonic.
%
%  Blatt:   xelatex ab-0-grundwissen-13.tex
%  Lösung:  xelatex -jobname=ab-0-grundwissen-13-lsg "\def\mitloesung{}\documentclass[11pt,a4paper]{article}
\input{ab-vorlage}
% Blatt:   tectonic ab-0-grundwissen-13.tex
% Lösung:  tectonic ab-0-grundwissen-13-lsg.tex  (Datei enthält nur: \def\mitloesung{}\input{ab-0-grundwissen-13})
% Seite 1: Grundwissen 6 bis 11 als ausgefüllte Übersicht (identisch in Blatt und Lösung)
% Seite 2: Übersichtsaufgaben zu den Inhalten der Jahrgangsstufe 12
\usetikzlibrary{shapes.geometric}
\newgeometry{left=14mm,right=14mm,top=21mm,bottom=12mm,headheight=11mm,headsep=3mm,footskip=5mm}
\ifloesung\tikzset{lsgbild/.style={red!75!black}}\else\tikzset{lsgbild/.style={opacity=0}}\fi
\newcommand{\gw}[1]{\par\vspace{0.6mm}{\bfseries\color{green!40!black}#1}\par}
\newcommand{\teil}[1]{{\bfseries #1}\par\vspace{0.5mm}}
\setlength{\columnsep}{5mm}
\setlength{\columnseprule}{0.3pt}
\newcommand{\kopf}[1]{\multicolumn{1}{c}{\textbf{#1}}}
\renewenvironment{aufgabe}[2][]{\begin{tcolorbox}[enhanced,breakable,colback=white,colframe=black,boxrule=0.7pt,arc=3mm,left=2mm,right=2mm,top=1mm,bottom=1mm,before skip=1.8mm,after skip=1.8mm]\textbf{Aufgabe #2)}\ifstrempty{#1}{}{ \textit{#1}}\par\vspace{0.3mm}}{\end{tcolorbox}}
\lstdefinestyle{klein}{basicstyle=\ttfamily\scriptsize,aboveskip=0.8mm,belowskip=0.8mm}
\begin{document}
\abtitel{0}{Grundwissen}
\vspace{-3mm}
\teil{Teil A: Grundwissen der Jahrgangsstufen 6 bis 11 (zum Nachschlagen)}

\begin{minipage}[t]{0.575\linewidth}
\vspace{0pt}
\begin{tikzpicture}[x=1cm,y=1cm,font=\footnotesize]
\node[anchor=north west] (kk) at (0,0) {\klassenkarte{Dreieck}{seitenlaenge\\farbe}{setzeFarbe(neu)}};
\node[anchor=south west,font=\footnotesize\bfseries] at (0,0.05) {Klasse (Bauplan)};
\node[anchor=north west,draw,rounded corners=2.5mm,inner sep=0pt] (ok) at (5.35,0) {\begin{tabular}{l}\rule{0pt}{3.5mm}\textbf{form1: Dreieck}\\[0.5mm]\hline\rule{0pt}{3.5mm}seitenlaenge = 20\\farbe = "blau"\\[0.5mm]\end{tabular}};
\node[anchor=south west,font=\footnotesize\bfseries] at (5.35,0.05) {Objekt (Instanz)};
\draw[decorate,decoration={brace,amplitude=3pt}] ([xshift=2pt]kk.north east) -- ([xshift=2pt]kk.north east|-{(0,-0.5)}) node[midway,right=4pt,font=\scriptsize] {Klassenname};
\draw[decorate,decoration={brace,amplitude=3pt}] ([xshift=2pt]kk.north east|-{(0,-0.58)}) -- ([xshift=2pt]kk.north east|-{(0,-1.35)}) node[midway,right=4pt,font=\scriptsize] {Attribute};
\draw[decorate,decoration={brace,amplitude=3pt}] ([xshift=2pt]kk.north east|-{(0,-1.43)}) -- ([xshift=2pt]kk.south east) node[midway,right=4pt,font=\scriptsize] {Methoden};
\draw[decorate,decoration={brace,amplitude=3pt}] ([xshift=2pt]ok.north east) -- ([xshift=2pt]ok.north east|-{(0,-0.55)}) node[midway,right=4pt,align=left,font=\scriptsize] {Objektname:\\Klasse};
\draw[decorate,decoration={brace,amplitude=3pt}] ([xshift=2pt]ok.north east|-{(0,-0.63)}) -- ([xshift=2pt]ok.south east) node[midway,right=4pt,align=left,font=\scriptsize] {Attribut-\\werte};
\end{tikzpicture}
\end{minipage}\hfill
\begin{minipage}[t]{0.405\linewidth}
\vspace{0pt}\footnotesize
\textbf{Objektorientierung}: Problemlösen durch Analyse der beteiligten \textbf{Objekte} und ihres Zusammenspiels. Die \textbf{Klasse} ist der Bauplan, ein \textbf{Objekt} eine konkrete Ausprägung (Instanz) der Klasse.\\[0.5mm]
\textbf{Punktnotation}: \code{objekt.methode(parameter);}\\ \code{objekt.attribut = wert;}\quad z.\,B. \code{karol.schritt();}
\end{minipage}

\vspace{-1.5mm}
\begin{multicols}{2}\footnotesize\setlength{\parskip}{0.8pt}
\gw{0.1 Beziehungen, Datentypen, Zugriffsrechte}
Objekte können \textbf{Referenzen} auf andere Objekte enthalten (\textbf{Referenzattribute}); die referenzierte Klasse wird zum \textbf{Datentyp}, z.\,B. \code{private Person eigentuemer;} in \code{Haus}:\\[0.3mm]
\begin{tikzpicture}[font=\scriptsize]
\node[draw,inner sep=2pt] (h) {\textbf{Haus}};
\node[draw,inner sep=2pt,right=17mm of h] (p) {\textbf{Person}};
\draw[-{Stealth}] (h) -- node[above,inner sep=1pt] {eigentuemer} node[below,pos=0.85,inner sep=1pt] {1} (p);
\end{tikzpicture}\quad {\scriptsize(Klassendiagramm ohne \code{String}, \code{int})}\\[0.3mm]
Weitere Datentypen: \code{boolean}, \code{char}, \code{double}, \code{float}, \code{long}, \code{byte}, Felder, jede Klasse.
Objekte kommunizieren über \textbf{Methodenaufrufe}. \textbf{Datenkapselung}: Attribute sind \code{private}, Zugriff nur über \textbf{Getter} und \textbf{Setter}.\\[0.5mm]
\renewcommand{\arraystretch}{1}
\begin{tabular}{@{}|c|l|l|@{}}\hline
\code{+} & \code{public} & öffentlich: jede Klasse darf zugreifen\\\hline
\code{\#} & \code{protected} & eigene Klasse und Unterklassen\\\hline
\code{-} & \code{private} & nur innerhalb der eigenen Klasse\\\hline
\end{tabular}

\gw{0.2 Bausteine von Algorithmen: Struktogramm und Java}
\begin{tikzpicture}[x=1cm,y=1cm,font=\scriptsize]
\draw (0,0) rectangle (3.95,1.1); \draw (0,1.1) -- (1.97,0.55) -- (3.95,1.1); \draw (0,0.55) -- (3.95,0.55); \draw (1.97,0) -- (1.97,0.55);
\node at (1.97,0.93) {\ttfamily\itshape istWand()}; \node at (0.33,0.66) {wahr}; \node at (3.55,0.66) {falsch};
\node at (0.98,0.27) {\ttfamily linksDrehen()}; \node at (2.96,0.27) {\ttfamily schritt()};
\begin{scope}[xshift=4.3cm]
\draw (0,0) rectangle (4.05,1.1); \draw (0.4,0) -- (0.4,0.75) -- (4.05,0.75);
\node[anchor=west] at (0.05,0.93) {Wiederhole solange {\ttfamily\itshape !istWand()}};
\node[anchor=west] at (0.5,0.53) {\ttfamily hinlegen()}; \node[anchor=west] at (0.5,0.2) {\ttfamily schritt()};
\end{scope}
\end{tikzpicture}\\[-1.2mm]
\begin{minipage}[t]{0.47\linewidth}
\begin{java}[style=klein]
if (istWand()) {
   linksDrehen();
} else {
   schritt();
}
\end{java}
\end{minipage}\hfill
\begin{minipage}[t]{0.49\linewidth}
\begin{java}[style=klein]
while (!istWand()) {
   hinlegen();
   schritt();
}
\end{java}
\end{minipage}\\[0.3mm]
\textbf{Sequenz}: Anweisungen untereinander; \textbf{einseitig}: \code{if} ohne \code{else}; „Wiederhole 5-mal“: \code{for (int i = 0; i < 5; i++) \{…\}}

\gw{0.3 Feld (Array)}
Folge \underline{fester Länge} von Werten bzw. Objekten \underline{desselben Typs}. Zugriff über den \textbf{Index}; die Zählung beginnt bei 0.
\begin{java}[style=klein]
int zahl[];            // Deklaration
zahl = new int[1000];  // Instanziierung (Index 0..999)
zahl[2] = 10;          // Wertzuweisung (3. Element)
\end{java}

\gw{0.4 Das EVA-Prinzip}
\begin{tikzpicture}[font=\scriptsize,node distance=2.5mm,every node/.style={align=center}]
\node[draw,minimum height=7mm] (e) {\textbf{Eingabe}\\\code{gewicht, groesse}};
\node[draw,ellipse,inner sep=0pt,right=of e] (v) {\textbf{Verarbeitung}\\\code{bmi = gewicht/}\\\code{(groesse*groesse)}};
\node[draw,minimum height=7mm,right=of v] (a) {\textbf{Ausgabe}\\\code{println(…)}};
\draw[-{Stealth}] (e) -- (v); \draw[-{Stealth}] (v) -- (a);
\end{tikzpicture}\\
Bei Methoden: Eingabe über \textbf{Parameter}, Ausgabe als \textbf{Rückgabewert} oder auf dem Bildschirm.

\gw{0.5 Vererbung, @Override und Polymorphismus}
Eine \textbf{Unterklasse} erbt Attribute und Methoden ihrer \textbf{Oberklasse}; in Java von genau \underline{einer} Klasse (keine Mehrfachvererbung).\\[0.3mm]
\begin{minipage}[t]{0.21\linewidth}
\vspace{0pt}
\begin{tikzpicture}[font=\scriptsize]
\node[draw,rectangle split,rectangle split parts=2,align=left,inner sep=2pt] (p) {\textbf{Person}\nodepart{second}alter\\name};
\node[draw,rectangle split,rectangle split parts=2,align=left,inner sep=2pt,below=4mm of p] (s) {\textbf{Schueler}\nodepart{second}klasse};
\draw[-{Triangle[open,length=2mm]}] (s) -- (p);
\end{tikzpicture}
\end{minipage}\hfill
\begin{minipage}[t]{0.78\linewidth}
\vspace{0pt}
\begin{java}[style=klein,aboveskip=0mm,belowskip=0mm]
public class Schueler extends Person {
   private String klasse;
   Schueler(int a, String n, String k) {
      super(a, n);  klasse = k;
   }
}
\end{java}
\end{minipage}\\[0.8mm]
\begin{tabular}{@{}|l|p{6.25cm}|@{}}\hline
\code{extends} & gibt die Oberklasse an\\\hline
\code{super(…)} & ruft den Konstruktor der Oberklasse auf\\\hline
\code{@Override} & Unterklasse \textbf{überschreibt} eine geerbte Methode (gleiche Signatur, neues Verhalten)\\\hline
Polymorphie & Variable vom Typ der Oberklasse kann auf Unterklassen-Objekte verweisen; ausgeführt wird die Methode der \emph{tatsächlichen} Klasse.\\\hline
\end{tabular}

\gw{0.6 Datenbanken und SQL}
Datenbanken speichern strukturierte Daten in \textbf{Tabellen}: Ein \textbf{Datensatz} (Zeile, Tupel) enthält die Attributwerte (Spalten $\to$ \textbf{Attribute}) eines Objekts. Der \textbf{Primärschlüssel} (unterstrichen; ein oder mehrere Attribute) identifiziert jeden Datensatz eindeutig. Beziehungen zwischen Tabellen: \textbf{Fremdschlüssel}.\\[0.3mm]
\textbf{SQL} (Structured Query Language):\\
\begin{tabular}{@{}l@{\quad}l@{}}
\code{SELECT [DISTINCT] a1, a2} & welche Spalten (ohne Duplikate)\\
\code{FROM tabelle} & aus welcher Tabelle\\
\code{WHERE bedingung} & welche Zeilen\\
\code{ORDER BY a [ASC|DESC]} & auf- bzw. absteigend sortieren\\
\end{tabular}\\[0.3mm]
Beispiel: Schüler:innen der 9a mit Geburtstag im Oktober:
\begin{java}[style=klein,language=SQL,aboveskip=0.3mm]
SELECT Nachname, Vorname, Geburtstag FROM Schueler
WHERE Klasse = '9a' AND Geburtstag LIKE '%-10-%'
ORDER BY Geburtstag ASC;
\end{java}

\gw{0.7 Graphen}
Ein Graph besteht aus \textbf{Knoten} und \textbf{Kanten}. \textbf{Pfad}: Folge durch Kanten verbundener Knoten, keine Kante doppelt. \textbf{Zyklus}: Pfad mit gleichem Start- und Endknoten ($\to$ zyklischer/azyklischer Graph). \textbf{Gerichtet}: Pfeile (eine Richtung), \textbf{ungerichtet}: Striche (beide Richtungen). \textbf{Gewichtet}: Kanten tragen Werte (z.\,B. Entfernung). \textbf{Länge eines Pfades}: Anzahl der Kanten bzw. Summe der Gewichte. Algorithmen: Breiten-, Tiefensuche, Dijkstra.

\gw{0.8 Client-Server-Modell}
Als Graph gerichtet und sternförmig mit dem \textbf{Server} in der Mitte: Jeder \textbf{Client} stellt eine \textbf{Anfrage}, der Server schickt eine \textbf{Antwort} (Zyklus der Länge 2). Protokolle: HTTP(S), DNS, SMTP, IMAP, FTP; darunter TCP/IP. Adressierung über die \textbf{IP-Adresse}, auf Hardwareseite über die \textbf{MAC-Adresse}.

\gw{0.9 Codierung und Verschlüsselung}
\textbf{Codierung}: Information nach festen, öffentlichen Regeln anders darstellen (Binärzahl, ASCII, Morsecode, QR-Code). \textbf{Verschlüsselung}: Nur wer den \textbf{Schlüssel} kennt, kann den Klartext lesen. Das Verfahren darf bekannt sein; es muss so viele Schlüssel geben, dass Ausprobieren (Brute Force) aussichtslos ist. \textbf{Symmetrisch}: ein Schlüssel zum Ver- und Entschlüsseln (Caesar, Vigenère, AES); Problem: Schlüsselaustausch. \textbf{Asymmetrisch}: \textbf{öffentlicher} Schlüssel verschlüsselt, nur der \textbf{private} entschlüsselt (RSA).

\gw{0.10 Künstliche Intelligenz und Maschinelles Lernen}
Informatik $\supset$ Data Science $\supset$ KI $\supset$ Maschinelles Lernen $\supset$ Deep Learning.
\textbf{Schwache KI} löst eng begrenzte Aufgaben (Bilderkennung, Chatbot) ohne Bewusstsein, wie alle heutigen Systeme. \textbf{Starke KI} wäre allgemein intelligent wie ein Mensch (gibt es nicht).\\[0.3mm]
\begin{minipage}[c]{0.37\linewidth}
\begin{tikzpicture}[font=\scriptsize,x=1cm,y=1cm]
\node[draw,circle,inner sep=1pt] (n) at (1.3,0) {$\Sigma\,|\,f$};
\node (x1) at (0,0.4) {$x_1$}; \node (x2) at (0,-0.4) {$x_2$}; \node (b) at (1.3,-0.8) {$b$};
\draw[-{Stealth}] (x1) -- node[above,pos=0.4] {$w_1$} (n); \draw[-{Stealth}] (x2) -- node[below,pos=0.4] {$w_2$} (n);
\draw[-{Stealth}] (b) -- (n); \draw[-{Stealth}] (n) -- ++(0.8,0) node[right] {$y$};
\end{tikzpicture}
\end{minipage}\hfill
\begin{minipage}[c]{0.62\linewidth}
\textbf{Neuron}: Eingaben $x_i$ mal \textbf{Gewichte} $w_i$, aufsummiert, plus \textbf{Bias} $b$; die \textbf{Aktivierungsfunktion} $f$ (z.\,B. Schwellenwert) liefert die Ausgabe $y$.
\end{minipage}\\[0.3mm]
\textbf{Überwachtes} Lernen: aus Beispielen mit Lösung (Label), z.\,B. Spamfilter. \textbf{Unüberwachtes}: findet Muster/Gruppen in Daten ohne Label. \textbf{Bestärkendes}: Agent lernt durch Belohnung und Bestrafung.
\end{multicols}

\newpage
\teil{Teil B: Neu in Jahrgangsstufe 12 – kurze Übersichtsaufgaben}
\footnotesize\setlength{\parskip}{1.5pt}
\begin{aufgabe}[Rekursion]{1}
\begin{minipage}[t]{0.35\linewidth}
\vspace{0pt}
\begin{java}[style=klein,aboveskip=0mm,belowskip=0mm]
public int summe(int n) {
   if (n == 0) {
      return 0;             // §\ding{172}§
   }
   return n + summe(n - 1); // §\ding{173}§
}
\end{java}
\end{minipage}\hfill
\begin{minipage}[t]{0.63\linewidth}
\vspace{0pt}
\textbf{a)} Benenne die beiden Bestandteile jeder rekursiven Methode:\\[0.5mm]
\ding{172} \lsg[4.2cm]{Rekursionsanfang (Abbruch)}\enspace \ding{173} \lsg[5.4cm]{rekursiver Aufruf, kleineres Problem}\\[0.5mm]
\textbf{b)} Was passiert, wenn \ding{172} fehlt? \lsg[6.6cm]{endlos viele Aufrufe $\to$ StackOverflowError}\\[0.5mm]
\textbf{c)} Berechne: \code{summe(3)} = \lsg[7.7cm]{3 + summe(2) = … = 3 + 2 + 1 + summe(0) = 6}
\end{minipage}
\end{aufgabe}

\begin{aufgabe}[Rekursive Datenstrukturen]{2}
\textbf{a)} Ergänze die Übersicht.\\[0.5mm]
\renewcommand{\arraystretch}{1.15}
\begin{tabularx}{\linewidth}{|p{2.9cm}|X|X|X|}\hline
 & \textbf{Warteschlange (Queue)} & \textbf{Liste} & \textbf{Binärbaum}\\\hline
Nachfolger je Knoten\rule[-2mm]{0pt}{6.2mm} & genau einer (am Ende \code{null}) & \lsgoder{\lsgfarbe genau einer}{} & \lsgoder{\lsgfarbe bis zu zwei (links, rechts)}{}\\\hline
Einfügen/Entnehmen\rule[-2mm]{0pt}{6.2mm} & \code{add} hinten, \code{poll} vorne: \lsg[1cm]{FIFO} & \lsgoder{\lsgfarbe an beliebiger Stelle}{} & \lsgoder{\lsgfarbe links kleiner, rechts größer}{}\\\hline
\end{tabularx}\\[0.5mm]
\begin{minipage}[t]{0.75\linewidth}
\vspace{0pt}
\textbf{b)} Beim Entwurfsmuster \textbf{Kompositum} erben \code{Knoten} und \code{Abschluss} von der abstrakten Klasse \code{Listenelement}. Welche Aufgabe hat der \code{Abschluss}, welchen Vorteil bringt er?
\lsglinien{2}{Er ersetzt \code{null} am Ende: Jeder Knoten hat immer einen Nachfolger und ruft die Methode dort auf; der Abschluss beendet die Rekursion. Die Abfragen auf \code{null} entfallen.}
\textbf{c)} Füge \textbf{45} in den Suchbaum ein. \textbf{Inorder}-Reihenfolge (links -- Wurzel -- rechts):\\[0.5mm] \lsg[\linewidth]{20, 30, 40, 45, 50, 60, 70 (sortiert!)}\\[0.8mm]
\textbf{d)} Warum findet man Daten im Suchbaum schneller als in einer Liste?\\[0.5mm] \lsg[\linewidth]{Jeder Vergleich halbiert etwa den Suchbereich: ca. $\log_2 n$ statt $n$ Schritte.}
\end{minipage}\hfill
\begin{minipage}[t]{0.23\linewidth}
\vspace{1mm}\centering
\begin{tikzpicture}[level distance=9mm,level 1/.style={sibling distance=19mm},level 2/.style={sibling distance=9.5mm},every node/.style={draw,circle,inner sep=0.5pt,minimum size=5.5mm,font=\scriptsize}]
\node {50}
  child {node {30} child {node {20}} child {node (v) {40}}}
  child {node {70} child {node {60}} child[missing]};
\node[lsgbild,draw,circle,minimum size=5.5mm,font=\scriptsize,below right=4mm and 0mm of v] (n) {45};
\draw[lsgbild] (v) -- (n);
\end{tikzpicture}
\end{minipage}
\end{aufgabe}

\begin{aufgabe}[Nebenläufigkeit]{3}
\begin{minipage}[t]{0.42\linewidth}
\vspace{0pt}
Zwei Kassen (Threads) buchen \textbf{gleichzeitig} ab. Guthaben: 50\,€; Kasse A bucht 5\,€ ab, Kasse B 10\,€.
\begin{java}[style=klein,belowskip=0mm]
public void abbuchen(int betrag) {
   int neu = guthaben - betrag;
   guthaben = neu;
}
\end{java}
\end{minipage}\hfill
\begin{minipage}[t]{0.55\linewidth}
\vspace{0pt}
\textbf{a)} Erkläre, wie am Ende 40\,€ statt 35\,€ auf dem Konto stehen können.
\lsglinien{2}{A und B lesen beide 50\,€. A schreibt 45\,€, danach überschreibt B mit 40\,€: Die Abbuchung von A geht verloren (\emph{Lost Update}, Race Condition).}
\end{minipage}\\[1mm]
\textbf{b)} Der Codeabschnitt, in dem auf eine gemeinsam genutzte Ressource zugegriffen wird, heißt \lsg[3.3cm]{kritischer Abschnitt}. Darin darf sich immer nur \lsg[1.9cm]{ein Thread} befinden (\lsg[3.7cm]{gegenseitiger Ausschluss}); in Java mit dem Schlüsselwort \lsg[2.2cm]{synchronized}. Wartet jeder von zwei Prozessen auf ein Betriebsmittel, das der andere belegt, blockieren beide für immer: \lsg[3.7cm]{Verklemmung (Deadlock)}.
\end{aufgabe}

\begin{aufgabe}[Softwareengineering]{4}
\textbf{a) Wasserfallmodell}: Nummeriere die Phasen:\enspace
\lsg[0.45cm]{\,4} Test\enspace \lsg[0.45cm]{\,1} Analyse\enspace \lsg[0.45cm]{\,5} Wartung\enspace \lsg[0.45cm]{\,3} Implementierung\enspace \lsg[0.45cm]{\,2} Entwurf\\[0.8mm]
Nachteil: \lsg[15.8cm]{Jede Phase wird nur einmal durchlaufen: Fehler und Änderungswünsche fallen erst spät auf und sind dann teuer.}\\[0.8mm]
\textbf{b) Agile Methoden (Scrum)}: Ein \textbf{Sprint} ist ein \lsg[5.3cm]{kurzer, fester Zeitraum (1–4 Wochen)}. Ergebnis: \lsg[4cm]{lauffähiges Teilprodukt}\\[0.8mm]
Vorteil gegenüber dem Wasserfallmodell: \lsg[11.2cm]{frühe Rückmeldung der Kunden; Anforderungen dürfen sich ändern}\\[0.8mm]
\textbf{c)} Erkläre kurz:\\[0.3mm]
\renewcommand{\arraystretch}{1.12}
\begin{tabularx}{\linewidth}{|p{2.3cm}|X|}\hline
Modularisierung\rule[-3.2mm]{0pt}{7.5mm} & \lsgoder{\lsgfarbe Zerlegung in unabhängige Teile mit klarer Aufgabe: Arbeitsteilung, einzeln testbar, austauschbar}{}\\\hline
Schnittstelle\rule[-3.2mm]{0pt}{7.5mm} & \lsgoder{\lsgfarbe legt fest, welche Methoden (Signaturen) ein Modul anbietet, ohne die Umsetzung; in Java \code{interface}}{}\\\hline
MVC\rule[-3.2mm]{0pt}{7.5mm} & \lsgoder{\lsgfarbe \textbf{Model}: Daten und Logik, \textbf{View}: Anzeige, \textbf{Controller}: verarbeitet Eingaben, steuert Model und View}{}\\\hline
\end{tabularx}
\end{aufgabe}

\begin{aufgabe}[Informationssicherheit]{5}
Welches \textbf{Schutzziel} (Vertraulichkeit, Integrität, Verfügbarkeit, Authentizität) ist verletzt? Nenne eine Maßnahme.\\[0.5mm]
\renewcommand{\arraystretch}{1.12}
\begin{tabularx}{\linewidth}{|p{7.6cm}|p{2.4cm}|X|}\hline
\textbf{Vorfall} & \textbf{Schutzziel} & \textbf{Maßnahme}\\\hline
\rule[-2mm]{0pt}{6.2mm}Ein Bagger durchtrennt die Internetleitung der Schule. & \lsgoder{\lsgfarbe Verfügbarkeit}{} & \lsgoder{\lsgfarbe zweite Leitung, Backups}{}\\\hline
\rule[-2mm]{0pt}{6.2mm}Jemand ändert heimlich Noten in der Datenbank. & \lsgoder{\lsgfarbe Integrität}{} & \lsgoder{\lsgfarbe Zugriffsrechte, Protokollierung}{}\\\hline
\rule[-2mm]{0pt}{6.2mm}Eine Phishing-Mail gibt sich als Schulleitung aus. & \lsgoder{\lsgfarbe Authentizität}{} & \lsgoder{\lsgfarbe digitale Signatur, Absender prüfen}{}\\\hline
\rule[-2mm]{0pt}{6.2mm}Fremde lesen Schülerdaten auf einem offenen Laptop. & \lsgoder{\lsgfarbe Vertraulichkeit}{} & \lsgoder{\lsgfarbe Bildschirmsperre, Verschlüsselung}{}\\\hline
\end{tabularx}
\end{aufgabe}
\end{document}
"
%           (füllt alle Lücken rot aus; siehe \lsg, \lsglinien, \lsgoder)
% =====================================================================
\usepackage[a4paper,left=18mm,right=18mm,top=26mm,bottom=17mm,headheight=15mm,headsep=4mm,footskip=8mm]{geometry}
\usepackage{iftex}
\ifPDFTeX
  \usepackage[T1]{fontenc}
  \usepackage[utf8]{inputenc}
  \usepackage{helvet}
  \usepackage[scaled=0.86]{beramono}
  \renewcommand{\familydefault}{\sfdefault}
\else
  \usepackage{fontspec}
  \setmainfont{texgyreheros}[Extension=.otf,UprightFont=*-regular,BoldFont=*-bold,ItalicFont=*-italic,BoldItalicFont=*-bolditalic]
  \setmonofont{DejaVuSansMono}[Extension=.ttf,UprightFont=*,BoldFont=*-Bold,ItalicFont=*-Oblique,BoldItalicFont=*-BoldOblique,Scale=0.86]
\fi
\usepackage[ngerman,shorthands=off]{babel}
\usepackage{xcolor,graphicx,tabularx,array,enumitem,fancyhdr,tikz,listings,multicol,calc,etoolbox,pifont}
\usepackage[most]{tcolorbox}
\usetikzlibrary{arrows.meta,positioning,calc,decorations.pathreplacing,shapes.multipart}
\usepackage[hidelinks]{hyperref}
\setlength{\parindent}{0pt}
\setlength{\parskip}{3pt}
\setlist{nosep,leftmargin=*}

% ---------- Lösungsschalter ----------
\newif\ifloesung
\ifdefined\mitloesung\loesungtrue\fi
\newcommand{\lsgfarbe}{\color{red!75!black}}

% ---------- Kopf- und Fußzeile (einheitliche Form) ----------
\newcommand{\lehrkraft}{Hau}
\newcommand{\themenbereich}{Grundwissen 6 bis 12}
\newcommand{\abnummer}{}
\pagestyle{fancy}\fancyhf{}
\renewcommand{\headrulewidth}{0.4pt}
\fancyhead[L]{\small Klasse 13\\Datum:}
\fancyhead[C]{\small Themenbereich:\\\themenbereich}
\fancyhead[R]{\small \lehrkraft\ifloesung\\{\lsgfarbe\bfseries Lösung}\fi}
\fancyfoot[C]{\scriptsize edu-mrh.de\,·\,Informatik 13\,·\,Blatt \abnummer\ifloesung\ (Lösung)\fi\,·\,Seite \thepage}
\newcommand{\abtitel}[2]{\renewcommand{\abnummer}{#1}\begin{center}\Large\bfseries #1\quad\underline{#2}\end{center}\vspace{-1mm}}

% ---------- Boxen ----------
\newenvironment{aufgabe}[2][]{\begin{tcolorbox}[enhanced,breakable,colback=white,colframe=black,boxrule=0.7pt,arc=3mm,left=2mm,right=2mm,top=1.5mm,bottom=1.5mm,before skip=2.5mm,after skip=2.5mm]\textbf{Aufgabe #2)}\ifstrempty{#1}{}{ \textit{#1}}\par\vspace{0.6mm}}{\end{tcolorbox}}
\newtcolorbox{merke}[1][Merke]{enhanced,breakable,colback=green!5,colframe=green!40!black,boxrule=0.8pt,arc=2mm,left=2mm,right=2mm,top=1.5mm,bottom=1.5mm,before skip=2.5mm,after skip=2.5mm,title={#1},fonttitle=\bfseries,coltitle=white}
\newtcolorbox{hinweis}{enhanced,breakable,colback=yellow!12,colframe=orange!70!black,boxrule=0.6pt,arc=2mm,left=2mm,right=2mm,top=1mm,bottom=1mm,before skip=2mm,after skip=2mm}
\newcommand{\web}[1]{{\small\color{gray}\textit{Website: #1}}}

% ---------- Lücken, Linien, Karos ----------
\newcommand{\luecke}[1][3.5cm]{\rule[-0.5ex]{#1}{0.4pt}}
\newcommand{\linien}[1]{\par\vspace{1.5mm}\foreach \i in {1,...,#1}{\noindent\rule[0pt]{\linewidth}{0.3pt}\par\vspace{5.5mm}}\vspace{-3mm}}
\newcommand{\kariert}[1]{\par\vspace{1.5mm}\noindent\begin{tikzpicture}\draw[step=5mm,gray!55,thin] (0,0) grid (\linewidth-1pt,#1*5mm);\end{tikzpicture}\par}
\newcommand{\karobox}[2]{\begin{tikzpicture}\draw[step=5mm,gray!55,thin] (0,0) grid (#1,#2);\end{tikzpicture}}
\newcommand{\klassenkarte}[3]{\begin{tabular}{|l|}\hline\multicolumn{1}{|c|}{\textbf{#1}}\\\hline #2\\\hline #3\\\hline\end{tabular}}
\newcommand{\code}[1]{\texttt{#1}}

% Lücke mit Lösung: \lsg[Mindestbreite]{Lösungstext}
\newlength{\lsgbreite}
\newcommand{\lsg}[2][3.5cm]{\ifloesung\settowidth{\lsgbreite}{\,#2\,}\ifdim\lsgbreite<#1\setlength{\lsgbreite}{#1}\fi\underline{\makebox[\lsgbreite][l]{\lsgfarbe\,#2}}\else\luecke[#1]\fi}
% Schreiblinien mit Lösung: \lsglinien{Anzahl}{Lösungstext} (gleiche Höhe wie \linien)
\newcommand{\lsglinien}[2]{\ifloesung\par\vspace{1.5mm}\noindent\begin{minipage}[t][#1\dimexpr 6.9mm\relax][t]{\linewidth}\lsgfarbe #2\end{minipage}\par\vspace{-1.5mm}\else\linien{#1}\fi}
% Beliebiger Inhalt: \lsgoder{nur in der Lösung}{nur im Blatt}
\newcommand{\lsgoder}[2]{\ifloesung\leavevmode#1\else#2\fi}

% ---------- Verkettete Strukturen zeichnen ----------
% Objekt mit Kopf und Attributen:  \node[obj] (k1) {k1: Knoten \nodepart{second} daten = Hans\\ nachfolger};
\tikzset{
  obj/.style={draw,rectangle split,rectangle split parts=2,rectangle split part fill={orange!25,orange!8},
              rounded corners=2pt,font=\scriptsize,inner sep=2pt,align=left},
  qobj/.style={obj,rectangle split part fill={teal!25,teal!8}},
  pobj/.style={draw,rounded corners=2pt,fill=gray!8,font=\scriptsize,inner sep=2pt,align=left},
  zeiger/.style={thick,-{Stealth[length=2.2mm]}},
  lsgzeiger/.style={zeiger,red!75!black},
  nullref/.style={font=\scriptsize\ttfamily,inner sep=1pt},
  schritt/.style={circle,draw,fill=white,inner sep=0.6pt,font=\scriptsize\bfseries},
  lsgschritt/.style={schritt,draw=red!75!black,text=red!75!black},
}

% ---------- Java-Listings ----------
\lstset{language=Java,basicstyle=\ttfamily\small,keywordstyle=\bfseries,commentstyle=\color{gray!80!black},stringstyle=\color{black},showstringspaces=false,columns=flexible,keepspaces=true,tabsize=3,frame=single,framerule=0.4pt,rulecolor=\color{black},xleftmargin=2mm,framexleftmargin=1mm,aboveskip=2mm,belowskip=2mm,breaklines=true,escapechar=§,
 literate={ä}{{\"a}}1 {ö}{{\"o}}1 {ü}{{\"u}}1 {Ä}{{\"A}}1 {Ö}{{\"O}}1 {Ü}{{\"U}}1 {ß}{{\ss}}1 {–}{{--}}1 {„}{{\glqq}}1 {“}{{\grqq}}1 {→}{{$\to$}}1 {①}{{\ding{172}}}1 {②}{{\ding{173}}}1 {③}{{\ding{174}}}1 {④}{{\ding{175}}}1}
\lstnewenvironment{java}[1][]{\lstset{#1}}{}

% Blatt:   tectonic ab-0-grundwissen-13.tex
% Lösung:  tectonic ab-0-grundwissen-13-lsg.tex  (Datei enthält nur: \def\mitloesung{}\documentclass[11pt,a4paper]{article}
% =====================================================================
%  Gemeinsame Vorlage der Arbeitsblätter – Grundwissen Q13
%  – Informatik 12 (gleiches Layout wie informatik/10/java/arbeitsblaetter/
%  und informatik/11/graphen/arbeitsblaetter/ab-vorlage.tex)
%  Einbinden in jedem Blatt direkt nach \documentclass: \input{ab-vorlage}
%  Übersetzen mit XeLaTeX, LuaLaTeX, pdfLaTeX oder Tectonic.
%
%  Blatt:   xelatex ab-0-grundwissen-13.tex
%  Lösung:  xelatex -jobname=ab-0-grundwissen-13-lsg "\def\mitloesung{}\input{ab-0-grundwissen-13}"
%           (füllt alle Lücken rot aus; siehe \lsg, \lsglinien, \lsgoder)
% =====================================================================
\usepackage[a4paper,left=18mm,right=18mm,top=26mm,bottom=17mm,headheight=15mm,headsep=4mm,footskip=8mm]{geometry}
\usepackage{iftex}
\ifPDFTeX
  \usepackage[T1]{fontenc}
  \usepackage[utf8]{inputenc}
  \usepackage{helvet}
  \usepackage[scaled=0.86]{beramono}
  \renewcommand{\familydefault}{\sfdefault}
\else
  \usepackage{fontspec}
  \setmainfont{texgyreheros}[Extension=.otf,UprightFont=*-regular,BoldFont=*-bold,ItalicFont=*-italic,BoldItalicFont=*-bolditalic]
  \setmonofont{DejaVuSansMono}[Extension=.ttf,UprightFont=*,BoldFont=*-Bold,ItalicFont=*-Oblique,BoldItalicFont=*-BoldOblique,Scale=0.86]
\fi
\usepackage[ngerman,shorthands=off]{babel}
\usepackage{xcolor,graphicx,tabularx,array,enumitem,fancyhdr,tikz,listings,multicol,calc,etoolbox,pifont}
\usepackage[most]{tcolorbox}
\usetikzlibrary{arrows.meta,positioning,calc,decorations.pathreplacing,shapes.multipart}
\usepackage[hidelinks]{hyperref}
\setlength{\parindent}{0pt}
\setlength{\parskip}{3pt}
\setlist{nosep,leftmargin=*}

% ---------- Lösungsschalter ----------
\newif\ifloesung
\ifdefined\mitloesung\loesungtrue\fi
\newcommand{\lsgfarbe}{\color{red!75!black}}

% ---------- Kopf- und Fußzeile (einheitliche Form) ----------
\newcommand{\lehrkraft}{Hau}
\newcommand{\themenbereich}{Grundwissen 6 bis 12}
\newcommand{\abnummer}{}
\pagestyle{fancy}\fancyhf{}
\renewcommand{\headrulewidth}{0.4pt}
\fancyhead[L]{\small Klasse 13\\Datum:}
\fancyhead[C]{\small Themenbereich:\\\themenbereich}
\fancyhead[R]{\small \lehrkraft\ifloesung\\{\lsgfarbe\bfseries Lösung}\fi}
\fancyfoot[C]{\scriptsize edu-mrh.de\,·\,Informatik 13\,·\,Blatt \abnummer\ifloesung\ (Lösung)\fi\,·\,Seite \thepage}
\newcommand{\abtitel}[2]{\renewcommand{\abnummer}{#1}\begin{center}\Large\bfseries #1\quad\underline{#2}\end{center}\vspace{-1mm}}

% ---------- Boxen ----------
\newenvironment{aufgabe}[2][]{\begin{tcolorbox}[enhanced,breakable,colback=white,colframe=black,boxrule=0.7pt,arc=3mm,left=2mm,right=2mm,top=1.5mm,bottom=1.5mm,before skip=2.5mm,after skip=2.5mm]\textbf{Aufgabe #2)}\ifstrempty{#1}{}{ \textit{#1}}\par\vspace{0.6mm}}{\end{tcolorbox}}
\newtcolorbox{merke}[1][Merke]{enhanced,breakable,colback=green!5,colframe=green!40!black,boxrule=0.8pt,arc=2mm,left=2mm,right=2mm,top=1.5mm,bottom=1.5mm,before skip=2.5mm,after skip=2.5mm,title={#1},fonttitle=\bfseries,coltitle=white}
\newtcolorbox{hinweis}{enhanced,breakable,colback=yellow!12,colframe=orange!70!black,boxrule=0.6pt,arc=2mm,left=2mm,right=2mm,top=1mm,bottom=1mm,before skip=2mm,after skip=2mm}
\newcommand{\web}[1]{{\small\color{gray}\textit{Website: #1}}}

% ---------- Lücken, Linien, Karos ----------
\newcommand{\luecke}[1][3.5cm]{\rule[-0.5ex]{#1}{0.4pt}}
\newcommand{\linien}[1]{\par\vspace{1.5mm}\foreach \i in {1,...,#1}{\noindent\rule[0pt]{\linewidth}{0.3pt}\par\vspace{5.5mm}}\vspace{-3mm}}
\newcommand{\kariert}[1]{\par\vspace{1.5mm}\noindent\begin{tikzpicture}\draw[step=5mm,gray!55,thin] (0,0) grid (\linewidth-1pt,#1*5mm);\end{tikzpicture}\par}
\newcommand{\karobox}[2]{\begin{tikzpicture}\draw[step=5mm,gray!55,thin] (0,0) grid (#1,#2);\end{tikzpicture}}
\newcommand{\klassenkarte}[3]{\begin{tabular}{|l|}\hline\multicolumn{1}{|c|}{\textbf{#1}}\\\hline #2\\\hline #3\\\hline\end{tabular}}
\newcommand{\code}[1]{\texttt{#1}}

% Lücke mit Lösung: \lsg[Mindestbreite]{Lösungstext}
\newlength{\lsgbreite}
\newcommand{\lsg}[2][3.5cm]{\ifloesung\settowidth{\lsgbreite}{\,#2\,}\ifdim\lsgbreite<#1\setlength{\lsgbreite}{#1}\fi\underline{\makebox[\lsgbreite][l]{\lsgfarbe\,#2}}\else\luecke[#1]\fi}
% Schreiblinien mit Lösung: \lsglinien{Anzahl}{Lösungstext} (gleiche Höhe wie \linien)
\newcommand{\lsglinien}[2]{\ifloesung\par\vspace{1.5mm}\noindent\begin{minipage}[t][#1\dimexpr 6.9mm\relax][t]{\linewidth}\lsgfarbe #2\end{minipage}\par\vspace{-1.5mm}\else\linien{#1}\fi}
% Beliebiger Inhalt: \lsgoder{nur in der Lösung}{nur im Blatt}
\newcommand{\lsgoder}[2]{\ifloesung\leavevmode#1\else#2\fi}

% ---------- Verkettete Strukturen zeichnen ----------
% Objekt mit Kopf und Attributen:  \node[obj] (k1) {k1: Knoten \nodepart{second} daten = Hans\\ nachfolger};
\tikzset{
  obj/.style={draw,rectangle split,rectangle split parts=2,rectangle split part fill={orange!25,orange!8},
              rounded corners=2pt,font=\scriptsize,inner sep=2pt,align=left},
  qobj/.style={obj,rectangle split part fill={teal!25,teal!8}},
  pobj/.style={draw,rounded corners=2pt,fill=gray!8,font=\scriptsize,inner sep=2pt,align=left},
  zeiger/.style={thick,-{Stealth[length=2.2mm]}},
  lsgzeiger/.style={zeiger,red!75!black},
  nullref/.style={font=\scriptsize\ttfamily,inner sep=1pt},
  schritt/.style={circle,draw,fill=white,inner sep=0.6pt,font=\scriptsize\bfseries},
  lsgschritt/.style={schritt,draw=red!75!black,text=red!75!black},
}

% ---------- Java-Listings ----------
\lstset{language=Java,basicstyle=\ttfamily\small,keywordstyle=\bfseries,commentstyle=\color{gray!80!black},stringstyle=\color{black},showstringspaces=false,columns=flexible,keepspaces=true,tabsize=3,frame=single,framerule=0.4pt,rulecolor=\color{black},xleftmargin=2mm,framexleftmargin=1mm,aboveskip=2mm,belowskip=2mm,breaklines=true,escapechar=§,
 literate={ä}{{\"a}}1 {ö}{{\"o}}1 {ü}{{\"u}}1 {Ä}{{\"A}}1 {Ö}{{\"O}}1 {Ü}{{\"U}}1 {ß}{{\ss}}1 {–}{{--}}1 {„}{{\glqq}}1 {“}{{\grqq}}1 {→}{{$\to$}}1 {①}{{\ding{172}}}1 {②}{{\ding{173}}}1 {③}{{\ding{174}}}1 {④}{{\ding{175}}}1}
\lstnewenvironment{java}[1][]{\lstset{#1}}{}

% Blatt:   tectonic ab-0-grundwissen-13.tex
% Lösung:  tectonic ab-0-grundwissen-13-lsg.tex  (Datei enthält nur: \def\mitloesung{}\documentclass[11pt,a4paper]{article}
\input{ab-vorlage}
% Blatt:   tectonic ab-0-grundwissen-13.tex
% Lösung:  tectonic ab-0-grundwissen-13-lsg.tex  (Datei enthält nur: \def\mitloesung{}\input{ab-0-grundwissen-13})
% Seite 1: Grundwissen 6 bis 11 als ausgefüllte Übersicht (identisch in Blatt und Lösung)
% Seite 2: Übersichtsaufgaben zu den Inhalten der Jahrgangsstufe 12
\usetikzlibrary{shapes.geometric}
\newgeometry{left=14mm,right=14mm,top=21mm,bottom=12mm,headheight=11mm,headsep=3mm,footskip=5mm}
\ifloesung\tikzset{lsgbild/.style={red!75!black}}\else\tikzset{lsgbild/.style={opacity=0}}\fi
\newcommand{\gw}[1]{\par\vspace{0.6mm}{\bfseries\color{green!40!black}#1}\par}
\newcommand{\teil}[1]{{\bfseries #1}\par\vspace{0.5mm}}
\setlength{\columnsep}{5mm}
\setlength{\columnseprule}{0.3pt}
\newcommand{\kopf}[1]{\multicolumn{1}{c}{\textbf{#1}}}
\renewenvironment{aufgabe}[2][]{\begin{tcolorbox}[enhanced,breakable,colback=white,colframe=black,boxrule=0.7pt,arc=3mm,left=2mm,right=2mm,top=1mm,bottom=1mm,before skip=1.8mm,after skip=1.8mm]\textbf{Aufgabe #2)}\ifstrempty{#1}{}{ \textit{#1}}\par\vspace{0.3mm}}{\end{tcolorbox}}
\lstdefinestyle{klein}{basicstyle=\ttfamily\scriptsize,aboveskip=0.8mm,belowskip=0.8mm}
\begin{document}
\abtitel{0}{Grundwissen}
\vspace{-3mm}
\teil{Teil A: Grundwissen der Jahrgangsstufen 6 bis 11 (zum Nachschlagen)}

\begin{minipage}[t]{0.575\linewidth}
\vspace{0pt}
\begin{tikzpicture}[x=1cm,y=1cm,font=\footnotesize]
\node[anchor=north west] (kk) at (0,0) {\klassenkarte{Dreieck}{seitenlaenge\\farbe}{setzeFarbe(neu)}};
\node[anchor=south west,font=\footnotesize\bfseries] at (0,0.05) {Klasse (Bauplan)};
\node[anchor=north west,draw,rounded corners=2.5mm,inner sep=0pt] (ok) at (5.35,0) {\begin{tabular}{l}\rule{0pt}{3.5mm}\textbf{form1: Dreieck}\\[0.5mm]\hline\rule{0pt}{3.5mm}seitenlaenge = 20\\farbe = "blau"\\[0.5mm]\end{tabular}};
\node[anchor=south west,font=\footnotesize\bfseries] at (5.35,0.05) {Objekt (Instanz)};
\draw[decorate,decoration={brace,amplitude=3pt}] ([xshift=2pt]kk.north east) -- ([xshift=2pt]kk.north east|-{(0,-0.5)}) node[midway,right=4pt,font=\scriptsize] {Klassenname};
\draw[decorate,decoration={brace,amplitude=3pt}] ([xshift=2pt]kk.north east|-{(0,-0.58)}) -- ([xshift=2pt]kk.north east|-{(0,-1.35)}) node[midway,right=4pt,font=\scriptsize] {Attribute};
\draw[decorate,decoration={brace,amplitude=3pt}] ([xshift=2pt]kk.north east|-{(0,-1.43)}) -- ([xshift=2pt]kk.south east) node[midway,right=4pt,font=\scriptsize] {Methoden};
\draw[decorate,decoration={brace,amplitude=3pt}] ([xshift=2pt]ok.north east) -- ([xshift=2pt]ok.north east|-{(0,-0.55)}) node[midway,right=4pt,align=left,font=\scriptsize] {Objektname:\\Klasse};
\draw[decorate,decoration={brace,amplitude=3pt}] ([xshift=2pt]ok.north east|-{(0,-0.63)}) -- ([xshift=2pt]ok.south east) node[midway,right=4pt,align=left,font=\scriptsize] {Attribut-\\werte};
\end{tikzpicture}
\end{minipage}\hfill
\begin{minipage}[t]{0.405\linewidth}
\vspace{0pt}\footnotesize
\textbf{Objektorientierung}: Problemlösen durch Analyse der beteiligten \textbf{Objekte} und ihres Zusammenspiels. Die \textbf{Klasse} ist der Bauplan, ein \textbf{Objekt} eine konkrete Ausprägung (Instanz) der Klasse.\\[0.5mm]
\textbf{Punktnotation}: \code{objekt.methode(parameter);}\\ \code{objekt.attribut = wert;}\quad z.\,B. \code{karol.schritt();}
\end{minipage}

\vspace{-1.5mm}
\begin{multicols}{2}\footnotesize\setlength{\parskip}{0.8pt}
\gw{0.1 Beziehungen, Datentypen, Zugriffsrechte}
Objekte können \textbf{Referenzen} auf andere Objekte enthalten (\textbf{Referenzattribute}); die referenzierte Klasse wird zum \textbf{Datentyp}, z.\,B. \code{private Person eigentuemer;} in \code{Haus}:\\[0.3mm]
\begin{tikzpicture}[font=\scriptsize]
\node[draw,inner sep=2pt] (h) {\textbf{Haus}};
\node[draw,inner sep=2pt,right=17mm of h] (p) {\textbf{Person}};
\draw[-{Stealth}] (h) -- node[above,inner sep=1pt] {eigentuemer} node[below,pos=0.85,inner sep=1pt] {1} (p);
\end{tikzpicture}\quad {\scriptsize(Klassendiagramm ohne \code{String}, \code{int})}\\[0.3mm]
Weitere Datentypen: \code{boolean}, \code{char}, \code{double}, \code{float}, \code{long}, \code{byte}, Felder, jede Klasse.
Objekte kommunizieren über \textbf{Methodenaufrufe}. \textbf{Datenkapselung}: Attribute sind \code{private}, Zugriff nur über \textbf{Getter} und \textbf{Setter}.\\[0.5mm]
\renewcommand{\arraystretch}{1}
\begin{tabular}{@{}|c|l|l|@{}}\hline
\code{+} & \code{public} & öffentlich: jede Klasse darf zugreifen\\\hline
\code{\#} & \code{protected} & eigene Klasse und Unterklassen\\\hline
\code{-} & \code{private} & nur innerhalb der eigenen Klasse\\\hline
\end{tabular}

\gw{0.2 Bausteine von Algorithmen: Struktogramm und Java}
\begin{tikzpicture}[x=1cm,y=1cm,font=\scriptsize]
\draw (0,0) rectangle (3.95,1.1); \draw (0,1.1) -- (1.97,0.55) -- (3.95,1.1); \draw (0,0.55) -- (3.95,0.55); \draw (1.97,0) -- (1.97,0.55);
\node at (1.97,0.93) {\ttfamily\itshape istWand()}; \node at (0.33,0.66) {wahr}; \node at (3.55,0.66) {falsch};
\node at (0.98,0.27) {\ttfamily linksDrehen()}; \node at (2.96,0.27) {\ttfamily schritt()};
\begin{scope}[xshift=4.3cm]
\draw (0,0) rectangle (4.05,1.1); \draw (0.4,0) -- (0.4,0.75) -- (4.05,0.75);
\node[anchor=west] at (0.05,0.93) {Wiederhole solange {\ttfamily\itshape !istWand()}};
\node[anchor=west] at (0.5,0.53) {\ttfamily hinlegen()}; \node[anchor=west] at (0.5,0.2) {\ttfamily schritt()};
\end{scope}
\end{tikzpicture}\\[-1.2mm]
\begin{minipage}[t]{0.47\linewidth}
\begin{java}[style=klein]
if (istWand()) {
   linksDrehen();
} else {
   schritt();
}
\end{java}
\end{minipage}\hfill
\begin{minipage}[t]{0.49\linewidth}
\begin{java}[style=klein]
while (!istWand()) {
   hinlegen();
   schritt();
}
\end{java}
\end{minipage}\\[0.3mm]
\textbf{Sequenz}: Anweisungen untereinander; \textbf{einseitig}: \code{if} ohne \code{else}; „Wiederhole 5-mal“: \code{for (int i = 0; i < 5; i++) \{…\}}

\gw{0.3 Feld (Array)}
Folge \underline{fester Länge} von Werten bzw. Objekten \underline{desselben Typs}. Zugriff über den \textbf{Index}; die Zählung beginnt bei 0.
\begin{java}[style=klein]
int zahl[];            // Deklaration
zahl = new int[1000];  // Instanziierung (Index 0..999)
zahl[2] = 10;          // Wertzuweisung (3. Element)
\end{java}

\gw{0.4 Das EVA-Prinzip}
\begin{tikzpicture}[font=\scriptsize,node distance=2.5mm,every node/.style={align=center}]
\node[draw,minimum height=7mm] (e) {\textbf{Eingabe}\\\code{gewicht, groesse}};
\node[draw,ellipse,inner sep=0pt,right=of e] (v) {\textbf{Verarbeitung}\\\code{bmi = gewicht/}\\\code{(groesse*groesse)}};
\node[draw,minimum height=7mm,right=of v] (a) {\textbf{Ausgabe}\\\code{println(…)}};
\draw[-{Stealth}] (e) -- (v); \draw[-{Stealth}] (v) -- (a);
\end{tikzpicture}\\
Bei Methoden: Eingabe über \textbf{Parameter}, Ausgabe als \textbf{Rückgabewert} oder auf dem Bildschirm.

\gw{0.5 Vererbung, @Override und Polymorphismus}
Eine \textbf{Unterklasse} erbt Attribute und Methoden ihrer \textbf{Oberklasse}; in Java von genau \underline{einer} Klasse (keine Mehrfachvererbung).\\[0.3mm]
\begin{minipage}[t]{0.21\linewidth}
\vspace{0pt}
\begin{tikzpicture}[font=\scriptsize]
\node[draw,rectangle split,rectangle split parts=2,align=left,inner sep=2pt] (p) {\textbf{Person}\nodepart{second}alter\\name};
\node[draw,rectangle split,rectangle split parts=2,align=left,inner sep=2pt,below=4mm of p] (s) {\textbf{Schueler}\nodepart{second}klasse};
\draw[-{Triangle[open,length=2mm]}] (s) -- (p);
\end{tikzpicture}
\end{minipage}\hfill
\begin{minipage}[t]{0.78\linewidth}
\vspace{0pt}
\begin{java}[style=klein,aboveskip=0mm,belowskip=0mm]
public class Schueler extends Person {
   private String klasse;
   Schueler(int a, String n, String k) {
      super(a, n);  klasse = k;
   }
}
\end{java}
\end{minipage}\\[0.8mm]
\begin{tabular}{@{}|l|p{6.25cm}|@{}}\hline
\code{extends} & gibt die Oberklasse an\\\hline
\code{super(…)} & ruft den Konstruktor der Oberklasse auf\\\hline
\code{@Override} & Unterklasse \textbf{überschreibt} eine geerbte Methode (gleiche Signatur, neues Verhalten)\\\hline
Polymorphie & Variable vom Typ der Oberklasse kann auf Unterklassen-Objekte verweisen; ausgeführt wird die Methode der \emph{tatsächlichen} Klasse.\\\hline
\end{tabular}

\gw{0.6 Datenbanken und SQL}
Datenbanken speichern strukturierte Daten in \textbf{Tabellen}: Ein \textbf{Datensatz} (Zeile, Tupel) enthält die Attributwerte (Spalten $\to$ \textbf{Attribute}) eines Objekts. Der \textbf{Primärschlüssel} (unterstrichen; ein oder mehrere Attribute) identifiziert jeden Datensatz eindeutig. Beziehungen zwischen Tabellen: \textbf{Fremdschlüssel}.\\[0.3mm]
\textbf{SQL} (Structured Query Language):\\
\begin{tabular}{@{}l@{\quad}l@{}}
\code{SELECT [DISTINCT] a1, a2} & welche Spalten (ohne Duplikate)\\
\code{FROM tabelle} & aus welcher Tabelle\\
\code{WHERE bedingung} & welche Zeilen\\
\code{ORDER BY a [ASC|DESC]} & auf- bzw. absteigend sortieren\\
\end{tabular}\\[0.3mm]
Beispiel: Schüler:innen der 9a mit Geburtstag im Oktober:
\begin{java}[style=klein,language=SQL,aboveskip=0.3mm]
SELECT Nachname, Vorname, Geburtstag FROM Schueler
WHERE Klasse = '9a' AND Geburtstag LIKE '%-10-%'
ORDER BY Geburtstag ASC;
\end{java}

\gw{0.7 Graphen}
Ein Graph besteht aus \textbf{Knoten} und \textbf{Kanten}. \textbf{Pfad}: Folge durch Kanten verbundener Knoten, keine Kante doppelt. \textbf{Zyklus}: Pfad mit gleichem Start- und Endknoten ($\to$ zyklischer/azyklischer Graph). \textbf{Gerichtet}: Pfeile (eine Richtung), \textbf{ungerichtet}: Striche (beide Richtungen). \textbf{Gewichtet}: Kanten tragen Werte (z.\,B. Entfernung). \textbf{Länge eines Pfades}: Anzahl der Kanten bzw. Summe der Gewichte. Algorithmen: Breiten-, Tiefensuche, Dijkstra.

\gw{0.8 Client-Server-Modell}
Als Graph gerichtet und sternförmig mit dem \textbf{Server} in der Mitte: Jeder \textbf{Client} stellt eine \textbf{Anfrage}, der Server schickt eine \textbf{Antwort} (Zyklus der Länge 2). Protokolle: HTTP(S), DNS, SMTP, IMAP, FTP; darunter TCP/IP. Adressierung über die \textbf{IP-Adresse}, auf Hardwareseite über die \textbf{MAC-Adresse}.

\gw{0.9 Codierung und Verschlüsselung}
\textbf{Codierung}: Information nach festen, öffentlichen Regeln anders darstellen (Binärzahl, ASCII, Morsecode, QR-Code). \textbf{Verschlüsselung}: Nur wer den \textbf{Schlüssel} kennt, kann den Klartext lesen. Das Verfahren darf bekannt sein; es muss so viele Schlüssel geben, dass Ausprobieren (Brute Force) aussichtslos ist. \textbf{Symmetrisch}: ein Schlüssel zum Ver- und Entschlüsseln (Caesar, Vigenère, AES); Problem: Schlüsselaustausch. \textbf{Asymmetrisch}: \textbf{öffentlicher} Schlüssel verschlüsselt, nur der \textbf{private} entschlüsselt (RSA).

\gw{0.10 Künstliche Intelligenz und Maschinelles Lernen}
Informatik $\supset$ Data Science $\supset$ KI $\supset$ Maschinelles Lernen $\supset$ Deep Learning.
\textbf{Schwache KI} löst eng begrenzte Aufgaben (Bilderkennung, Chatbot) ohne Bewusstsein, wie alle heutigen Systeme. \textbf{Starke KI} wäre allgemein intelligent wie ein Mensch (gibt es nicht).\\[0.3mm]
\begin{minipage}[c]{0.37\linewidth}
\begin{tikzpicture}[font=\scriptsize,x=1cm,y=1cm]
\node[draw,circle,inner sep=1pt] (n) at (1.3,0) {$\Sigma\,|\,f$};
\node (x1) at (0,0.4) {$x_1$}; \node (x2) at (0,-0.4) {$x_2$}; \node (b) at (1.3,-0.8) {$b$};
\draw[-{Stealth}] (x1) -- node[above,pos=0.4] {$w_1$} (n); \draw[-{Stealth}] (x2) -- node[below,pos=0.4] {$w_2$} (n);
\draw[-{Stealth}] (b) -- (n); \draw[-{Stealth}] (n) -- ++(0.8,0) node[right] {$y$};
\end{tikzpicture}
\end{minipage}\hfill
\begin{minipage}[c]{0.62\linewidth}
\textbf{Neuron}: Eingaben $x_i$ mal \textbf{Gewichte} $w_i$, aufsummiert, plus \textbf{Bias} $b$; die \textbf{Aktivierungsfunktion} $f$ (z.\,B. Schwellenwert) liefert die Ausgabe $y$.
\end{minipage}\\[0.3mm]
\textbf{Überwachtes} Lernen: aus Beispielen mit Lösung (Label), z.\,B. Spamfilter. \textbf{Unüberwachtes}: findet Muster/Gruppen in Daten ohne Label. \textbf{Bestärkendes}: Agent lernt durch Belohnung und Bestrafung.
\end{multicols}

\newpage
\teil{Teil B: Neu in Jahrgangsstufe 12 – kurze Übersichtsaufgaben}
\footnotesize\setlength{\parskip}{1.5pt}
\begin{aufgabe}[Rekursion]{1}
\begin{minipage}[t]{0.35\linewidth}
\vspace{0pt}
\begin{java}[style=klein,aboveskip=0mm,belowskip=0mm]
public int summe(int n) {
   if (n == 0) {
      return 0;             // §\ding{172}§
   }
   return n + summe(n - 1); // §\ding{173}§
}
\end{java}
\end{minipage}\hfill
\begin{minipage}[t]{0.63\linewidth}
\vspace{0pt}
\textbf{a)} Benenne die beiden Bestandteile jeder rekursiven Methode:\\[0.5mm]
\ding{172} \lsg[4.2cm]{Rekursionsanfang (Abbruch)}\enspace \ding{173} \lsg[5.4cm]{rekursiver Aufruf, kleineres Problem}\\[0.5mm]
\textbf{b)} Was passiert, wenn \ding{172} fehlt? \lsg[6.6cm]{endlos viele Aufrufe $\to$ StackOverflowError}\\[0.5mm]
\textbf{c)} Berechne: \code{summe(3)} = \lsg[7.7cm]{3 + summe(2) = … = 3 + 2 + 1 + summe(0) = 6}
\end{minipage}
\end{aufgabe}

\begin{aufgabe}[Rekursive Datenstrukturen]{2}
\textbf{a)} Ergänze die Übersicht.\\[0.5mm]
\renewcommand{\arraystretch}{1.15}
\begin{tabularx}{\linewidth}{|p{2.9cm}|X|X|X|}\hline
 & \textbf{Warteschlange (Queue)} & \textbf{Liste} & \textbf{Binärbaum}\\\hline
Nachfolger je Knoten\rule[-2mm]{0pt}{6.2mm} & genau einer (am Ende \code{null}) & \lsgoder{\lsgfarbe genau einer}{} & \lsgoder{\lsgfarbe bis zu zwei (links, rechts)}{}\\\hline
Einfügen/Entnehmen\rule[-2mm]{0pt}{6.2mm} & \code{add} hinten, \code{poll} vorne: \lsg[1cm]{FIFO} & \lsgoder{\lsgfarbe an beliebiger Stelle}{} & \lsgoder{\lsgfarbe links kleiner, rechts größer}{}\\\hline
\end{tabularx}\\[0.5mm]
\begin{minipage}[t]{0.75\linewidth}
\vspace{0pt}
\textbf{b)} Beim Entwurfsmuster \textbf{Kompositum} erben \code{Knoten} und \code{Abschluss} von der abstrakten Klasse \code{Listenelement}. Welche Aufgabe hat der \code{Abschluss}, welchen Vorteil bringt er?
\lsglinien{2}{Er ersetzt \code{null} am Ende: Jeder Knoten hat immer einen Nachfolger und ruft die Methode dort auf; der Abschluss beendet die Rekursion. Die Abfragen auf \code{null} entfallen.}
\textbf{c)} Füge \textbf{45} in den Suchbaum ein. \textbf{Inorder}-Reihenfolge (links -- Wurzel -- rechts):\\[0.5mm] \lsg[\linewidth]{20, 30, 40, 45, 50, 60, 70 (sortiert!)}\\[0.8mm]
\textbf{d)} Warum findet man Daten im Suchbaum schneller als in einer Liste?\\[0.5mm] \lsg[\linewidth]{Jeder Vergleich halbiert etwa den Suchbereich: ca. $\log_2 n$ statt $n$ Schritte.}
\end{minipage}\hfill
\begin{minipage}[t]{0.23\linewidth}
\vspace{1mm}\centering
\begin{tikzpicture}[level distance=9mm,level 1/.style={sibling distance=19mm},level 2/.style={sibling distance=9.5mm},every node/.style={draw,circle,inner sep=0.5pt,minimum size=5.5mm,font=\scriptsize}]
\node {50}
  child {node {30} child {node {20}} child {node (v) {40}}}
  child {node {70} child {node {60}} child[missing]};
\node[lsgbild,draw,circle,minimum size=5.5mm,font=\scriptsize,below right=4mm and 0mm of v] (n) {45};
\draw[lsgbild] (v) -- (n);
\end{tikzpicture}
\end{minipage}
\end{aufgabe}

\begin{aufgabe}[Nebenläufigkeit]{3}
\begin{minipage}[t]{0.42\linewidth}
\vspace{0pt}
Zwei Kassen (Threads) buchen \textbf{gleichzeitig} ab. Guthaben: 50\,€; Kasse A bucht 5\,€ ab, Kasse B 10\,€.
\begin{java}[style=klein,belowskip=0mm]
public void abbuchen(int betrag) {
   int neu = guthaben - betrag;
   guthaben = neu;
}
\end{java}
\end{minipage}\hfill
\begin{minipage}[t]{0.55\linewidth}
\vspace{0pt}
\textbf{a)} Erkläre, wie am Ende 40\,€ statt 35\,€ auf dem Konto stehen können.
\lsglinien{2}{A und B lesen beide 50\,€. A schreibt 45\,€, danach überschreibt B mit 40\,€: Die Abbuchung von A geht verloren (\emph{Lost Update}, Race Condition).}
\end{minipage}\\[1mm]
\textbf{b)} Der Codeabschnitt, in dem auf eine gemeinsam genutzte Ressource zugegriffen wird, heißt \lsg[3.3cm]{kritischer Abschnitt}. Darin darf sich immer nur \lsg[1.9cm]{ein Thread} befinden (\lsg[3.7cm]{gegenseitiger Ausschluss}); in Java mit dem Schlüsselwort \lsg[2.2cm]{synchronized}. Wartet jeder von zwei Prozessen auf ein Betriebsmittel, das der andere belegt, blockieren beide für immer: \lsg[3.7cm]{Verklemmung (Deadlock)}.
\end{aufgabe}

\begin{aufgabe}[Softwareengineering]{4}
\textbf{a) Wasserfallmodell}: Nummeriere die Phasen:\enspace
\lsg[0.45cm]{\,4} Test\enspace \lsg[0.45cm]{\,1} Analyse\enspace \lsg[0.45cm]{\,5} Wartung\enspace \lsg[0.45cm]{\,3} Implementierung\enspace \lsg[0.45cm]{\,2} Entwurf\\[0.8mm]
Nachteil: \lsg[15.8cm]{Jede Phase wird nur einmal durchlaufen: Fehler und Änderungswünsche fallen erst spät auf und sind dann teuer.}\\[0.8mm]
\textbf{b) Agile Methoden (Scrum)}: Ein \textbf{Sprint} ist ein \lsg[5.3cm]{kurzer, fester Zeitraum (1–4 Wochen)}. Ergebnis: \lsg[4cm]{lauffähiges Teilprodukt}\\[0.8mm]
Vorteil gegenüber dem Wasserfallmodell: \lsg[11.2cm]{frühe Rückmeldung der Kunden; Anforderungen dürfen sich ändern}\\[0.8mm]
\textbf{c)} Erkläre kurz:\\[0.3mm]
\renewcommand{\arraystretch}{1.12}
\begin{tabularx}{\linewidth}{|p{2.3cm}|X|}\hline
Modularisierung\rule[-3.2mm]{0pt}{7.5mm} & \lsgoder{\lsgfarbe Zerlegung in unabhängige Teile mit klarer Aufgabe: Arbeitsteilung, einzeln testbar, austauschbar}{}\\\hline
Schnittstelle\rule[-3.2mm]{0pt}{7.5mm} & \lsgoder{\lsgfarbe legt fest, welche Methoden (Signaturen) ein Modul anbietet, ohne die Umsetzung; in Java \code{interface}}{}\\\hline
MVC\rule[-3.2mm]{0pt}{7.5mm} & \lsgoder{\lsgfarbe \textbf{Model}: Daten und Logik, \textbf{View}: Anzeige, \textbf{Controller}: verarbeitet Eingaben, steuert Model und View}{}\\\hline
\end{tabularx}
\end{aufgabe}

\begin{aufgabe}[Informationssicherheit]{5}
Welches \textbf{Schutzziel} (Vertraulichkeit, Integrität, Verfügbarkeit, Authentizität) ist verletzt? Nenne eine Maßnahme.\\[0.5mm]
\renewcommand{\arraystretch}{1.12}
\begin{tabularx}{\linewidth}{|p{7.6cm}|p{2.4cm}|X|}\hline
\textbf{Vorfall} & \textbf{Schutzziel} & \textbf{Maßnahme}\\\hline
\rule[-2mm]{0pt}{6.2mm}Ein Bagger durchtrennt die Internetleitung der Schule. & \lsgoder{\lsgfarbe Verfügbarkeit}{} & \lsgoder{\lsgfarbe zweite Leitung, Backups}{}\\\hline
\rule[-2mm]{0pt}{6.2mm}Jemand ändert heimlich Noten in der Datenbank. & \lsgoder{\lsgfarbe Integrität}{} & \lsgoder{\lsgfarbe Zugriffsrechte, Protokollierung}{}\\\hline
\rule[-2mm]{0pt}{6.2mm}Eine Phishing-Mail gibt sich als Schulleitung aus. & \lsgoder{\lsgfarbe Authentizität}{} & \lsgoder{\lsgfarbe digitale Signatur, Absender prüfen}{}\\\hline
\rule[-2mm]{0pt}{6.2mm}Fremde lesen Schülerdaten auf einem offenen Laptop. & \lsgoder{\lsgfarbe Vertraulichkeit}{} & \lsgoder{\lsgfarbe Bildschirmsperre, Verschlüsselung}{}\\\hline
\end{tabularx}
\end{aufgabe}
\end{document}
)
% Seite 1: Grundwissen 6 bis 11 als ausgefüllte Übersicht (identisch in Blatt und Lösung)
% Seite 2: Übersichtsaufgaben zu den Inhalten der Jahrgangsstufe 12
\usetikzlibrary{shapes.geometric}
\newgeometry{left=14mm,right=14mm,top=21mm,bottom=12mm,headheight=11mm,headsep=3mm,footskip=5mm}
\ifloesung\tikzset{lsgbild/.style={red!75!black}}\else\tikzset{lsgbild/.style={opacity=0}}\fi
\newcommand{\gw}[1]{\par\vspace{0.6mm}{\bfseries\color{green!40!black}#1}\par}
\newcommand{\teil}[1]{{\bfseries #1}\par\vspace{0.5mm}}
\setlength{\columnsep}{5mm}
\setlength{\columnseprule}{0.3pt}
\newcommand{\kopf}[1]{\multicolumn{1}{c}{\textbf{#1}}}
\renewenvironment{aufgabe}[2][]{\begin{tcolorbox}[enhanced,breakable,colback=white,colframe=black,boxrule=0.7pt,arc=3mm,left=2mm,right=2mm,top=1mm,bottom=1mm,before skip=1.8mm,after skip=1.8mm]\textbf{Aufgabe #2)}\ifstrempty{#1}{}{ \textit{#1}}\par\vspace{0.3mm}}{\end{tcolorbox}}
\lstdefinestyle{klein}{basicstyle=\ttfamily\scriptsize,aboveskip=0.8mm,belowskip=0.8mm}
\begin{document}
\abtitel{0}{Grundwissen}
\vspace{-3mm}
\teil{Teil A: Grundwissen der Jahrgangsstufen 6 bis 11 (zum Nachschlagen)}

\begin{minipage}[t]{0.575\linewidth}
\vspace{0pt}
\begin{tikzpicture}[x=1cm,y=1cm,font=\footnotesize]
\node[anchor=north west] (kk) at (0,0) {\klassenkarte{Dreieck}{seitenlaenge\\farbe}{setzeFarbe(neu)}};
\node[anchor=south west,font=\footnotesize\bfseries] at (0,0.05) {Klasse (Bauplan)};
\node[anchor=north west,draw,rounded corners=2.5mm,inner sep=0pt] (ok) at (5.35,0) {\begin{tabular}{l}\rule{0pt}{3.5mm}\textbf{form1: Dreieck}\\[0.5mm]\hline\rule{0pt}{3.5mm}seitenlaenge = 20\\farbe = "blau"\\[0.5mm]\end{tabular}};
\node[anchor=south west,font=\footnotesize\bfseries] at (5.35,0.05) {Objekt (Instanz)};
\draw[decorate,decoration={brace,amplitude=3pt}] ([xshift=2pt]kk.north east) -- ([xshift=2pt]kk.north east|-{(0,-0.5)}) node[midway,right=4pt,font=\scriptsize] {Klassenname};
\draw[decorate,decoration={brace,amplitude=3pt}] ([xshift=2pt]kk.north east|-{(0,-0.58)}) -- ([xshift=2pt]kk.north east|-{(0,-1.35)}) node[midway,right=4pt,font=\scriptsize] {Attribute};
\draw[decorate,decoration={brace,amplitude=3pt}] ([xshift=2pt]kk.north east|-{(0,-1.43)}) -- ([xshift=2pt]kk.south east) node[midway,right=4pt,font=\scriptsize] {Methoden};
\draw[decorate,decoration={brace,amplitude=3pt}] ([xshift=2pt]ok.north east) -- ([xshift=2pt]ok.north east|-{(0,-0.55)}) node[midway,right=4pt,align=left,font=\scriptsize] {Objektname:\\Klasse};
\draw[decorate,decoration={brace,amplitude=3pt}] ([xshift=2pt]ok.north east|-{(0,-0.63)}) -- ([xshift=2pt]ok.south east) node[midway,right=4pt,align=left,font=\scriptsize] {Attribut-\\werte};
\end{tikzpicture}
\end{minipage}\hfill
\begin{minipage}[t]{0.405\linewidth}
\vspace{0pt}\footnotesize
\textbf{Objektorientierung}: Problemlösen durch Analyse der beteiligten \textbf{Objekte} und ihres Zusammenspiels. Die \textbf{Klasse} ist der Bauplan, ein \textbf{Objekt} eine konkrete Ausprägung (Instanz) der Klasse.\\[0.5mm]
\textbf{Punktnotation}: \code{objekt.methode(parameter);}\\ \code{objekt.attribut = wert;}\quad z.\,B. \code{karol.schritt();}
\end{minipage}

\vspace{-1.5mm}
\begin{multicols}{2}\footnotesize\setlength{\parskip}{0.8pt}
\gw{0.1 Beziehungen, Datentypen, Zugriffsrechte}
Objekte können \textbf{Referenzen} auf andere Objekte enthalten (\textbf{Referenzattribute}); die referenzierte Klasse wird zum \textbf{Datentyp}, z.\,B. \code{private Person eigentuemer;} in \code{Haus}:\\[0.3mm]
\begin{tikzpicture}[font=\scriptsize]
\node[draw,inner sep=2pt] (h) {\textbf{Haus}};
\node[draw,inner sep=2pt,right=17mm of h] (p) {\textbf{Person}};
\draw[-{Stealth}] (h) -- node[above,inner sep=1pt] {eigentuemer} node[below,pos=0.85,inner sep=1pt] {1} (p);
\end{tikzpicture}\quad {\scriptsize(Klassendiagramm ohne \code{String}, \code{int})}\\[0.3mm]
Weitere Datentypen: \code{boolean}, \code{char}, \code{double}, \code{float}, \code{long}, \code{byte}, Felder, jede Klasse.
Objekte kommunizieren über \textbf{Methodenaufrufe}. \textbf{Datenkapselung}: Attribute sind \code{private}, Zugriff nur über \textbf{Getter} und \textbf{Setter}.\\[0.5mm]
\renewcommand{\arraystretch}{1}
\begin{tabular}{@{}|c|l|l|@{}}\hline
\code{+} & \code{public} & öffentlich: jede Klasse darf zugreifen\\\hline
\code{\#} & \code{protected} & eigene Klasse und Unterklassen\\\hline
\code{-} & \code{private} & nur innerhalb der eigenen Klasse\\\hline
\end{tabular}

\gw{0.2 Bausteine von Algorithmen: Struktogramm und Java}
\begin{tikzpicture}[x=1cm,y=1cm,font=\scriptsize]
\draw (0,0) rectangle (3.95,1.1); \draw (0,1.1) -- (1.97,0.55) -- (3.95,1.1); \draw (0,0.55) -- (3.95,0.55); \draw (1.97,0) -- (1.97,0.55);
\node at (1.97,0.93) {\ttfamily\itshape istWand()}; \node at (0.33,0.66) {wahr}; \node at (3.55,0.66) {falsch};
\node at (0.98,0.27) {\ttfamily linksDrehen()}; \node at (2.96,0.27) {\ttfamily schritt()};
\begin{scope}[xshift=4.3cm]
\draw (0,0) rectangle (4.05,1.1); \draw (0.4,0) -- (0.4,0.75) -- (4.05,0.75);
\node[anchor=west] at (0.05,0.93) {Wiederhole solange {\ttfamily\itshape !istWand()}};
\node[anchor=west] at (0.5,0.53) {\ttfamily hinlegen()}; \node[anchor=west] at (0.5,0.2) {\ttfamily schritt()};
\end{scope}
\end{tikzpicture}\\[-1.2mm]
\begin{minipage}[t]{0.47\linewidth}
\begin{java}[style=klein]
if (istWand()) {
   linksDrehen();
} else {
   schritt();
}
\end{java}
\end{minipage}\hfill
\begin{minipage}[t]{0.49\linewidth}
\begin{java}[style=klein]
while (!istWand()) {
   hinlegen();
   schritt();
}
\end{java}
\end{minipage}\\[0.3mm]
\textbf{Sequenz}: Anweisungen untereinander; \textbf{einseitig}: \code{if} ohne \code{else}; „Wiederhole 5-mal“: \code{for (int i = 0; i < 5; i++) \{…\}}

\gw{0.3 Feld (Array)}
Folge \underline{fester Länge} von Werten bzw. Objekten \underline{desselben Typs}. Zugriff über den \textbf{Index}; die Zählung beginnt bei 0.
\begin{java}[style=klein]
int zahl[];            // Deklaration
zahl = new int[1000];  // Instanziierung (Index 0..999)
zahl[2] = 10;          // Wertzuweisung (3. Element)
\end{java}

\gw{0.4 Das EVA-Prinzip}
\begin{tikzpicture}[font=\scriptsize,node distance=2.5mm,every node/.style={align=center}]
\node[draw,minimum height=7mm] (e) {\textbf{Eingabe}\\\code{gewicht, groesse}};
\node[draw,ellipse,inner sep=0pt,right=of e] (v) {\textbf{Verarbeitung}\\\code{bmi = gewicht/}\\\code{(groesse*groesse)}};
\node[draw,minimum height=7mm,right=of v] (a) {\textbf{Ausgabe}\\\code{println(…)}};
\draw[-{Stealth}] (e) -- (v); \draw[-{Stealth}] (v) -- (a);
\end{tikzpicture}\\
Bei Methoden: Eingabe über \textbf{Parameter}, Ausgabe als \textbf{Rückgabewert} oder auf dem Bildschirm.

\gw{0.5 Vererbung, @Override und Polymorphismus}
Eine \textbf{Unterklasse} erbt Attribute und Methoden ihrer \textbf{Oberklasse}; in Java von genau \underline{einer} Klasse (keine Mehrfachvererbung).\\[0.3mm]
\begin{minipage}[t]{0.21\linewidth}
\vspace{0pt}
\begin{tikzpicture}[font=\scriptsize]
\node[draw,rectangle split,rectangle split parts=2,align=left,inner sep=2pt] (p) {\textbf{Person}\nodepart{second}alter\\name};
\node[draw,rectangle split,rectangle split parts=2,align=left,inner sep=2pt,below=4mm of p] (s) {\textbf{Schueler}\nodepart{second}klasse};
\draw[-{Triangle[open,length=2mm]}] (s) -- (p);
\end{tikzpicture}
\end{minipage}\hfill
\begin{minipage}[t]{0.78\linewidth}
\vspace{0pt}
\begin{java}[style=klein,aboveskip=0mm,belowskip=0mm]
public class Schueler extends Person {
   private String klasse;
   Schueler(int a, String n, String k) {
      super(a, n);  klasse = k;
   }
}
\end{java}
\end{minipage}\\[0.8mm]
\begin{tabular}{@{}|l|p{6.25cm}|@{}}\hline
\code{extends} & gibt die Oberklasse an\\\hline
\code{super(…)} & ruft den Konstruktor der Oberklasse auf\\\hline
\code{@Override} & Unterklasse \textbf{überschreibt} eine geerbte Methode (gleiche Signatur, neues Verhalten)\\\hline
Polymorphie & Variable vom Typ der Oberklasse kann auf Unterklassen-Objekte verweisen; ausgeführt wird die Methode der \emph{tatsächlichen} Klasse.\\\hline
\end{tabular}

\gw{0.6 Datenbanken und SQL}
Datenbanken speichern strukturierte Daten in \textbf{Tabellen}: Ein \textbf{Datensatz} (Zeile, Tupel) enthält die Attributwerte (Spalten $\to$ \textbf{Attribute}) eines Objekts. Der \textbf{Primärschlüssel} (unterstrichen; ein oder mehrere Attribute) identifiziert jeden Datensatz eindeutig. Beziehungen zwischen Tabellen: \textbf{Fremdschlüssel}.\\[0.3mm]
\textbf{SQL} (Structured Query Language):\\
\begin{tabular}{@{}l@{\quad}l@{}}
\code{SELECT [DISTINCT] a1, a2} & welche Spalten (ohne Duplikate)\\
\code{FROM tabelle} & aus welcher Tabelle\\
\code{WHERE bedingung} & welche Zeilen\\
\code{ORDER BY a [ASC|DESC]} & auf- bzw. absteigend sortieren\\
\end{tabular}\\[0.3mm]
Beispiel: Schüler:innen der 9a mit Geburtstag im Oktober:
\begin{java}[style=klein,language=SQL,aboveskip=0.3mm]
SELECT Nachname, Vorname, Geburtstag FROM Schueler
WHERE Klasse = '9a' AND Geburtstag LIKE '%-10-%'
ORDER BY Geburtstag ASC;
\end{java}

\gw{0.7 Graphen}
Ein Graph besteht aus \textbf{Knoten} und \textbf{Kanten}. \textbf{Pfad}: Folge durch Kanten verbundener Knoten, keine Kante doppelt. \textbf{Zyklus}: Pfad mit gleichem Start- und Endknoten ($\to$ zyklischer/azyklischer Graph). \textbf{Gerichtet}: Pfeile (eine Richtung), \textbf{ungerichtet}: Striche (beide Richtungen). \textbf{Gewichtet}: Kanten tragen Werte (z.\,B. Entfernung). \textbf{Länge eines Pfades}: Anzahl der Kanten bzw. Summe der Gewichte. Algorithmen: Breiten-, Tiefensuche, Dijkstra.

\gw{0.8 Client-Server-Modell}
Als Graph gerichtet und sternförmig mit dem \textbf{Server} in der Mitte: Jeder \textbf{Client} stellt eine \textbf{Anfrage}, der Server schickt eine \textbf{Antwort} (Zyklus der Länge 2). Protokolle: HTTP(S), DNS, SMTP, IMAP, FTP; darunter TCP/IP. Adressierung über die \textbf{IP-Adresse}, auf Hardwareseite über die \textbf{MAC-Adresse}.

\gw{0.9 Codierung und Verschlüsselung}
\textbf{Codierung}: Information nach festen, öffentlichen Regeln anders darstellen (Binärzahl, ASCII, Morsecode, QR-Code). \textbf{Verschlüsselung}: Nur wer den \textbf{Schlüssel} kennt, kann den Klartext lesen. Das Verfahren darf bekannt sein; es muss so viele Schlüssel geben, dass Ausprobieren (Brute Force) aussichtslos ist. \textbf{Symmetrisch}: ein Schlüssel zum Ver- und Entschlüsseln (Caesar, Vigenère, AES); Problem: Schlüsselaustausch. \textbf{Asymmetrisch}: \textbf{öffentlicher} Schlüssel verschlüsselt, nur der \textbf{private} entschlüsselt (RSA).

\gw{0.10 Künstliche Intelligenz und Maschinelles Lernen}
Informatik $\supset$ Data Science $\supset$ KI $\supset$ Maschinelles Lernen $\supset$ Deep Learning.
\textbf{Schwache KI} löst eng begrenzte Aufgaben (Bilderkennung, Chatbot) ohne Bewusstsein, wie alle heutigen Systeme. \textbf{Starke KI} wäre allgemein intelligent wie ein Mensch (gibt es nicht).\\[0.3mm]
\begin{minipage}[c]{0.37\linewidth}
\begin{tikzpicture}[font=\scriptsize,x=1cm,y=1cm]
\node[draw,circle,inner sep=1pt] (n) at (1.3,0) {$\Sigma\,|\,f$};
\node (x1) at (0,0.4) {$x_1$}; \node (x2) at (0,-0.4) {$x_2$}; \node (b) at (1.3,-0.8) {$b$};
\draw[-{Stealth}] (x1) -- node[above,pos=0.4] {$w_1$} (n); \draw[-{Stealth}] (x2) -- node[below,pos=0.4] {$w_2$} (n);
\draw[-{Stealth}] (b) -- (n); \draw[-{Stealth}] (n) -- ++(0.8,0) node[right] {$y$};
\end{tikzpicture}
\end{minipage}\hfill
\begin{minipage}[c]{0.62\linewidth}
\textbf{Neuron}: Eingaben $x_i$ mal \textbf{Gewichte} $w_i$, aufsummiert, plus \textbf{Bias} $b$; die \textbf{Aktivierungsfunktion} $f$ (z.\,B. Schwellenwert) liefert die Ausgabe $y$.
\end{minipage}\\[0.3mm]
\textbf{Überwachtes} Lernen: aus Beispielen mit Lösung (Label), z.\,B. Spamfilter. \textbf{Unüberwachtes}: findet Muster/Gruppen in Daten ohne Label. \textbf{Bestärkendes}: Agent lernt durch Belohnung und Bestrafung.
\end{multicols}

\newpage
\teil{Teil B: Neu in Jahrgangsstufe 12 – kurze Übersichtsaufgaben}
\footnotesize\setlength{\parskip}{1.5pt}
\begin{aufgabe}[Rekursion]{1}
\begin{minipage}[t]{0.35\linewidth}
\vspace{0pt}
\begin{java}[style=klein,aboveskip=0mm,belowskip=0mm]
public int summe(int n) {
   if (n == 0) {
      return 0;             // §\ding{172}§
   }
   return n + summe(n - 1); // §\ding{173}§
}
\end{java}
\end{minipage}\hfill
\begin{minipage}[t]{0.63\linewidth}
\vspace{0pt}
\textbf{a)} Benenne die beiden Bestandteile jeder rekursiven Methode:\\[0.5mm]
\ding{172} \lsg[4.2cm]{Rekursionsanfang (Abbruch)}\enspace \ding{173} \lsg[5.4cm]{rekursiver Aufruf, kleineres Problem}\\[0.5mm]
\textbf{b)} Was passiert, wenn \ding{172} fehlt? \lsg[6.6cm]{endlos viele Aufrufe $\to$ StackOverflowError}\\[0.5mm]
\textbf{c)} Berechne: \code{summe(3)} = \lsg[7.7cm]{3 + summe(2) = … = 3 + 2 + 1 + summe(0) = 6}
\end{minipage}
\end{aufgabe}

\begin{aufgabe}[Rekursive Datenstrukturen]{2}
\textbf{a)} Ergänze die Übersicht.\\[0.5mm]
\renewcommand{\arraystretch}{1.15}
\begin{tabularx}{\linewidth}{|p{2.9cm}|X|X|X|}\hline
 & \textbf{Warteschlange (Queue)} & \textbf{Liste} & \textbf{Binärbaum}\\\hline
Nachfolger je Knoten\rule[-2mm]{0pt}{6.2mm} & genau einer (am Ende \code{null}) & \lsgoder{\lsgfarbe genau einer}{} & \lsgoder{\lsgfarbe bis zu zwei (links, rechts)}{}\\\hline
Einfügen/Entnehmen\rule[-2mm]{0pt}{6.2mm} & \code{add} hinten, \code{poll} vorne: \lsg[1cm]{FIFO} & \lsgoder{\lsgfarbe an beliebiger Stelle}{} & \lsgoder{\lsgfarbe links kleiner, rechts größer}{}\\\hline
\end{tabularx}\\[0.5mm]
\begin{minipage}[t]{0.75\linewidth}
\vspace{0pt}
\textbf{b)} Beim Entwurfsmuster \textbf{Kompositum} erben \code{Knoten} und \code{Abschluss} von der abstrakten Klasse \code{Listenelement}. Welche Aufgabe hat der \code{Abschluss}, welchen Vorteil bringt er?
\lsglinien{2}{Er ersetzt \code{null} am Ende: Jeder Knoten hat immer einen Nachfolger und ruft die Methode dort auf; der Abschluss beendet die Rekursion. Die Abfragen auf \code{null} entfallen.}
\textbf{c)} Füge \textbf{45} in den Suchbaum ein. \textbf{Inorder}-Reihenfolge (links -- Wurzel -- rechts):\\[0.5mm] \lsg[\linewidth]{20, 30, 40, 45, 50, 60, 70 (sortiert!)}\\[0.8mm]
\textbf{d)} Warum findet man Daten im Suchbaum schneller als in einer Liste?\\[0.5mm] \lsg[\linewidth]{Jeder Vergleich halbiert etwa den Suchbereich: ca. $\log_2 n$ statt $n$ Schritte.}
\end{minipage}\hfill
\begin{minipage}[t]{0.23\linewidth}
\vspace{1mm}\centering
\begin{tikzpicture}[level distance=9mm,level 1/.style={sibling distance=19mm},level 2/.style={sibling distance=9.5mm},every node/.style={draw,circle,inner sep=0.5pt,minimum size=5.5mm,font=\scriptsize}]
\node {50}
  child {node {30} child {node {20}} child {node (v) {40}}}
  child {node {70} child {node {60}} child[missing]};
\node[lsgbild,draw,circle,minimum size=5.5mm,font=\scriptsize,below right=4mm and 0mm of v] (n) {45};
\draw[lsgbild] (v) -- (n);
\end{tikzpicture}
\end{minipage}
\end{aufgabe}

\begin{aufgabe}[Nebenläufigkeit]{3}
\begin{minipage}[t]{0.42\linewidth}
\vspace{0pt}
Zwei Kassen (Threads) buchen \textbf{gleichzeitig} ab. Guthaben: 50\,€; Kasse A bucht 5\,€ ab, Kasse B 10\,€.
\begin{java}[style=klein,belowskip=0mm]
public void abbuchen(int betrag) {
   int neu = guthaben - betrag;
   guthaben = neu;
}
\end{java}
\end{minipage}\hfill
\begin{minipage}[t]{0.55\linewidth}
\vspace{0pt}
\textbf{a)} Erkläre, wie am Ende 40\,€ statt 35\,€ auf dem Konto stehen können.
\lsglinien{2}{A und B lesen beide 50\,€. A schreibt 45\,€, danach überschreibt B mit 40\,€: Die Abbuchung von A geht verloren (\emph{Lost Update}, Race Condition).}
\end{minipage}\\[1mm]
\textbf{b)} Der Codeabschnitt, in dem auf eine gemeinsam genutzte Ressource zugegriffen wird, heißt \lsg[3.3cm]{kritischer Abschnitt}. Darin darf sich immer nur \lsg[1.9cm]{ein Thread} befinden (\lsg[3.7cm]{gegenseitiger Ausschluss}); in Java mit dem Schlüsselwort \lsg[2.2cm]{synchronized}. Wartet jeder von zwei Prozessen auf ein Betriebsmittel, das der andere belegt, blockieren beide für immer: \lsg[3.7cm]{Verklemmung (Deadlock)}.
\end{aufgabe}

\begin{aufgabe}[Softwareengineering]{4}
\textbf{a) Wasserfallmodell}: Nummeriere die Phasen:\enspace
\lsg[0.45cm]{\,4} Test\enspace \lsg[0.45cm]{\,1} Analyse\enspace \lsg[0.45cm]{\,5} Wartung\enspace \lsg[0.45cm]{\,3} Implementierung\enspace \lsg[0.45cm]{\,2} Entwurf\\[0.8mm]
Nachteil: \lsg[15.8cm]{Jede Phase wird nur einmal durchlaufen: Fehler und Änderungswünsche fallen erst spät auf und sind dann teuer.}\\[0.8mm]
\textbf{b) Agile Methoden (Scrum)}: Ein \textbf{Sprint} ist ein \lsg[5.3cm]{kurzer, fester Zeitraum (1–4 Wochen)}. Ergebnis: \lsg[4cm]{lauffähiges Teilprodukt}\\[0.8mm]
Vorteil gegenüber dem Wasserfallmodell: \lsg[11.2cm]{frühe Rückmeldung der Kunden; Anforderungen dürfen sich ändern}\\[0.8mm]
\textbf{c)} Erkläre kurz:\\[0.3mm]
\renewcommand{\arraystretch}{1.12}
\begin{tabularx}{\linewidth}{|p{2.3cm}|X|}\hline
Modularisierung\rule[-3.2mm]{0pt}{7.5mm} & \lsgoder{\lsgfarbe Zerlegung in unabhängige Teile mit klarer Aufgabe: Arbeitsteilung, einzeln testbar, austauschbar}{}\\\hline
Schnittstelle\rule[-3.2mm]{0pt}{7.5mm} & \lsgoder{\lsgfarbe legt fest, welche Methoden (Signaturen) ein Modul anbietet, ohne die Umsetzung; in Java \code{interface}}{}\\\hline
MVC\rule[-3.2mm]{0pt}{7.5mm} & \lsgoder{\lsgfarbe \textbf{Model}: Daten und Logik, \textbf{View}: Anzeige, \textbf{Controller}: verarbeitet Eingaben, steuert Model und View}{}\\\hline
\end{tabularx}
\end{aufgabe}

\begin{aufgabe}[Informationssicherheit]{5}
Welches \textbf{Schutzziel} (Vertraulichkeit, Integrität, Verfügbarkeit, Authentizität) ist verletzt? Nenne eine Maßnahme.\\[0.5mm]
\renewcommand{\arraystretch}{1.12}
\begin{tabularx}{\linewidth}{|p{7.6cm}|p{2.4cm}|X|}\hline
\textbf{Vorfall} & \textbf{Schutzziel} & \textbf{Maßnahme}\\\hline
\rule[-2mm]{0pt}{6.2mm}Ein Bagger durchtrennt die Internetleitung der Schule. & \lsgoder{\lsgfarbe Verfügbarkeit}{} & \lsgoder{\lsgfarbe zweite Leitung, Backups}{}\\\hline
\rule[-2mm]{0pt}{6.2mm}Jemand ändert heimlich Noten in der Datenbank. & \lsgoder{\lsgfarbe Integrität}{} & \lsgoder{\lsgfarbe Zugriffsrechte, Protokollierung}{}\\\hline
\rule[-2mm]{0pt}{6.2mm}Eine Phishing-Mail gibt sich als Schulleitung aus. & \lsgoder{\lsgfarbe Authentizität}{} & \lsgoder{\lsgfarbe digitale Signatur, Absender prüfen}{}\\\hline
\rule[-2mm]{0pt}{6.2mm}Fremde lesen Schülerdaten auf einem offenen Laptop. & \lsgoder{\lsgfarbe Vertraulichkeit}{} & \lsgoder{\lsgfarbe Bildschirmsperre, Verschlüsselung}{}\\\hline
\end{tabularx}
\end{aufgabe}
\end{document}
)
% Seite 1: Grundwissen 6 bis 11 als ausgefüllte Übersicht (identisch in Blatt und Lösung)
% Seite 2: Übersichtsaufgaben zu den Inhalten der Jahrgangsstufe 12
\usetikzlibrary{shapes.geometric}
\newgeometry{left=14mm,right=14mm,top=21mm,bottom=12mm,headheight=11mm,headsep=3mm,footskip=5mm}
\ifloesung\tikzset{lsgbild/.style={red!75!black}}\else\tikzset{lsgbild/.style={opacity=0}}\fi
\newcommand{\gw}[1]{\par\vspace{0.6mm}{\bfseries\color{green!40!black}#1}\par}
\newcommand{\teil}[1]{{\bfseries #1}\par\vspace{0.5mm}}
\setlength{\columnsep}{5mm}
\setlength{\columnseprule}{0.3pt}
\newcommand{\kopf}[1]{\multicolumn{1}{c}{\textbf{#1}}}
\renewenvironment{aufgabe}[2][]{\begin{tcolorbox}[enhanced,breakable,colback=white,colframe=black,boxrule=0.7pt,arc=3mm,left=2mm,right=2mm,top=1mm,bottom=1mm,before skip=1.8mm,after skip=1.8mm]\textbf{Aufgabe #2)}\ifstrempty{#1}{}{ \textit{#1}}\par\vspace{0.3mm}}{\end{tcolorbox}}
\lstdefinestyle{klein}{basicstyle=\ttfamily\scriptsize,aboveskip=0.8mm,belowskip=0.8mm}
\begin{document}
\abtitel{0}{Grundwissen}
\vspace{-3mm}
\teil{Teil A: Grundwissen der Jahrgangsstufen 6 bis 11 (zum Nachschlagen)}

\begin{minipage}[t]{0.575\linewidth}
\vspace{0pt}
\begin{tikzpicture}[x=1cm,y=1cm,font=\footnotesize]
\node[anchor=north west] (kk) at (0,0) {\klassenkarte{Dreieck}{seitenlaenge\\farbe}{setzeFarbe(neu)}};
\node[anchor=south west,font=\footnotesize\bfseries] at (0,0.05) {Klasse (Bauplan)};
\node[anchor=north west,draw,rounded corners=2.5mm,inner sep=0pt] (ok) at (5.35,0) {\begin{tabular}{l}\rule{0pt}{3.5mm}\textbf{form1: Dreieck}\\[0.5mm]\hline\rule{0pt}{3.5mm}seitenlaenge = 20\\farbe = "blau"\\[0.5mm]\end{tabular}};
\node[anchor=south west,font=\footnotesize\bfseries] at (5.35,0.05) {Objekt (Instanz)};
\draw[decorate,decoration={brace,amplitude=3pt}] ([xshift=2pt]kk.north east) -- ([xshift=2pt]kk.north east|-{(0,-0.5)}) node[midway,right=4pt,font=\scriptsize] {Klassenname};
\draw[decorate,decoration={brace,amplitude=3pt}] ([xshift=2pt]kk.north east|-{(0,-0.58)}) -- ([xshift=2pt]kk.north east|-{(0,-1.35)}) node[midway,right=4pt,font=\scriptsize] {Attribute};
\draw[decorate,decoration={brace,amplitude=3pt}] ([xshift=2pt]kk.north east|-{(0,-1.43)}) -- ([xshift=2pt]kk.south east) node[midway,right=4pt,font=\scriptsize] {Methoden};
\draw[decorate,decoration={brace,amplitude=3pt}] ([xshift=2pt]ok.north east) -- ([xshift=2pt]ok.north east|-{(0,-0.55)}) node[midway,right=4pt,align=left,font=\scriptsize] {Objektname:\\Klasse};
\draw[decorate,decoration={brace,amplitude=3pt}] ([xshift=2pt]ok.north east|-{(0,-0.63)}) -- ([xshift=2pt]ok.south east) node[midway,right=4pt,align=left,font=\scriptsize] {Attribut-\\werte};
\end{tikzpicture}
\end{minipage}\hfill
\begin{minipage}[t]{0.405\linewidth}
\vspace{0pt}\footnotesize
\textbf{Objektorientierung}: Problemlösen durch Analyse der beteiligten \textbf{Objekte} und ihres Zusammenspiels. Die \textbf{Klasse} ist der Bauplan, ein \textbf{Objekt} eine konkrete Ausprägung (Instanz) der Klasse.\\[0.5mm]
\textbf{Punktnotation}: \code{objekt.methode(parameter);}\\ \code{objekt.attribut = wert;}\quad z.\,B. \code{karol.schritt();}
\end{minipage}

\vspace{-1.5mm}
\begin{multicols}{2}\footnotesize\setlength{\parskip}{0.8pt}
\gw{0.1 Beziehungen, Datentypen, Zugriffsrechte}
Objekte können \textbf{Referenzen} auf andere Objekte enthalten (\textbf{Referenzattribute}); die referenzierte Klasse wird zum \textbf{Datentyp}, z.\,B. \code{private Person eigentuemer;} in \code{Haus}:\\[0.3mm]
\begin{tikzpicture}[font=\scriptsize]
\node[draw,inner sep=2pt] (h) {\textbf{Haus}};
\node[draw,inner sep=2pt,right=17mm of h] (p) {\textbf{Person}};
\draw[-{Stealth}] (h) -- node[above,inner sep=1pt] {eigentuemer} node[below,pos=0.85,inner sep=1pt] {1} (p);
\end{tikzpicture}\quad {\scriptsize(Klassendiagramm ohne \code{String}, \code{int})}\\[0.3mm]
Weitere Datentypen: \code{boolean}, \code{char}, \code{double}, \code{float}, \code{long}, \code{byte}, Felder, jede Klasse.
Objekte kommunizieren über \textbf{Methodenaufrufe}. \textbf{Datenkapselung}: Attribute sind \code{private}, Zugriff nur über \textbf{Getter} und \textbf{Setter}.\\[0.5mm]
\renewcommand{\arraystretch}{1}
\begin{tabular}{@{}|c|l|l|@{}}\hline
\code{+} & \code{public} & öffentlich: jede Klasse darf zugreifen\\\hline
\code{\#} & \code{protected} & eigene Klasse und Unterklassen\\\hline
\code{-} & \code{private} & nur innerhalb der eigenen Klasse\\\hline
\end{tabular}

\gw{0.2 Bausteine von Algorithmen: Struktogramm und Java}
\begin{tikzpicture}[x=1cm,y=1cm,font=\scriptsize]
\draw (0,0) rectangle (3.95,1.1); \draw (0,1.1) -- (1.97,0.55) -- (3.95,1.1); \draw (0,0.55) -- (3.95,0.55); \draw (1.97,0) -- (1.97,0.55);
\node at (1.97,0.93) {\ttfamily\itshape istWand()}; \node at (0.33,0.66) {wahr}; \node at (3.55,0.66) {falsch};
\node at (0.98,0.27) {\ttfamily linksDrehen()}; \node at (2.96,0.27) {\ttfamily schritt()};
\begin{scope}[xshift=4.3cm]
\draw (0,0) rectangle (4.05,1.1); \draw (0.4,0) -- (0.4,0.75) -- (4.05,0.75);
\node[anchor=west] at (0.05,0.93) {Wiederhole solange {\ttfamily\itshape !istWand()}};
\node[anchor=west] at (0.5,0.53) {\ttfamily hinlegen()}; \node[anchor=west] at (0.5,0.2) {\ttfamily schritt()};
\end{scope}
\end{tikzpicture}\\[-1.2mm]
\begin{minipage}[t]{0.47\linewidth}
\begin{java}[style=klein]
if (istWand()) {
   linksDrehen();
} else {
   schritt();
}
\end{java}
\end{minipage}\hfill
\begin{minipage}[t]{0.49\linewidth}
\begin{java}[style=klein]
while (!istWand()) {
   hinlegen();
   schritt();
}
\end{java}
\end{minipage}\\[0.3mm]
\textbf{Sequenz}: Anweisungen untereinander; \textbf{einseitig}: \code{if} ohne \code{else}; „Wiederhole 5-mal“: \code{for (int i = 0; i < 5; i++) \{…\}}

\gw{0.3 Feld (Array)}
Folge \underline{fester Länge} von Werten bzw. Objekten \underline{desselben Typs}. Zugriff über den \textbf{Index}; die Zählung beginnt bei 0.
\begin{java}[style=klein]
int zahl[];            // Deklaration
zahl = new int[1000];  // Instanziierung (Index 0..999)
zahl[2] = 10;          // Wertzuweisung (3. Element)
\end{java}

\gw{0.4 Das EVA-Prinzip}
\begin{tikzpicture}[font=\scriptsize,node distance=2.5mm,every node/.style={align=center}]
\node[draw,minimum height=7mm] (e) {\textbf{Eingabe}\\\code{gewicht, groesse}};
\node[draw,ellipse,inner sep=0pt,right=of e] (v) {\textbf{Verarbeitung}\\\code{bmi = gewicht/}\\\code{(groesse*groesse)}};
\node[draw,minimum height=7mm,right=of v] (a) {\textbf{Ausgabe}\\\code{println(…)}};
\draw[-{Stealth}] (e) -- (v); \draw[-{Stealth}] (v) -- (a);
\end{tikzpicture}\\
Bei Methoden: Eingabe über \textbf{Parameter}, Ausgabe als \textbf{Rückgabewert} oder auf dem Bildschirm.

\gw{0.5 Vererbung, @Override und Polymorphismus}
Eine \textbf{Unterklasse} erbt Attribute und Methoden ihrer \textbf{Oberklasse}; in Java von genau \underline{einer} Klasse (keine Mehrfachvererbung).\\[0.3mm]
\begin{minipage}[t]{0.21\linewidth}
\vspace{0pt}
\begin{tikzpicture}[font=\scriptsize]
\node[draw,rectangle split,rectangle split parts=2,align=left,inner sep=2pt] (p) {\textbf{Person}\nodepart{second}alter\\name};
\node[draw,rectangle split,rectangle split parts=2,align=left,inner sep=2pt,below=4mm of p] (s) {\textbf{Schueler}\nodepart{second}klasse};
\draw[-{Triangle[open,length=2mm]}] (s) -- (p);
\end{tikzpicture}
\end{minipage}\hfill
\begin{minipage}[t]{0.78\linewidth}
\vspace{0pt}
\begin{java}[style=klein,aboveskip=0mm,belowskip=0mm]
public class Schueler extends Person {
   private String klasse;
   Schueler(int a, String n, String k) {
      super(a, n);  klasse = k;
   }
}
\end{java}
\end{minipage}\\[0.8mm]
\begin{tabular}{@{}|l|p{6.25cm}|@{}}\hline
\code{extends} & gibt die Oberklasse an\\\hline
\code{super(…)} & ruft den Konstruktor der Oberklasse auf\\\hline
\code{@Override} & Unterklasse \textbf{überschreibt} eine geerbte Methode (gleiche Signatur, neues Verhalten)\\\hline
Polymorphie & Variable vom Typ der Oberklasse kann auf Unterklassen-Objekte verweisen; ausgeführt wird die Methode der \emph{tatsächlichen} Klasse.\\\hline
\end{tabular}

\gw{0.6 Datenbanken und SQL}
Datenbanken speichern strukturierte Daten in \textbf{Tabellen}: Ein \textbf{Datensatz} (Zeile, Tupel) enthält die Attributwerte (Spalten $\to$ \textbf{Attribute}) eines Objekts. Der \textbf{Primärschlüssel} (unterstrichen; ein oder mehrere Attribute) identifiziert jeden Datensatz eindeutig. Beziehungen zwischen Tabellen: \textbf{Fremdschlüssel}.\\[0.3mm]
\textbf{SQL} (Structured Query Language):\\
\begin{tabular}{@{}l@{\quad}l@{}}
\code{SELECT [DISTINCT] a1, a2} & welche Spalten (ohne Duplikate)\\
\code{FROM tabelle} & aus welcher Tabelle\\
\code{WHERE bedingung} & welche Zeilen\\
\code{ORDER BY a [ASC|DESC]} & auf- bzw. absteigend sortieren\\
\end{tabular}\\[0.3mm]
Beispiel: Schüler:innen der 9a mit Geburtstag im Oktober:
\begin{java}[style=klein,language=SQL,aboveskip=0.3mm]
SELECT Nachname, Vorname, Geburtstag FROM Schueler
WHERE Klasse = '9a' AND Geburtstag LIKE '%-10-%'
ORDER BY Geburtstag ASC;
\end{java}

\gw{0.7 Graphen}
Ein Graph besteht aus \textbf{Knoten} und \textbf{Kanten}. \textbf{Pfad}: Folge durch Kanten verbundener Knoten, keine Kante doppelt. \textbf{Zyklus}: Pfad mit gleichem Start- und Endknoten ($\to$ zyklischer/azyklischer Graph). \textbf{Gerichtet}: Pfeile (eine Richtung), \textbf{ungerichtet}: Striche (beide Richtungen). \textbf{Gewichtet}: Kanten tragen Werte (z.\,B. Entfernung). \textbf{Länge eines Pfades}: Anzahl der Kanten bzw. Summe der Gewichte. Algorithmen: Breiten-, Tiefensuche, Dijkstra.

\gw{0.8 Client-Server-Modell}
Als Graph gerichtet und sternförmig mit dem \textbf{Server} in der Mitte: Jeder \textbf{Client} stellt eine \textbf{Anfrage}, der Server schickt eine \textbf{Antwort} (Zyklus der Länge 2). Protokolle: HTTP(S), DNS, SMTP, IMAP, FTP; darunter TCP/IP. Adressierung über die \textbf{IP-Adresse}, auf Hardwareseite über die \textbf{MAC-Adresse}.

\gw{0.9 Codierung und Verschlüsselung}
\textbf{Codierung}: Information nach festen, öffentlichen Regeln anders darstellen (Binärzahl, ASCII, Morsecode, QR-Code). \textbf{Verschlüsselung}: Nur wer den \textbf{Schlüssel} kennt, kann den Klartext lesen. Das Verfahren darf bekannt sein; es muss so viele Schlüssel geben, dass Ausprobieren (Brute Force) aussichtslos ist. \textbf{Symmetrisch}: ein Schlüssel zum Ver- und Entschlüsseln (Caesar, Vigenère, AES); Problem: Schlüsselaustausch. \textbf{Asymmetrisch}: \textbf{öffentlicher} Schlüssel verschlüsselt, nur der \textbf{private} entschlüsselt (RSA).

\gw{0.10 Künstliche Intelligenz und Maschinelles Lernen}
Informatik $\supset$ Data Science $\supset$ KI $\supset$ Maschinelles Lernen $\supset$ Deep Learning.
\textbf{Schwache KI} löst eng begrenzte Aufgaben (Bilderkennung, Chatbot) ohne Bewusstsein, wie alle heutigen Systeme. \textbf{Starke KI} wäre allgemein intelligent wie ein Mensch (gibt es nicht).\\[0.3mm]
\begin{minipage}[c]{0.37\linewidth}
\begin{tikzpicture}[font=\scriptsize,x=1cm,y=1cm]
\node[draw,circle,inner sep=1pt] (n) at (1.3,0) {$\Sigma\,|\,f$};
\node (x1) at (0,0.4) {$x_1$}; \node (x2) at (0,-0.4) {$x_2$}; \node (b) at (1.3,-0.8) {$b$};
\draw[-{Stealth}] (x1) -- node[above,pos=0.4] {$w_1$} (n); \draw[-{Stealth}] (x2) -- node[below,pos=0.4] {$w_2$} (n);
\draw[-{Stealth}] (b) -- (n); \draw[-{Stealth}] (n) -- ++(0.8,0) node[right] {$y$};
\end{tikzpicture}
\end{minipage}\hfill
\begin{minipage}[c]{0.62\linewidth}
\textbf{Neuron}: Eingaben $x_i$ mal \textbf{Gewichte} $w_i$, aufsummiert, plus \textbf{Bias} $b$; die \textbf{Aktivierungsfunktion} $f$ (z.\,B. Schwellenwert) liefert die Ausgabe $y$.
\end{minipage}\\[0.3mm]
\textbf{Überwachtes} Lernen: aus Beispielen mit Lösung (Label), z.\,B. Spamfilter. \textbf{Unüberwachtes}: findet Muster/Gruppen in Daten ohne Label. \textbf{Bestärkendes}: Agent lernt durch Belohnung und Bestrafung.
\end{multicols}

\newpage
\teil{Teil B: Neu in Jahrgangsstufe 12 – kurze Übersichtsaufgaben}
\footnotesize\setlength{\parskip}{1.5pt}
\begin{aufgabe}[Rekursion]{1}
\begin{minipage}[t]{0.35\linewidth}
\vspace{0pt}
\begin{java}[style=klein,aboveskip=0mm,belowskip=0mm]
public int summe(int n) {
   if (n == 0) {
      return 0;             // §\ding{172}§
   }
   return n + summe(n - 1); // §\ding{173}§
}
\end{java}
\end{minipage}\hfill
\begin{minipage}[t]{0.63\linewidth}
\vspace{0pt}
\textbf{a)} Benenne die beiden Bestandteile jeder rekursiven Methode:\\[0.5mm]
\ding{172} \lsg[4.2cm]{Rekursionsanfang (Abbruch)}\enspace \ding{173} \lsg[5.4cm]{rekursiver Aufruf, kleineres Problem}\\[0.5mm]
\textbf{b)} Was passiert, wenn \ding{172} fehlt? \lsg[6.6cm]{endlos viele Aufrufe $\to$ StackOverflowError}\\[0.5mm]
\textbf{c)} Berechne: \code{summe(3)} = \lsg[7.7cm]{3 + summe(2) = … = 3 + 2 + 1 + summe(0) = 6}
\end{minipage}
\end{aufgabe}

\begin{aufgabe}[Rekursive Datenstrukturen]{2}
\textbf{a)} Ergänze die Übersicht.\\[0.5mm]
\renewcommand{\arraystretch}{1.15}
\begin{tabularx}{\linewidth}{|p{2.9cm}|X|X|X|}\hline
 & \textbf{Warteschlange (Queue)} & \textbf{Liste} & \textbf{Binärbaum}\\\hline
Nachfolger je Knoten\rule[-2mm]{0pt}{6.2mm} & genau einer (am Ende \code{null}) & \lsgoder{\lsgfarbe genau einer}{} & \lsgoder{\lsgfarbe bis zu zwei (links, rechts)}{}\\\hline
Einfügen/Entnehmen\rule[-2mm]{0pt}{6.2mm} & \code{add} hinten, \code{poll} vorne: \lsg[1cm]{FIFO} & \lsgoder{\lsgfarbe an beliebiger Stelle}{} & \lsgoder{\lsgfarbe links kleiner, rechts größer}{}\\\hline
\end{tabularx}\\[0.5mm]
\begin{minipage}[t]{0.75\linewidth}
\vspace{0pt}
\textbf{b)} Beim Entwurfsmuster \textbf{Kompositum} erben \code{Knoten} und \code{Abschluss} von der abstrakten Klasse \code{Listenelement}. Welche Aufgabe hat der \code{Abschluss}, welchen Vorteil bringt er?
\lsglinien{2}{Er ersetzt \code{null} am Ende: Jeder Knoten hat immer einen Nachfolger und ruft die Methode dort auf; der Abschluss beendet die Rekursion. Die Abfragen auf \code{null} entfallen.}
\textbf{c)} Füge \textbf{45} in den Suchbaum ein. \textbf{Inorder}-Reihenfolge (links -- Wurzel -- rechts):\\[0.5mm] \lsg[\linewidth]{20, 30, 40, 45, 50, 60, 70 (sortiert!)}\\[0.8mm]
\textbf{d)} Warum findet man Daten im Suchbaum schneller als in einer Liste?\\[0.5mm] \lsg[\linewidth]{Jeder Vergleich halbiert etwa den Suchbereich: ca. $\log_2 n$ statt $n$ Schritte.}
\end{minipage}\hfill
\begin{minipage}[t]{0.23\linewidth}
\vspace{1mm}\centering
\begin{tikzpicture}[level distance=9mm,level 1/.style={sibling distance=19mm},level 2/.style={sibling distance=9.5mm},every node/.style={draw,circle,inner sep=0.5pt,minimum size=5.5mm,font=\scriptsize}]
\node {50}
  child {node {30} child {node {20}} child {node (v) {40}}}
  child {node {70} child {node {60}} child[missing]};
\node[lsgbild,draw,circle,minimum size=5.5mm,font=\scriptsize,below right=4mm and 0mm of v] (n) {45};
\draw[lsgbild] (v) -- (n);
\end{tikzpicture}
\end{minipage}
\end{aufgabe}

\begin{aufgabe}[Nebenläufigkeit]{3}
\begin{minipage}[t]{0.42\linewidth}
\vspace{0pt}
Zwei Kassen (Threads) buchen \textbf{gleichzeitig} ab. Guthaben: 50\,€; Kasse A bucht 5\,€ ab, Kasse B 10\,€.
\begin{java}[style=klein,belowskip=0mm]
public void abbuchen(int betrag) {
   int neu = guthaben - betrag;
   guthaben = neu;
}
\end{java}
\end{minipage}\hfill
\begin{minipage}[t]{0.55\linewidth}
\vspace{0pt}
\textbf{a)} Erkläre, wie am Ende 40\,€ statt 35\,€ auf dem Konto stehen können.
\lsglinien{2}{A und B lesen beide 50\,€. A schreibt 45\,€, danach überschreibt B mit 40\,€: Die Abbuchung von A geht verloren (\emph{Lost Update}, Race Condition).}
\end{minipage}\\[1mm]
\textbf{b)} Der Codeabschnitt, in dem auf eine gemeinsam genutzte Ressource zugegriffen wird, heißt \lsg[3.3cm]{kritischer Abschnitt}. Darin darf sich immer nur \lsg[1.9cm]{ein Thread} befinden (\lsg[3.7cm]{gegenseitiger Ausschluss}); in Java mit dem Schlüsselwort \lsg[2.2cm]{synchronized}. Wartet jeder von zwei Prozessen auf ein Betriebsmittel, das der andere belegt, blockieren beide für immer: \lsg[3.7cm]{Verklemmung (Deadlock)}.
\end{aufgabe}

\begin{aufgabe}[Softwareengineering]{4}
\textbf{a) Wasserfallmodell}: Nummeriere die Phasen:\enspace
\lsg[0.45cm]{\,4} Test\enspace \lsg[0.45cm]{\,1} Analyse\enspace \lsg[0.45cm]{\,5} Wartung\enspace \lsg[0.45cm]{\,3} Implementierung\enspace \lsg[0.45cm]{\,2} Entwurf\\[0.8mm]
Nachteil: \lsg[15.8cm]{Jede Phase wird nur einmal durchlaufen: Fehler und Änderungswünsche fallen erst spät auf und sind dann teuer.}\\[0.8mm]
\textbf{b) Agile Methoden (Scrum)}: Ein \textbf{Sprint} ist ein \lsg[5.3cm]{kurzer, fester Zeitraum (1–4 Wochen)}. Ergebnis: \lsg[4cm]{lauffähiges Teilprodukt}\\[0.8mm]
Vorteil gegenüber dem Wasserfallmodell: \lsg[11.2cm]{frühe Rückmeldung der Kunden; Anforderungen dürfen sich ändern}\\[0.8mm]
\textbf{c)} Erkläre kurz:\\[0.3mm]
\renewcommand{\arraystretch}{1.12}
\begin{tabularx}{\linewidth}{|p{2.3cm}|X|}\hline
Modularisierung\rule[-3.2mm]{0pt}{7.5mm} & \lsgoder{\lsgfarbe Zerlegung in unabhängige Teile mit klarer Aufgabe: Arbeitsteilung, einzeln testbar, austauschbar}{}\\\hline
Schnittstelle\rule[-3.2mm]{0pt}{7.5mm} & \lsgoder{\lsgfarbe legt fest, welche Methoden (Signaturen) ein Modul anbietet, ohne die Umsetzung; in Java \code{interface}}{}\\\hline
MVC\rule[-3.2mm]{0pt}{7.5mm} & \lsgoder{\lsgfarbe \textbf{Model}: Daten und Logik, \textbf{View}: Anzeige, \textbf{Controller}: verarbeitet Eingaben, steuert Model und View}{}\\\hline
\end{tabularx}
\end{aufgabe}

\begin{aufgabe}[Informationssicherheit]{5}
Welches \textbf{Schutzziel} (Vertraulichkeit, Integrität, Verfügbarkeit, Authentizität) ist verletzt? Nenne eine Maßnahme.\\[0.5mm]
\renewcommand{\arraystretch}{1.12}
\begin{tabularx}{\linewidth}{|p{7.6cm}|p{2.4cm}|X|}\hline
\textbf{Vorfall} & \textbf{Schutzziel} & \textbf{Maßnahme}\\\hline
\rule[-2mm]{0pt}{6.2mm}Ein Bagger durchtrennt die Internetleitung der Schule. & \lsgoder{\lsgfarbe Verfügbarkeit}{} & \lsgoder{\lsgfarbe zweite Leitung, Backups}{}\\\hline
\rule[-2mm]{0pt}{6.2mm}Jemand ändert heimlich Noten in der Datenbank. & \lsgoder{\lsgfarbe Integrität}{} & \lsgoder{\lsgfarbe Zugriffsrechte, Protokollierung}{}\\\hline
\rule[-2mm]{0pt}{6.2mm}Eine Phishing-Mail gibt sich als Schulleitung aus. & \lsgoder{\lsgfarbe Authentizität}{} & \lsgoder{\lsgfarbe digitale Signatur, Absender prüfen}{}\\\hline
\rule[-2mm]{0pt}{6.2mm}Fremde lesen Schülerdaten auf einem offenen Laptop. & \lsgoder{\lsgfarbe Vertraulichkeit}{} & \lsgoder{\lsgfarbe Bildschirmsperre, Verschlüsselung}{}\\\hline
\end{tabularx}
\end{aufgabe}
\end{document}
"
%           (füllt alle Lücken rot aus; siehe \lsg, \lsglinien, \lsgoder)
% =====================================================================
\usepackage[a4paper,left=18mm,right=18mm,top=26mm,bottom=17mm,headheight=15mm,headsep=4mm,footskip=8mm]{geometry}
\usepackage{iftex}
\ifPDFTeX
  \usepackage[T1]{fontenc}
  \usepackage[utf8]{inputenc}
  \usepackage{helvet}
  \usepackage[scaled=0.86]{beramono}
  \renewcommand{\familydefault}{\sfdefault}
\else
  \usepackage{fontspec}
  \setmainfont{texgyreheros}[Extension=.otf,UprightFont=*-regular,BoldFont=*-bold,ItalicFont=*-italic,BoldItalicFont=*-bolditalic]
  \setmonofont{DejaVuSansMono}[Extension=.ttf,UprightFont=*,BoldFont=*-Bold,ItalicFont=*-Oblique,BoldItalicFont=*-BoldOblique,Scale=0.86]
\fi
\usepackage[ngerman,shorthands=off]{babel}
\usepackage{xcolor,graphicx,tabularx,array,enumitem,fancyhdr,tikz,listings,multicol,calc,etoolbox,pifont}
\usepackage[most]{tcolorbox}
\usetikzlibrary{arrows.meta,positioning,calc,decorations.pathreplacing,shapes.multipart}
\usepackage[hidelinks]{hyperref}
\setlength{\parindent}{0pt}
\setlength{\parskip}{3pt}
\setlist{nosep,leftmargin=*}

% ---------- Lösungsschalter ----------
\newif\ifloesung
\ifdefined\mitloesung\loesungtrue\fi
\newcommand{\lsgfarbe}{\color{red!75!black}}

% ---------- Kopf- und Fußzeile (einheitliche Form) ----------
\newcommand{\lehrkraft}{Hau}
\newcommand{\themenbereich}{Grundwissen 6 bis 12}
\newcommand{\abnummer}{}
\pagestyle{fancy}\fancyhf{}
\renewcommand{\headrulewidth}{0.4pt}
\fancyhead[L]{\small Klasse 13\\Datum:}
\fancyhead[C]{\small Themenbereich:\\\themenbereich}
\fancyhead[R]{\small \lehrkraft\ifloesung\\{\lsgfarbe\bfseries Lösung}\fi}
\fancyfoot[C]{\scriptsize edu-mrh.de\,·\,Informatik 13\,·\,Blatt \abnummer\ifloesung\ (Lösung)\fi\,·\,Seite \thepage}
\newcommand{\abtitel}[2]{\renewcommand{\abnummer}{#1}\begin{center}\Large\bfseries #1\quad\underline{#2}\end{center}\vspace{-1mm}}

% ---------- Boxen ----------
\newenvironment{aufgabe}[2][]{\begin{tcolorbox}[enhanced,breakable,colback=white,colframe=black,boxrule=0.7pt,arc=3mm,left=2mm,right=2mm,top=1.5mm,bottom=1.5mm,before skip=2.5mm,after skip=2.5mm]\textbf{Aufgabe #2)}\ifstrempty{#1}{}{ \textit{#1}}\par\vspace{0.6mm}}{\end{tcolorbox}}
\newtcolorbox{merke}[1][Merke]{enhanced,breakable,colback=green!5,colframe=green!40!black,boxrule=0.8pt,arc=2mm,left=2mm,right=2mm,top=1.5mm,bottom=1.5mm,before skip=2.5mm,after skip=2.5mm,title={#1},fonttitle=\bfseries,coltitle=white}
\newtcolorbox{hinweis}{enhanced,breakable,colback=yellow!12,colframe=orange!70!black,boxrule=0.6pt,arc=2mm,left=2mm,right=2mm,top=1mm,bottom=1mm,before skip=2mm,after skip=2mm}
\newcommand{\web}[1]{{\small\color{gray}\textit{Website: #1}}}

% ---------- Lücken, Linien, Karos ----------
\newcommand{\luecke}[1][3.5cm]{\rule[-0.5ex]{#1}{0.4pt}}
\newcommand{\linien}[1]{\par\vspace{1.5mm}\foreach \i in {1,...,#1}{\noindent\rule[0pt]{\linewidth}{0.3pt}\par\vspace{5.5mm}}\vspace{-3mm}}
\newcommand{\kariert}[1]{\par\vspace{1.5mm}\noindent\begin{tikzpicture}\draw[step=5mm,gray!55,thin] (0,0) grid (\linewidth-1pt,#1*5mm);\end{tikzpicture}\par}
\newcommand{\karobox}[2]{\begin{tikzpicture}\draw[step=5mm,gray!55,thin] (0,0) grid (#1,#2);\end{tikzpicture}}
\newcommand{\klassenkarte}[3]{\begin{tabular}{|l|}\hline\multicolumn{1}{|c|}{\textbf{#1}}\\\hline #2\\\hline #3\\\hline\end{tabular}}
\newcommand{\code}[1]{\texttt{#1}}

% Lücke mit Lösung: \lsg[Mindestbreite]{Lösungstext}
\newlength{\lsgbreite}
\newcommand{\lsg}[2][3.5cm]{\ifloesung\settowidth{\lsgbreite}{\,#2\,}\ifdim\lsgbreite<#1\setlength{\lsgbreite}{#1}\fi\underline{\makebox[\lsgbreite][l]{\lsgfarbe\,#2}}\else\luecke[#1]\fi}
% Schreiblinien mit Lösung: \lsglinien{Anzahl}{Lösungstext} (gleiche Höhe wie \linien)
\newcommand{\lsglinien}[2]{\ifloesung\par\vspace{1.5mm}\noindent\begin{minipage}[t][#1\dimexpr 6.9mm\relax][t]{\linewidth}\lsgfarbe #2\end{minipage}\par\vspace{-1.5mm}\else\linien{#1}\fi}
% Beliebiger Inhalt: \lsgoder{nur in der Lösung}{nur im Blatt}
\newcommand{\lsgoder}[2]{\ifloesung\leavevmode#1\else#2\fi}

% ---------- Verkettete Strukturen zeichnen ----------
% Objekt mit Kopf und Attributen:  \node[obj] (k1) {k1: Knoten \nodepart{second} daten = Hans\\ nachfolger};
\tikzset{
  obj/.style={draw,rectangle split,rectangle split parts=2,rectangle split part fill={orange!25,orange!8},
              rounded corners=2pt,font=\scriptsize,inner sep=2pt,align=left},
  qobj/.style={obj,rectangle split part fill={teal!25,teal!8}},
  pobj/.style={draw,rounded corners=2pt,fill=gray!8,font=\scriptsize,inner sep=2pt,align=left},
  zeiger/.style={thick,-{Stealth[length=2.2mm]}},
  lsgzeiger/.style={zeiger,red!75!black},
  nullref/.style={font=\scriptsize\ttfamily,inner sep=1pt},
  schritt/.style={circle,draw,fill=white,inner sep=0.6pt,font=\scriptsize\bfseries},
  lsgschritt/.style={schritt,draw=red!75!black,text=red!75!black},
}

% ---------- Java-Listings ----------
\lstset{language=Java,basicstyle=\ttfamily\small,keywordstyle=\bfseries,commentstyle=\color{gray!80!black},stringstyle=\color{black},showstringspaces=false,columns=flexible,keepspaces=true,tabsize=3,frame=single,framerule=0.4pt,rulecolor=\color{black},xleftmargin=2mm,framexleftmargin=1mm,aboveskip=2mm,belowskip=2mm,breaklines=true,escapechar=§,
 literate={ä}{{\"a}}1 {ö}{{\"o}}1 {ü}{{\"u}}1 {Ä}{{\"A}}1 {Ö}{{\"O}}1 {Ü}{{\"U}}1 {ß}{{\ss}}1 {–}{{--}}1 {„}{{\glqq}}1 {“}{{\grqq}}1 {→}{{$\to$}}1 {①}{{\ding{172}}}1 {②}{{\ding{173}}}1 {③}{{\ding{174}}}1 {④}{{\ding{175}}}1}
\lstnewenvironment{java}[1][]{\lstset{#1}}{}
